\section{Part IX: weakened walks, regularity tiers and the weakened background map}\label{d9:sec}
\setlist{itemsep=1pt,topsep=3pt}

This chapter expands Part~IX of the note, Sections~\Nref{sec:weakenable}--\Nref{sec:weak-background}. Each Detail names the result of the note that it expands and ends with a status (Section~\ref{subsec:guide-purpose}). Part~IX localizes the inputs of the step with \emph{volume-type} decay, a fixed small factor for every cube of the dependence region, which the polymer expansion of Part~X needs. The main line through this Part is:
\begin{itemize}
\item the weakened families and the walk on domains (Lemmas~\Nref{lem:weak-chain}, \Nref{lem:step-small}, Theorem~\Nref{thm:weak-inverse});
\item the five levels, the weakened covariance and the determinant (Corollary~\Nref{cor:weak-covariance}), the one-loop pieces (Proposition~\Nref{prop:one-loop-polymer}) and their torus independence (Lemma~\Nref{lem:torus-transfer}, Proposition~\Nref{prop:one-loop-local});
\item the tier norms and the pointwise bounds (Proposition~\Nref{prop:nl-tiers}, Theorems~\Nref{thm:tier-walk} and~\Nref{thm:HC-pointwise});
\item the weakened background map (Theorem~\Nref{thm:weak-bg}).
\end{itemize}
Section~\Nref{sec:pointwise-weak} is a discussion; only its exponent chain is checked. Section~\ref{d9:status} summarizes what this chapter adds to the note.

% IX.1
\subsection{Conventions, imported results and the \texorpdfstring{$\delta$}{delta}-budget}\label{d9:conv}

\paragraph{Conventions.}
\begin{itemize}
\item The global order of choice of the parameters, the full constants ledger and the table of disambiguated names are those of the parameter ledger (Section~\ref{sec:ledger}; Sections~\ref{subsec:ledger-order} and~\ref{subsec:ledger-names}). Here we state only the constants local to Part~IX and the loss in the weakening rate, the $\delta$-budget.
\item \emph{One collision matters here.} $C_T$ denotes the tier constant of Lemmas~\Nref{lem:int-factors} and~\Nref{lem:endpoint-factors}. It is not the bound on $C_k^{1/2}(s)$ of Section~\Nref{sec:step-quantities}, which the note also calls $C_T$.
\item Distances between blocks and sets are $\infty$-distances; the exponential weights use the Euclidean distance (the Setting of Section~\Nref{sec:tiers}).
\item Two hats occur. Definition~\Nref{def:domain} uses the points at $\infty$-distance $\le1$. Part~VIII (Section~\Nref{sec:dirichlet-averaged}, paragraph ``Conventions for Part VIII'') writes $\hat X$ for $X^{(4)}$. We say which is meant.
\end{itemize}

\paragraph{Constants local to Part IX.}
\begin{center}\small
\begin{tabularx}{\textwidth}{@{}>{\raggedright\arraybackslash}p{0.2\textwidth}>{\raggedright\arraybackslash}X>{\raggedright\arraybackslash}p{0.27\textwidth}@{}}
\toprule
constant & meaning; depends on & where\\
\midrule
$\rho$, $C_v$, $R_v$ & size of $\Lam^v_j$; margin of $X(j)$, $R_v\ge R_0+4+2r_{\rm loc}$; on $L$ & Detail~\ref{d9:rho}\\
$\mathsf N_{\rm b}$, $C_j$ & boxes of $X(j)$; $\#\{j:\Box\subset X(j)\}$; on $d,\rho,C_v$ & Section~\Nref{sec:weakenable}, paragraph ``Boxes and polymers'', Detail~\ref{d9:step-small}\\
$w$ & seam width: $10M$, and $R_0+10M+2L$ at (L5) & Detail~\ref{d9:rho}\\
$\gamma_0$, $C_G$, $C_{\rm loc}$, $C_{\mathfrak N}$ & inputs of a walk: rate, (W2), local commutator, (W3); free of $M$ (Detail~\ref{d9:wi-a}, Step~2) & Details~\ref{d9:L1}--\ref{d9:L5}\\
$\theta$, $M_1$ & step constant and walk threshold & Details~\ref{d9:step-small}, \ref{d9:wi-a}\\
$C_{\rm cpt}$ & compatibility constant of a walk & Detail~\ref{d9:wi-b}\\
$C_T$, $C'_{\mathfrak N}$, $\theta_{\mathfrak N}$, $C^D_{\mathfrak N}$, $\theta'$, $M_{\mathfrak N}$, $M_1'$ & tier constants and thresholds & Details~\ref{d9:int-factors}--\ref{d9:tier-walk}\\
$K_H$, $r_P$, $r'_P$, $\varrho_P$, $R_P$, $\eps_P$, $C_P$ & background map & Details~\ref{d9:Ds}--\ref{d9:weak-bg}\\
$\kappa'$, $\kappa_2=3^{-d}\kappa'$, $q_1=2p+6$, $c_a$ & rates and constants of the one-loop pieces & Details~\ref{d9:cov-cor}, \ref{d9:one-loop}\\
\bottomrule
\end{tabularx}
\end{center}

\begin{detail}[The $\delta$-budget: Section~\Nref{ss:41.3}, Corollary~\Nref{cor:weak-covariance}, Proposition~\Nref{prop:nl-tiers}, Theorem~\Nref{thm:tier-walk}]\label{d9:budget}\Ntarget{ss:41.3,cor:weak-covariance,prop:nl-tiers,thm:tier-walk}
Three rules give the cost of each operation in the weakening rate:
\begin{itemize}
\item a walk with output at $\kappa$ needs its non-local part at $\kappa+2\delta$ (Lemma~\Nref{lem:step-small} at $\kappa+\delta$, Theorem~\Nref{thm:weak-inverse});
\item a flattened product at $\kappa$ needs its factors at $\kappa+\delta$ (Lemma~\Nref{lem:weak-chain});
\item local factors, additive terms and parameter integrals cost nothing.
\end{itemize}
\begin{step}[The $L^2$ route: $C_k^{1/2}(s)$ and $\ell_k(b)$ at $\kappa'$]
Follow the deepest chain:
\begin{center}\small
\begin{tabular}{@{}lll@{}}
\toprule
object & rate & reason\\
\midrule
$R_r(s)$, (L5) & $\kappa'$ & output\\
$\Delta^{\rm c}(s)$ & $\kappa'+2\delta$ & non-local part of (L5); local factors of $\Delta^0$, $\omega_0$\\
$\boldsymbol\Gamma[W_2](s)$, (L4$'$) & $\kappa'+2\delta$ & $\Delta^0=\boldsymbol\Gamma[W_2]-a$\\
$\bG[W_2](s)$, (L3$'$) & $\kappa'+4\delta$ & non-local part of (L4$'$) is $\bQ\bG[W_2]\bQ^T$\\
$\mathfrak N[W_2]$ & $\kappa'+6\delta$ & flattened products of (L1), (L2), $W_2$\\
$W_2(s)$, $\omega_0(s)$, $\bH_0(s)$ & $\kappa'+7\delta$ & factors of these products; $\bH_0=\bG\bQ^T\boldsymbol\Gamma$\\
$\boldsymbol\Gamma(s)$, (L4) & $\kappa'+8\delta$ & factor of $\bH_0$\\
$\bG(s)$, (L3) & $\kappa'+10\delta$ & non-local part of (L4)\\
$\mathfrak N$ of (L3) & $\kappa'+12\delta$ & non-local part of (L3)\\
(L2) outputs & $\kappa'+13\delta$ & factors of the flattened middle factors\\
$Q'G'Q'^T$, $M'$ & $\kappa'+15\delta$ & non-local parts of (L2)\\
(L1) output & $\kappa'+16\delta$ & factors of these products\\
\bottomrule
\end{tabular}
\end{center}
(L1) has $\mathfrak N=0$. So constructing everything from the base rate $\kappa=\kappa'+16\delta$ suffices. The note's count, ``seven walks at $2\delta$ and at most six products at $\delta$, $C\le20$'' (Corollary~\Nref{cor:weak-covariance}), is an upper bound for this, and $C=20$ may be used. The rate conditions on $\kappa'$ are in the parameter ledger (Detail~\ref{det:ledger-rates}).
\end{step}
\begin{step}[The tier route: $\bH(s)$, $\bC(s)$ at $\kappa_{\rm bg}$]
$\bC(s)$ and $\bH(s)$ are flattened products, so their factors are at $\kappa_{\rm bg}+\delta$.
\begin{itemize}
\item The tier walk $\bG(s)$ at $\kappa_{\rm bg}+\delta$ uses its steps at $+2\delta$, and Proposition~\Nref{prop:nl-tiers} at $+2\delta$.
\item That uses Lemma~\Nref{lem:N-tier} at $+3\delta$ and Lemma~\Nref{lem:lift} at $+4\delta$, whose products use the walks of (L1), (L2) at $+5\delta$. Those need (L1) at $+8\delta$.
\item The factor $\boldsymbol\Gamma(s)$ at $+\delta$ needs $\bG$ in $L^2$ at $+3\delta$, its non-local part at $+5\delta$, (L2) at $+6\delta$, and (L1) at $+9\delta$.
\end{itemize}
So the base rate $\kappa_{\rm bg}+9\delta$ suffices. With one base rate $\kappa_0$ for all constructions, the background map lives at $\kappa_0-9\delta$, and the covariance and determinant at $\kappa_0-16\delta$. Theorems~\Nref{thm:HC-pointwise} and~\Nref{thm:weak-bg} are stated ``for $s\in\mathbb D_\kappa$'' with $C(\kappa,\delta)$. The budget above is the meaning of that statement; the note does not write it.
\end{step}
\status{New argument (bookkeeping).}
\end{detail}

\paragraph{Imported results.} Parts~I--VIII are background; their results are used as stated. The table records each result used in this Part, the hypothesis that has to be checked, and where it is checked. Forward imports are marked~(f).
{\footnotesize
\begin{longtable}{@{}>{\raggedright\arraybackslash}p{0.24\textwidth}>{\raggedright\arraybackslash}p{0.42\textwidth}>{\raggedright\arraybackslash}p{0.27\textwidth}@{}}
\toprule
result & used for; hypothesis & checked in\\
\midrule
\endhead
Lemma~\Nref{lem:partition}, Def.~\Nref{def:domain}, $m_X$ (Section~\Nref{ss:19.1}) & $h_j$, $K_j$, the $3^d$ overlap, $J(X)$; $n\ge8M$ & Details~\ref{d9:parametrix}, \ref{d9:fam-setting}\\
Thm.~\Nref{thm:reg-scalar}(a)--(e) & (W2) at (L1), (L2); $(1,\eps)$-regular, $\eps\le\eps_0$, any domain & Details~\ref{d9:L1}, \ref{d9:L2}\\
Thms.~\Nref{thm:reg-coercive}, \Nref{thm:reg-gauge}(a), \Nref{thm:reg-locality}(a), Lemma~\Nref{lem:full-gradient} & (W2) at (L3), with $\nabla\bG$; $(M,\eps)$-regular, $M\ge M_0$; constants free of $M$ & Details~\ref{d9:L3}, \ref{d9:L4}, Step~4\\
Prop.~\Nref{prop:reg-complex} & complex perturbations; $\delta\le\delta_0$, range-$3$ clause & Details~\ref{d9:L3}, \ref{d9:L4}, Step~7\\
Thm.~\Nref{thm:random-walk}(a) (proof) & step hypothesis at (L1) & Detail~\ref{d9:L1}\\
Props.~\Nref{prop:dirichlet-averaged}, \Nref{prop:dirichlet-hessian} & (W2), (W3) at (L4), (L4$'$); $\|W\|\le\eps$, range $\le3$ & Detail~\ref{d9:L4}\\
Lemma~\Nref{lem:coercive-general}, Prop.~\Nref{prop:local-fluct}, \Neqref{eq:gf-local-part}, Lemma~\Nref{lem:avg-transport}(a),(c),(f) & (L5); $R_0\ge R_0(L,M)$ & Details~\ref{d9:rho}, \ref{d9:L5}\\
Prop.~\Nref{prop:det-densities} & the densities; $\Rea\Delta_k\ge\frac18c_L$ & Detail~\ref{d9:cov-cor}\\
Defs.~\Nref{def:bg-activity}, \Nref{def:graded}, Lemma~\Nref{lem:quadratic-bound}, Rem.~\Nref{rem:bg-to-graded} & one-loop pieces; every hypothesis of Lemma~\Nref{lem:quadratic-bound} & Detail~\ref{d9:one-loop}\\
Thm.~\Nref{thm:local-regularity}, Prop.~\Nref{prop:green-eta}, Lemma~\Nref{lem:lattice-sums} & flat $L^p$ estimates; spacing $\le1$, differences of order $\le3$ & Section~\ref{d9:lp}\\
Lemmas~\Nref{lem:iteration}, \Nref{lem:rescaled}, Props.~\Nref{prop:scalar-interior}, \Nref{prop:system-interior}, Lemma~\Nref{lem:local-system}, Thm.~\Nref{thm:ID} & interior estimates; (B) on $\Lam_4$, $\eps$ small & Section~\ref{d9:lp}, Detail~\ref{d9:int-factors}\\
Lemmas~\Nref{lem:cov-expansion}, \Nref{lem:curvature}, \Nref{lem:hodge} & normal form, $c_{\mu\nu}$, Hodge identity; local gauge & Details~\ref{d9:N-tier}, \ref{d9:trace}\\
Lemmas~\Nref{lem:slice}(c), \Nref{lem:Api-identities}(a), Prop.~\Nref{prop:linearized-averages}(d),(e), Lemma~\Nref{lem:avg} & $E=0$, the identity for $\bC$, sup bounds of $\bQ^T$ & Details~\ref{d9:C-identity}, \ref{d9:crit-current}\\
Props.~\Nref{prop:action}, \Nref{prop:Q-expansion}, Lemmas~\Nref{lem:current}, \Nref{lem:linearizing}, \Nref{lem:V-bounds}, Prop.~\Nref{prop:contraction}, Cor.~\Nref{cor:analytic}, Thm.~\Nref{thm:variational}(d), Thm.~\Nref{thm:reg-min} & background map; $|A|_\infty\le c_5$, $|A|\le c_3$, $\eps+\beta\le c_7$ & Details~\ref{d9:crit-current}--\ref{d9:weak-bg}\\
Lemmas~\Nref{lem:regular-monotone}, \Nref{lem:uniformity}, Rem.~\Nref{rem:step-M}, Rem.~\Nref{rem:general-averaging} & $M$- and $L$-independence of constants & Details~\ref{d9:crit-current}, \ref{d9:fam-setting}; Section~\ref{d7:unif}\\
Lemmas~\Nref{lem:eliminate}, \Nref{lem:one-step-basic}(c), Prop.~\Nref{prop:jacobians}, Prop.~\Nref{prop:first-step}(b),(c) & ranges of the local operators & Detail~\ref{d9:inventory}\\
(f) Props.~\Nref{prop:AB}, \Nref{prop:complex-ext}, Lemma~\Nref{lem:norm-conversion}, Lemma~\Nref{lem:raw-new}, Lemma~\Nref{lem:torus}, Rem.~\Nref{rem:domain-constants} & complex references, flat points, wrapping & Details~\ref{xa:d-cmp}, \ref{xa:d-majorant}; Section~\ref{d8:sec}\\
\bottomrule
\end{longtable}}
Textbook facts, stated where used: the Schur test; the Hardy--Littlewood--Sobolev inequality and the Calder\'on--Zygmund theorem on spaces of homogeneous type (Stein, \emph{Singular integrals}, Ch.~V, \S1.2; \emph{Harmonic analysis}, Ch.~I); Marcinkiewicz interpolation; the formula $A^{-1/2}=\frac1\pi\int_0^\infty r^{-1/2}(A+r)^{-1}dr$; Cauchy's estimate; the identity theorem; Neumann series; contractions depending analytically on parameters.

% IX.2
\subsection{Weakenable families}\label{d9:fam}

\begin{detail}[The Setting of Section~\Nref{sec:weakenable}: boxes and polymers]\label{d9:fam-setting}\Ntarget{sec:weakenable}
\begin{step}[The partition]
By Lemma~\Nref{lem:partition}, $\sum_j\chi_j^2=1$, $\chi_j$ vanishes off $K_j=\{|x-z_j|_\infty\le M\}$, and $|\partial^\beta\chi_j|\le C_\beta M^{-|\beta|}$. So $h_j=\chi_j^2$ satisfies $|h_j(x)-h_j(x')|\le2C_1M^{-1}|x-x'|$, and every point lies in at most $3^d$ sets $K_j$. When the partition is at the scale of the polymer grid, $M\mid n$ gives $\check M=M$. The side condition $n\ge C(\rho+C_v)M$, $C\ge8$, gives $n\ge8M$ for Lemma~\Nref{lem:partition}, and $4R_{\rm g}<n$ (Detail~\ref{d9:torus-setting}).
\end{step}
\begin{step}[Boxes]
Each $\{I^\mu_a\}_a$ partitions the circle into intervals of length in $[M/4,M]$.
\begin{itemize}
\item Two boxes that do not touch differ in some direction $\mu$ by non-adjacent intervals. Between them lies a whole interval, so their distance is $\ge M/4$.
\item A set of $\infty$-diameter $<M/4$ meets at most two intervals per direction. The $r$-neighbourhood of a unit block has diameter $1+2r<M/4$ for $r\le M/16$ and $M>16$, so at most $2^d$ boxes lie within $r$ of a block. The note says ``of a point''; the block version is the one used in Lemma~\Nref{lem:weak-basic}(a).
\item A box touches the boxes formed from the same or adjacent intervals in each direction: at most $3^d$, itself included.
\end{itemize}
\end{step}
\begin{step}[$X(j)$]
$(\Lam^v_j)^{(R_v)}$ is the set of points within $\rho M+1+R_v$ of $z_j$, a cube. The boxes meeting a cube form a product of interval sets, which is $*$-connected. So $X(j)$ is a polymer. In each direction it has at most $2(\rho M+1+R_v)/(M/4)+2\le8(\rho+C_v)+10$ intervals. So $\mathsf N_{\rm b}\le(8(\rho+C_v)+10)^d$, which depends on $d,\rho,C_v$ only.
\end{step}
\begin{step}[The margin $R_v$]
The note takes $R_v\ge R_0+4+2r_{\rm loc}$ (Section~\Nref{sec:weakenable}, \emph{Boxes and polymers}) and says that, for $Y=\Lam^v_j$, $(\Lam^v_j)^{(R_v)}$ contains $\hat Y$ and $\widehat{Y^+}$ ``together with their $2r_{\rm loc}$-neighbourhoods''. With Part~VIII's hat, $\widehat{Y^+}=Y^{(R_0+4)}$, so its $2r_{\rm loc}$-neighbourhood is $Y^{(R_0+4+2r_{\rm loc})}\subset Y^{(R_v)}$, and the sentence holds literally. With the hat of Definition~\Nref{def:domain} the sets are smaller. The bound $R_v\ge R_0+2+2r_{\rm loc}$ would be $2$ short for this sentence.

What is used is less:
\begin{itemize}
\item the dependence region $Y^{(R_0+4)}$ ((W2) at (L4), (L5));
\item the neighbourhood $(Y^+)^{(2r_{\mathcal L})}=Y^{(R_0+2r_{\mathcal L})}$ (Step~2 of Lemma~\Nref{lem:torus-transfer}, Detail~\ref{d9:transfer}).
\end{itemize}
Both lie in $Y^{(R_0+\max(4,2r_{\rm loc}))}$, since $r_{\rm loc}\ge2L\ge2$. So $R_0+2+2r_{\rm loc}$ would suffice for the uses. The extra $2$ costs nothing: $R_v\le C_vM$ is a choice, and $C_v$ absorbs it.
\end{step}
\status{Expanded.}
\end{detail}

\begin{detail}[Lemmas~\Nref{lem:decay-norm} and~\Nref{lem:weak-basic}, the norms $|\cdot|^D$]\label{d9:decay}\Ntarget{lem:decay-norm,lem:weak-basic}
\begin{itemize}
\item \emph{Lemma~\Nref{lem:decay-norm}.}
\begin{itemize}
\item (a) is the Schur test and $e^{\gamma|y-y''|}\le e^{\gamma|y-y'|}e^{\gamma|y'-y''|}$.
\item (b): decay $(C,\gamma')$ gives row and column sums $\le CS_{\gamma'-\gamma}$.
\item (c): $h_j$ commutes with $\one_y$, so $\one_y[F,h_j]\one_{y'}=\one_yF\one_{y'}\psi-\psi\one_yF\one_{y'}$ with $\psi=h_j-h_j(c_y)$. Here $|\psi|\le CM^{-1}$ on $B(y)$ and $|\psi|\le CM^{-1}(1+|y-y'|)$ on $B(y')$. With $t=|y-y'|$, $(1+t)e^{-\gamma t/2}\le1+2(e\gamma)^{-1}\le C\gamma^{-1}$ for $\gamma\le1$.
\end{itemize}
\item \emph{$|F|^D_\gamma$.} ``$D$'' is the list of the derivatives occurring in a level:
\begin{itemize}
\item at the scalar levels, $D$ and $D^T_X$;
\item at the gauge level, $D$, $D^T_{X_j}$, the torus divergence and $\nabla$ (Section~\Nref{ss:41.2}).
\end{itemize}
Mixed spaces use the same four combinations.
\item \emph{Lemma~\Nref{lem:weak-basic}.}
\begin{itemize}
\item (a) A term contributes to $\one_yO(s)\one_{y'}$ only if $B(y)$ meets $X^{(r)}_\omega$. Then $X_\omega$ contains one of the $\le2^d$ boxes within $r$ of $B(y)$ (Detail~\ref{d9:fam-setting}), and $|s^{n_\omega}|\le e^{\kappa\deg\omega}$. Column sums are bounded in the same way. $\mathcal P$ is finite, so the monomials are polynomials, and uniform convergence gives analyticity (Weierstrass).
\item (b) $s^{n_\omega}=0$ unless $X_\omega\subset Z$. Then $X_\omega$ lies in one component, and distinct components are at distance $\ge M/4>2r$.
\item (c) $R(u)$ commutes with the monomials.
\end{itemize}
\end{itemize}
\status{Expanded.}
\end{detail}

\begin{detail}[Definition~\Nref{def:weak-family}, the conversion to margin $0$, and additive terms]\label{d9:fam-def}\Ntarget{def:weak-family}
\begin{step}[Margin $0$ (with a norm bound, New argument)]
Let $X^\sharp$ be the union of the boxes meeting $X^{(r)}$. Since $r<M/4$, such a box is within $r$ of a box of $X$, hence touches it. So $|X^\sharp|\le3^d|X|$, and $X^\sharp$ is a polymer. The new monomial $s^{n_\omega}s^{X^\sharp\setminus X}$ has degree $\le\deg\omega+(3^d-1)|X|\le3^d\deg\omega$.

For the norm, a box $\Box\subset X^\sharp_\omega$ either lies in $X_\omega$ or touches a box $\Box'\subset X_\omega$, and there are at most $3^d$ such $\Box'$. So
\[
\|O^\sharp\|_{\kappa,\gamma}=\sup_\Box\sum_{\omega:\Box\subset X^\sharp_\omega}e^{\kappa\deg^\sharp\omega}|O_\omega|_\gamma\le3^d\|O\|_{3^d\kappa,\gamma} .
\]
This is why the one-loop pieces decay at $\kappa_2=3^{-d}\kappa'$ (Detail~\ref{d9:one-loop}).
\end{step}
\begin{step}[Additive terms]
Each site (point, bond, plaquette, unit bond) is based at one point, and the boxes partition $\Tf$. So $\sum_\Box\one_\Box=1$ and $T=\sum_\Box T\one_\Box$. For $T$ of range $\le r$, $T\one_\Box$ has its kernel in $\Box^{(r)}\times\Box$ and its data in $\Box^{(r)}$, so it is a term with polymer $\Box$, monomial $s_\Box$ and margin $r$. The family norm is $\sup_\Box e^\kappa|T\one_\Box|_\gamma\le e^\kappa|T|_\gamma$. The split uses only the paving and the base points, so it commutes with paving-preserving translations and with the transfers of Section~\Nref{subsec:torus-independence}.

A split by indicators does not preserve H\"older tiers. No tier bound is claimed for an additive term alone. In every tier bound, $R(s)=1-P_\cV(s)$ occurs inside a product, with the literal identity (Detail~\ref{d9:lift}, (Z3)). So no indicator split has to preserve a tier.
\end{step}
\status{Expanded; the norm bound for the conversion is a New argument.}
\end{detail}

\begin{detail}[Lemma~\Nref{lem:weak-chain}]\label{d9:chain}\Ntarget{lem:weak-chain}
\begin{step}[Terms]
If $X^{(r)}_{\omega_i}\cap X^{(r)}_{\omega_{i+1}}=\emptyset$, the product vanishes, because the input support of $F_{i,\omega_i}$ and the output support of $F_{i+1,\omega_{i+1}}$ are disjoint. Otherwise the two polymers have boxes within $2r<M/4$ of each other, which therefore touch. So the union is a polymer. The monomial $\prod s^{n_{\omega_i}}$ has support equal to the union of the boxes and degree $\sum\deg\omega_i$. The kernel lies in $X^{(r)}_{\omega_1}\times X^{(r)}_{\omega_n}$, and the data in $\bigcup X^{(r)}_{\omega_i}$. The index set is geometric.
\end{step}
\begin{step}[Counting]
Fix $\Box$ and the first $i$ with $\Box\subset X_{\omega_i}$; the sum over $i$ costs at most $n$. Sum outward from the two ends towards $i$. Given $\omega_{l-1}$, a term $\omega_l$ must have a box meeting $X^{(2r)}_{\omega_{l-1}}$, which is covered by $\le3^d|X_{\omega_{l-1}}|$ boxes. So
\[
\sum_{\omega_l}e^{\kappa\deg\omega_l}|F_{l,\omega_l}|\,w_l\le3^d|X_{\omega_{l-1}}|\sup_{\Box'}\sum_{\omega_l\ni\Box'}e^{\kappa\deg\omega_l}|F_{l,\omega_l}|\,w_l ,
\]
where $w_l$ is the weight $3^d|X_{\omega_l}|$ passed on from the next factor (or $1$ at the end). Each factor receives at most two weights, the anchored factor one from each side, and $|X_\omega|\le\deg\omega$. With $t^2e^{-\delta t}\le4e^{-2}\delta^{-2}$, each factor contributes at most $C_d\delta^{-2}\|F_l\|_{\kappa+\delta,\gamma}$, with $C_d=4\cdot9^de^{-2}$. The decay norm of a product of terms is at most the product of the decay norms (Lemma~\Nref{lem:decay-norm}(a)). The same holds for $|\cdot|^D$, since $|D\Phi_1\cdots\Phi_nD^T|\le|D\Phi_1|\cdots|\Phi_nD^T|$.
\end{step}
\status{Expanded.}
\end{detail}

\begin{detail}[Lemma~\Nref{lem:weak-local-ops}]\label{d9:local-ops}\Ntarget{lem:weak-local-ops}
\begin{step}[(a), the margin]
The kernel of $\mathcal LO_\omega$ at $(x,x')$ is $\sum_z\mathcal L(x,z)O_\omega(z,x')$, with $z\in X^{(r)}_\omega$ and $|x-z|\le r_0$. So the output lies in $X^{(r+r_0)}_\omega$. The note says that $\mathcal L$ depends on the data only within distance $r_0$ of its support. So the entry $\mathcal L(x,z)$ reads data within $r_0$ of $x$ or of $z$, hence within $2r_0$ of $z$. For $O\mathcal L$ the same holds with the roles of the two sites exchanged. So both $\mathcal LO$ and $O\mathcal L$ have margin at most $r+2r_0$, under every reading of the convention. The finer readings give:
\begin{itemize}
\item if $\mathcal L(x,z)$ depends on the data within $r_0$ of each of $x$ and $z$, both have margin $r+r_0$;
\item under the convention of Lemma~\Nref{lem:local-inventory}(c), where the data lie in the ball of radius $r_{\mathcal L}$ about the input site, $\mathcal LO$ has margin $r+r_0$ but $O\mathcal L$ has $r+2r_0$. This is the asymmetry of the consequence of that lemma.
\end{itemize}
The note states $r+2r_0$ (Lemma~\Nref{lem:weak-local-ops}(a), Theorem~\Nref{thm:weak-inverse}(a)). Only boundedly many local factors occur, and margins need only be $\le M/16$, so the margin bookkeeping of Corollary~\Nref{cor:weak-covariance} and Theorem~\Nref{thm:weak-bg}(d) is not affected.
\end{step}
\begin{step}[(a), the norms]
$|\mathcal LO_\omega|_\gamma\le|\mathcal L|_\gamma|O_\omega|_\gamma$. If $\mathcal L$ has range $r'_0$ in block units, then $|\mathcal L|_\gamma\le e^{\gamma r'_0}|\mathcal L|_0$. A row of the block kernel has at most $(2r'_0+1)^d$ non-zero entries, so $|\mathcal L|_0\le(2r'_0+1)^d\max_{y,y'}\|\one_y\mathcal L\one_{y'}\|$. The block Schur norm can exceed $\|\mathcal L\|$; the operator norm is not enough.
\end{step}
\begin{step}[(b)]
Apply Lemma~\Nref{lem:decay-norm}(c) term by term. For the derivative versions:
\[
\begin{aligned}
D[O,h]&=[DO,h]-[D,h]O,\qquad[O,h]D^T=[OD^T,h]+O[h,D^T],\\
D[O,h]D^T&=[DOD^T,h]-[D,h]OD^T+DO[h,D^T],
\end{aligned}
\]
where $[D,h]$ and $[h,D^T]$ are multiplications by $\eta^{-1}$ times differences of $h$, of size $CM^{-1}$. This gives the extra $CM^{-1}\|O\|^D_{\kappa,\gamma}$.
\end{step}
\status{Expanded. The margin is $r+2r_0$, as Lemma~\Nref{lem:weak-local-ops}(a) and Theorem~\Nref{thm:weak-inverse}(a) state.}
\end{detail}

% IX.3
\subsection{The walk on domains}\label{d9:walk}

\paragraph{The walk in one paragraph.} The inverse of an operator $\fA^X$ on a domain $X$ is built from local Dirichlet inverses. Around each grid point $z_j$ the cube $\Lam^v_j$ (the blocks meeting $\{|x-z_j|_\infty\le\rho M\}$) gives the local domain $X_j=\Lam^v_j\cap X$ and the local propagator $G_j=(\fA^{X_j})^{-1}$, extended by zero. With the partition $h_j=\chi_j^2$ of Section~\Nref{sec:domains}, the parametrix $\cK^X=\sum_jh_jG_j$ satisfies $\fA^X\cK^X=I-\sum_j\mathcal R^X_j$, where $\mathcal R^X_j$ collects the commutators $[\fA^X,h_j]$ and a seam term (Lemma~\Nref{lem:parametrix-X}). Each step of the Neumann series $G^X(s)=\sum_{n\ge0}\big(\sum_js^{X(j)}h_jG_j\big)\mathcal R^X(s)^n$ carries the monomial $s^{X(j)}$, where $X(j)$ is the union of the paving boxes that meet $(\Lam^v_j)^{(R_v)}$ (Section~\Nref{sec:weakenable}, paragraph ``Boxes and polymers''; the paragraph ``The weakened step'' of Section~\Nref{ss:41.2} attaches it to the steps). So a term that reaches far carries many factors $s_\Box$; this gives analyticity in $s\in\mathbb D_\kappa$ and decay of volume type. The commutators are small because $|\nabla h_j|\le C/M$ (Lemma~\Nref{lem:step-small}), and the walk converges for $M\ge M_1$ (Theorem~\Nref{thm:weak-inverse}). The hypotheses are (W1) the structure of $\fA^X$ (local part, mass term, non-local part $\mathfrak N^X$), (W2) the bounds and the comparison of domains of the local propagators, and (W3) compatible weakened representations of $\mathfrak N^X$ (Section~\Nref{ss:41.2}). Comparisons between $X$ and $X'\subset X$ are measured in the seam norm $|\cdot|_{\gamma,\Sigma}$ of Section~\Nref{ss:41.1}, where $\Sigma=\Sigma_w(X,X')$ is the set of blocks within distance $w$ of $X\setminus X'$.

Throughout, $X_j=\Lam^v_j\cap X$, $G_j=G^X_j=(\fA^{X_j})^{-1}$, extended by zero, and $\gamma=\gamma_0/8$. Distances between blocks and sets are $\infty$-distances; the weights $e^{\gamma|y-y'|}$ use the Euclidean distance. Since $|\cdot|_\infty\le|\cdot|$, every inequality $d_\Sigma(y)\le|y-w|_\infty$ below also holds with $|y-w|$.

\begin{detail}[\Neqref{eq:rho-choice} and the seam widths, Section~\Nref{sec:nested-walks}]\label{d9:rho}\Ntarget{eq:rho-choice,sec:nested-walks}
\emph{Claim.} The seam width is $w=10M$ at (L1)--(L4) and $w=R_0+10M+2L$ at (L5). The margin satisfies $R_0(L,M)\le C^{R}_LM$. With $\rho\ge 2R_0/M+30$, inequality \Neqref{eq:rho-choice} holds at every level. $\rho$ can be chosen depending on $L$ only.
\begin{step}[$R_0(L,M)$]
Proposition~\Nref{prop:local-fluct}(a) requires $R_0$ large in two ways. First, $R_0\ge 2L+2$ (Section~\Nref{ss:37.4}). Second, the entry bound in the proof of Proposition~\Nref{prop:local-fluct} (Detail~\ref{d8:lf-a}),
\[
|\Delta_k(b,b')-\Delta^{(Y^+)}_k(b,b')|\le C_Le^{-\gamma'(R_0-2L-10M)/2}e^{-\gamma'|b-b'|/2},
\]
in which the sum over the points of the seam costs half of the rate in the distance to the seam, gives by a Schur bound
\[
\|\Delta_Y-\one_{\mathfrak F(Y)}\Delta_k\one_{\mathfrak F(Y)}\|\le C_{\rm seam}e^{-\gamma'(R_0-2L-10M)/2},
\]
at real $U$ and at the complex perturbations of Lemma~\Nref{lem:coercive-general}. This norm must be at most $\frac1{16}c_L$. Since $\one\Delta_k\one\ge\frac14c_L$ at real $U$ and $\Rea\one\Delta_k\one\ge\frac18c_L$ at these perturbations (Lemma~\Nref{lem:coercive-general}), this gives $\Delta_Y\ge\frac18c_L$ and $\Rea\Delta_Y\ge\frac1{16}c_L$. Both conditions hold for
\[
R_0(L,M)=2L+10M+\frac2{\gamma'}\log^+\frac{16C_{\rm seam}}{c_L},
\]
the value fixed in the proof of Proposition~\Nref{prop:local-fluct}; it satisfies $R_0(L,M)\ge2L+2$, since $M\ge1$. The constants $C_{\rm seam}$, $\gamma'$ and $c_L$ do not depend on $M\ge M_0$ (Remark~\Nref{rem:step-M} and the proof of Proposition~\Nref{prop:local-fluct}). With $L<M$ we get $R_0(L,M)\le C^{R}_LM$, where $C^{R}_L=12+\frac2{\gamma'}\log^+(16C_{\rm seam}/c_L)$.
\end{step}
\begin{step}[(L1)--(L4)]
For $w=10M$, \Neqref{eq:rho-choice} reads $\frac34(\rho-1)\ge10$, i.e.\ $\rho\ge\frac{43}3$. This holds for $\rho\ge30$.
\end{step}
\begin{step}[(L5)]
We need $\frac34(\rho-1)M\ge R_0+10M+2L$. From $\rho\ge2R_0/M+30$,
\[
\tfrac34(\rho-1)M\ge\tfrac32R_0+\tfrac{87}4M\ge R_0+10M+2L,
\]
because $\frac{47}4M\ge2L$ for $M>L$.
\end{step}
\begin{step}[Order]
Take $\rho=2C^{R}_L+30=54+\frac4{\gamma'}\log^+(16C_{\rm seam}/c_L)$, with $C^{R}_L$ of Step~1. Then $\rho\ge2R_0/M+30$, and $\rho$ depends on $L$ only, and so does $R_v\le C_vM$ through $R_v\ge R_0+4+2r_{\rm loc}$. Hence $\mathsf N_{\rm b}=\mathsf N_{\rm b}(d,\rho,C_v)$ is independent of $M$, as Section~\Nref{sec:nested-walks} claims.
\end{step}
\status{Expanded. The value of $R_0(L,M)$ is that of the proof of Proposition~\Nref{prop:local-fluct} (Detail~\ref{d8:lf-a}).}
\end{detail}

\begin{detail}[Seam norms and \Neqref{eq:seam-product}, Section~\Nref{ss:41.1}]\label{d9:seam}\Ntarget{eq:seam-product,ss:41.1}
Let $\Sigma$ be a set of blocks and $d_\Sigma=\operatorname{dist}_\infty(\cdot,\Sigma)$.
\begin{step}[Seam form gives a seam norm]
Let $\|\one_yF\one_{y'}\|\le C\sum_{w\in\Sigma}e^{-\gamma'(|y-w|+|w-y'|)}$. For $w\in\Sigma$ we have $|y-y'|\le|y-w|+|w-y'|$, $d_\Sigma(y)\le|y-w|$ and $d_\Sigma(y')\le|w-y'|$. So the weight $e^{\frac{\gamma'}4(|y-y'|+d_\Sigma(y)+d_\Sigma(y'))}$ is at most $e^{\frac{\gamma'}2(|y-w|+|w-y'|)}$. Hence
\[
\sum_{y'}e^{\frac{\gamma'}4(\cdots)}\|\one_yF\one_{y'}\|\le C\sum_{w}e^{-\frac{\gamma'}2|y-w|}\sum_{y'}e^{-\frac{\gamma'}2|w-y'|}\le CS_{\gamma'/2}^2 ,
\]
and the column sums are bounded in the same way. So $|F|_{\gamma'/4,\Sigma}\le CS^2_{\gamma'/2}$.
\end{step}
\begin{step}[Converse, and support on the seam]
For sets $A,A'$ of blocks, the Schur test applied to the weighted kernel gives
\[
\|\one_AF\one_{A'}\|\le|F|_{\gamma,\Sigma}e^{-\gamma(d_\Sigma(A)+d_\Sigma(A'))} .
\]
Now let $F$ be supported on $\Sigma\times\Tf$. Then $d_\Sigma(y)=0$ for the output blocks, and $d_\Sigma(y')\le|y-y'|$. So the seam weight is at most $e^{2\gamma|y-y'|}$, which gives $|F|_{\gamma,\Sigma}\le|F|_{2\gamma}$. The same holds for support on $\Tf\times\Sigma$.
\end{step}
\begin{step}[\Neqref{eq:seam-product}]
For blocks $y,y_1,y_2,y_3$,
\[
|y-y_3|+d_\Sigma(y)+d_\Sigma(y_3)\le2|y-y_1|+\big(|y_1-y_2|+d_\Sigma(y_1)+d_\Sigma(y_2)\big)+2|y_2-y_3| ,
\]
by $d_\Sigma(y)\le|y-y_1|+d_\Sigma(y_1)$ and the triangle inequality. Composing the weighted kernels as matrices gives $|AFB|_{\gamma,\Sigma}\le|A|_{2\gamma}|F|_{\gamma,\Sigma}|B|_{2\gamma}$.
\end{step}
\begin{step}[Products]
The proof of Lemma~\Nref{lem:weak-chain} uses only $|\prod_iF_{i,\omega_i}|_\gamma\le\prod_i|F_{i,\omega_i}|_\gamma$, together with the supports and the counting. Replacing this estimate by \Neqref{eq:seam-product} gives the seam version: one factor is measured in $\|\cdot\|_{\kappa+\delta,\gamma,\Sigma}$, the others in $\|\cdot\|_{\kappa+\delta,2\gamma}$, and the product is bounded in $\|\cdot\|_{\kappa,\gamma,\Sigma}$ with the same constant.
\end{step}
\status{Expanded.}
\end{detail}

\begin{detail}[Lemma~\Nref{lem:parametrix-X}]\label{d9:parametrix}\Ntarget{lem:parametrix-X}
\begin{step}[Distance of $K_j$ from the cut]
$K_j=\{x:|x-z_j|_\infty\le M\}$ (Lemma~\Nref{lem:partition}). A point $x\notin\Lam^v_j$ lies in a block that does not meet $\{|x'-z_j|_\infty\le\rho M\}$, so $|x-z_j|_\infty>\rho M$. Hence $\operatorname{dist}_\infty(K_j,\Tf\setminus\Lam^v_j)>(\rho-1)M$. Between blocks this distance is at least $(\rho-1)M-1$, the form used in Section~\Nref{sec:tiers}.
\end{step}
\begin{step}[Commutator]
$\fA^Xh_jG_j=h_j\fA^XG_j+[\fA^X,h_j]G_j$. Here $[\fA^X,h_j]=[\fA^X_{\rm loc},h_j]+[\mathfrak N^X,h_j]$, because the mass term $\kappa_mm_X$ is a multiplication operator and commutes with $h_j$.
\end{step}
\begin{step}[Local parts]
Let $f$ be supported in $X_j$. By (W1), $\fA^X_{\rm loc}$ and $\fA^{X_j}_{\rm loc}$ differ only by cut terms of range $r_0$:
\begin{itemize}
\item on $\partial X\cap\partial X_j$ the two cuts are the same;
\item on the rest of $\partial X_j$, which lies inside $X$ at $\partial\Lam^v_j$, they differ.
\end{itemize}
So $(\fA^X_{\rm loc}-\fA^{X_j}_{\rm loc})f$ is supported within $r_0$ of $X\setminus X_j\subset\Tf\setminus\Lam^v_j$. Since $(\rho-1)M>r_0$, this support is disjoint from $K_j$. Hence $h_j(\fA^X_{\rm loc}-\fA^{X_j}_{\rm loc})G_j=0$.
\end{step}
\begin{step}[Mass profiles]
$m_Y=\sum_{i\notin J(Y)}\chi_i^2$ with $J(Y)=\{i:\Lam_i\subset Y\}$. At $x\in K_j$ only the $i$ with $x\in K_i$ contribute. For these, $|z_i-z_j|_\infty\le2M$. Every block of $\Lam_i$ meets $\{|x'-z_i|_\infty\le3M\}\subset\{|x'-z_j|_\infty\le5M\}$, hence meets the $\rho M$-cube about $z_j$ for $\rho\ge5$. So $\Lam_i\subset\Lam^v_j$, and therefore $\Lam_i\subset X$ iff $\Lam_i\subset X_j$. Thus $m_X=m_{X_j}$ on $K_j$. (The note states this for $\rho\ge6$; $\rho\ge5$ suffices.)
\end{step}
\begin{step}[Conclusion]
On $\ell^2(X)$ we have $h_j\one_{X_j}=h_j$. In the gauge case, a bond of $\ell^2_1(X)$ with initial point in $K_j$ has both endpoints in $\Lam^v_j$, and its endpoint in $X$ lies in $X_j$, so the bond belongs to $\ell^2_1(X_j)$. Since $\fA^{X_j}G_j=\one_{X_j}$, Steps 3--4 give
\[
h_j\fA^XG_j=h_j+h_j\big(\mathfrak N^X\one_{X_j}-\one_{X_j}\mathfrak N^{X_j}\one_{X_j}\big)G_j .
\]
Summing over $j$ with $\sum_jh_j=1$ gives $\fA^X\cK^X=I-\sum_j\mathcal R^X_j$. Since $G_j=\one_{X_j}G_j$ and $h_j=h_j\one_{X_j}$, the seam term equals $h_j\one_{X_j}(\mathfrak N^X-\mathfrak N^{X_j})\one_{X_j}G_j$ (Section~\Nref{ss:41.2}).
\end{step}
\status{Expanded.}
\end{detail}

\begin{detail}[The weakened step, Section~\Nref{ss:41.2}]\label{d9:wstep}\Ntarget{ss:41.2}
\begin{step}[Polymers]
The first factor and every step of index $j$ carry the monomial $s^{X(j)}$. The local term $s^{X(j)}[\fA^X_{\rm loc},h_j]G_j$ has its kernel in $K_j^{(r_0)}\times X_j$. It depends on the data only on $(X_j)^\wedge\cup K_j^{(r_0)}\subset(\Lam^v_j)^{(R_v)}$. Since $(\Lam^v_j)^{(R_v)}\subset X(j)$, it is a term with polymer $X(j)$ and margin $0$.

Now let $\omega$ be a term of $\mathfrak N(s)$, with margin $r$. A non-local or seam term built from $\omega$ vanishes unless $X_\omega^{(r+1)}$ meets $\Lam^v_j$. The $+1$ is the range of $Y_i$ and $X_i^T$.
\begin{itemize}
\item In that case some box $\Box$ of $X_\omega$ lies within $r+1$ of $\Lam^v_j$, so it meets $(\Lam^v_j)^{(R_v)}$ when $R_v\ge r+1$.
\item Then $\Box\subset X(j)$, so $X(j)\cup X_\omega$ is $*$-connected and is a polymer.
\item Its monomial $s^{X(j)}s^{n_\omega}$ has support equal to the boxes of $X(j)\cup X_\omega$.
\end{itemize}
The margin is at most $r+1$.
\end{step}
\begin{step}[The form $D^T\mathcal A+\mathcal B$]
A left $D^T$ comes from $X_i^T=D^T$, or, in the gauge local commutator, from $D^T[D,h_j]G_j$. Expanding $\mathcal F_0\prod_{i=1}^n(D^T\mathcal A+\mathcal B)$ gives $2^n$ words. Move each $D^T$ onto the factor on its left. A factor then carries $D^T$ on the right exactly when the next factor is of type $\mathcal A$. The words are therefore products $\Phi_0\cdots\Phi_n$ with
\[
\Phi_0\in\{\mathcal F_0,\mathcal F_0D^T\},\qquad\Phi_i\in\{\mathcal A,\mathcal AD^T,\mathcal B,\mathcal BD^T\}\quad(i\ge1),
\]
and the last factor carries no $D^T$ unless an external one is applied.
\end{step}
\status{Expanded.}
\end{detail}

\begin{detail}[Lemma~\Nref{lem:step-small}]\label{d9:step-small}\Ntarget{lem:step-small}
\emph{The idea.} The step $\mathcal R^X(s)$ has three parts, each small for its own reason. The local commutator $[\fA^X_{\rm loc},h_j]G_j$ carries $|\nabla h_j|\le C/M$; the non-local commutator $[\mathfrak N^X,h_j]G_j$ is small because $\mathfrak N$ decays while $h_j$ varies only on the scale $M$; and the seam term lives at distance of order $(\rho-1)M$ from the support of $h_j$. Step~1 regroups the divergences so that only derivatives controlled by (W2) occur, Steps~2--4 bound the three parts, and Steps~5--6 count the overlaps and collect the constant $\theta$, which is linear in $C_G$.
Inputs:
\begin{itemize}
\item (W2): $|G_j|^D_{\gamma_0}\le C_G$, with seam form at the rate $\gamma_0$;
\item (W3): $\|M_i\|_{\kappa+2\delta,2\gamma_0}\le C_{\mathfrak N}$ for the middle factors, with seam norms at $\gamma_0$;
\item $\gamma_0\le1$.
\end{itemize}
Let $C_j=\max_\Box\#\{j:\Box\subset X(j)\}$. Since $X(j)$ lies within $(\rho+C_v+2)M$ of $z_j$ (Lemma~\Nref{lem:local-inventory}, proof of (a)) and the grid has spacing $M$, $C_j\le(2\rho+2C_v+7)^d$, independent of $M$.
\begin{step}[The regrouping identity]
Put $X'=X_j$ and $\sigma_j=(D^T_X-D^T_{X'})G_j$. At points of $X'$ the two divergences agree, so $\sigma_j$ is supported in $X\setminus X'$. Hence $h_j\sigma_j=0$, i.e.\ $h_jD^T_XG_j=h_jD^T_{X'}G_j$. Then
\begin{align*}
[X_i^TM_iD^T_X,h_j]G_j&=X_i^TM_i\big(h_jD^T_X+[D^T_X,h_j]\big)G_j-h_jX_i^TM_iD^T_XG_j\\
&=[X_i^TM_i,h_j]D^T_{X'}G_j+X_i^TM_i[D^T_X,h_j]G_j-h_jX_i^TM_i\sigma_j ,
\end{align*}
because $X_i^TM_ih_jD^T_XG_j=X_i^TM_ih_jD^T_{X'}G_j$ and $h_jX_i^TM_iD^T_XG_j=h_jX_i^TM_iD^T_{X'}G_j+h_jX_i^TM_i\sigma_j$. Similarly,
\[
h_j\big(X_i^TM^X_iD^T_X-X_i^TM^{X'}_iD^T_{X'}\big)G_j=h_jX_i^T(M^X_i-M^{X'}_i)D^T_{X'}G_j+h_jX_i^TM^X_i\sigma_j .
\]
In the sum the $\sigma_j$-terms cancel exactly.

There is one step the note does not write. To compare $M^X_i$ with $M^{X'}_i$ through compatibility, we need the form $\one_{X'}M^X_i\one_{X'}-M^{X'}_i$.
\begin{itemize}
\item On the right, $D^T_{X'}G_j$ is supported in $X'$.
\item On the left, $h_jX_i^T\one_{X\setminus X'}=0$, because $X_i^T\in\{1,D^T\}$ has range one lattice step and $K_j$ is at distance $>(\rho-1)M$ from $X\setminus X'$.
\end{itemize}
\end{step}
\begin{step}[The local commutator]
By hypothesis, each $j$-term has $|\cdot|_{\gamma_0}\le C_GC_{\rm loc}M^{-1}$, with polymer $X(j)$ and monomial weight $e^{\kappa|X(j)|}\le e^{\kappa\mathsf N_{\rm b}}$. The anchored sum over $j$ costs $C_j$. So the contribution to $\|\mathcal A\|+\|\mathcal B\|$, also with $D^T$ on the right, is at most $C_je^{\kappa\mathsf N_{\rm b}}C_GC_{\rm loc}M^{-1}$.
\end{step}
\begin{step}[The non-local commutator]
For a term $F=X^T_iM_{i,\omega}Y_i$,
\[
[F,h_j]G_j=X_i^T[M_{i,\omega},h_j]Y_iG_j+[X_i^T,h_j]M_{i,\omega}Y_iG_j+X_i^TM_{i,\omega}[Y_i,h_j]G_j .
\]
\begin{itemize}
\item In the first and third terms a left $D^T$, if present, stays in front. The rest is the $\mathcal A$-part $[M_{i,\omega},h_j]Y_iG_j$, respectively $M_{i,\omega}[Y_i,h_j]G_j$.
\item $[X_i^T,h_j]$ and $[Y_i,h_j]$ are multiplications by $\eta^{-1}$ times differences of $h_j$, composed with shifts by at most $\eta$. They have size $CM^{-1}$ and are supported in $K_j^{(1)}$.
\item Lemma~\Nref{lem:decay-norm}(c) at the rate $\gamma_0$, applied to $|M_{i,\omega}|_{\gamma_0}\le|M_{i,\omega}|_{2\gamma_0}$, gives $|[M_{i,\omega},h_j]|_{\gamma_0/2}\le C(\gamma_0M)^{-1}|M_{i,\omega}|_{2\gamma_0}$.
\end{itemize}
With $|Y_iG_j|_{\gamma_0}\le C_G$ (W2), each term is bounded in $|\cdot|_{\gamma_0/2}$ by $C(\gamma_0M)^{-1}|M_{i,\omega}|_{2\gamma_0}C_G$ or by $CM^{-1}|M_{i,\omega}|_{2\gamma_0}C_G$. With $D^T$ on the right, $Y_iG_j$ is replaced by $Y_iG_jD^T$ and $[Y_i,h_j]G_jD^T$, which (W2) also bounds. No derivative of $M_i$ enters.
\end{step}
\begin{step}[The seam term]
Let $\Delta_\omega=\one_{X'}M^X_{i,\omega}\one_{X'}-M^{X'}_{i,\omega}$. By (W3), $\sum_\omega$ of $|\Delta_\omega|_{\gamma_0,\Sigma}$ (anchored, with monomial weights) is at most $C_{\mathfrak N}$, with $\Sigma=\Sigma_w(X,X')$. Here $X\setminus X'\subset\Tf\setminus\Lam^v_j$.
\begin{itemize}
\item The output of $h_jX_i^T$ lies in blocks $y$ within distance $1$ of $K_j$. For these, by Step~1 of Detail~\ref{d9:parametrix} and \Neqref{eq:rho-choice},
\[
d_\Sigma(y)\ge(\rho-1)M-2-w\ge\tfrac14(\rho-1)M-2=:D_0 .
\]
\item By \Neqref{eq:seam-product}, $|X_i^T\Delta_\omega Y_iG_j|_{\gamma_0/2,\Sigma}\le C|\Delta_\omega|_{\gamma_0,\Sigma}C_G$.
\item Restricting the output to blocks with $d_\Sigma\ge D_0$ multiplies every entry by at most $e^{-\gamma_0D_0/2}$, in the row sums and in the column sums.
\end{itemize}
Hence
\[
|h_jX_i^T\Delta_\omega Y_iG_j|_{\gamma_0/2}\le Ce^{\gamma_0}e^{-\gamma_0(\rho-1)M/8}C_G|\Delta_\omega|_{\gamma_0,\Sigma},
\]
which is smaller than the note's $e^{-\gamma_0(\rho-1)M/16}$. For the term (ii), $h_jD^T=D^Th_j-[D^T,h_j]$ splits this into an $\mathcal A$-part and a $\mathcal B$-part, both bounded in the same way.
\end{step}
\begin{step}[Counting]
For the anchored norm $\sup_\Box\sum_{(j,\omega):\Box\subset X(j)\cup X_\omega}e^{\kappa(|X(j)|+\deg\omega)}|\cdot|$, split the sum according to where $\Box$ lies.
\begin{itemize}
\item If $\Box\subset X(j)$: there are at most $C_j$ such $j$. A non-zero term needs $X_\omega^{(r+2)}$ to meet $K_j$, so $X_\omega$ contains one of at most $C'$ boxes near $K_j$. Here $C'$ depends on $d$ only, since $K_j$ has side $2M$ and the boxes have sides $\ge M/4$. With the weight $e^{\kappa|X(j)|}\le e^{\kappa\mathsf N_{\rm b}}$ of $s^{X(j)}$, the sum over these $j$ and $\omega$ is then at most $C_jC'e^{\kappa\mathsf N_{\rm b}}\|M_i\|_{\kappa,2\gamma_0}$.
\item If $\Box\subset X_\omega$: for fixed $\omega$ at most $C|X_\omega|\le C\deg\omega$ values of $j$ contribute. Since $te^{-\delta t}\le(e\delta)^{-1}$, and $s^{X(j)}$ again contributes $e^{\kappa|X(j)|}\le e^{\kappa\mathsf N_{\rm b}}$, the sum is at most $Ce^{\kappa\mathsf N_{\rm b}}\delta^{-1}\|M_i\|_{\kappa+\delta,2\gamma_0}$.
\end{itemize}
\end{step}
\begin{step}[The constant]
Steps 2--5 give
\[
\theta=Ce^{\kappa\mathsf N_{\rm b}}\Big(C_GC_{\rm loc}M^{-1}+C\delta^{-1}(\gamma_0M)^{-1}C_{\mathfrak N}C_G+C\delta^{-1}C_{\mathfrak N}C_Ge^{-\gamma_0(\rho-1)M/16}\Big),
\]
which is \emph{linear} in $C_G$: each term of $\mathcal A$, $\mathcal B$, $\mathcal AD^T$ and $\mathcal BD^T$ contains exactly one factor $PG_jP'$ with $P,P'\in\{1,D,D^T\}$. This is the statement of Lemma~\Nref{lem:step-small}, and its proof says why (the paragraph ``The power of $C_G$''). The linear form matters: a bound with $C_G^2$ in the last two terms would dominate the proved one for $C_G\ge1$, but at (L5) $C_G=C_L(\frac1{16}c_L+r)^{-1}$, which is $<1$ for large $r$.

The later use does not depend on the form. Since $C_G(t)\le16C_L/c_L$, also $C_G(t)^2\le(16C_L/c_L)\,C_G(t)$. So either form gives $\theta(t)\le C\,C_G(t)$, and with it the bound $\|R^{(\ge1)}_t\|\le C(c_L+t)^{-2}$ of Detail~\ref{d9:cov-cor}.
\end{step}
\status{Expanded.}
\end{detail}

\begin{detail}[Theorem~\Nref{thm:weak-inverse}(a)]\label{d9:wi-a}\Ntarget{thm:weak-inverse}
\begin{step}[The bound]
A word of length $n$ (Detail~\ref{d9:wstep}) is a product of $n+1$ bounded factors.
\begin{itemize}
\item $\|\mathcal F_0\|_{\kappa+\delta,\gamma}\le C_je^{(\kappa+\delta)\mathsf N_{\rm b}}C_G$, because $|h_jG_j|_\gamma\le|G_j|_{\gamma_0}$ and the anchored sum over $j$ costs $C_j$. The same holds for $\mathcal F_0D^T$.
\item The step factors have norm at most $\theta(\kappa+\delta)$ by Lemma~\Nref{lem:step-small} at $\kappa+\delta$. This uses $\mathfrak N$ at $\kappa+2\delta$, as (W3) provides.
\end{itemize}
Lemma~\Nref{lem:weak-chain} at $\kappa$, with the factors at $\kappa+\delta$, bounds a word by
\[
(n+1)\,C_d\delta^{-2}C_je^{(\kappa+\delta)\mathsf N_{\rm b}}C_G\,(C_d\delta^{-2}\theta)^n .
\]
There are $2^n$ words. If $2C_d\delta^{-2}\theta(\kappa+\delta)\le\frac12$, then
\[
\|G^X\|_{\kappa,\gamma}\le C_d\delta^{-2}C_je^{(\kappa+\delta)\mathsf N_{\rm b}}C_G\sum_{n\ge0}(n+1)2^{-n}=4C_d\delta^{-2}C_je^{(\kappa+\delta)\mathsf N_{\rm b}}C_G .
\]
\end{step}
\begin{step}[$M_1$]
For fixed $\kappa,\delta,C_G,C_{\rm loc},C_{\mathfrak N},\gamma_0,\rho,C_v$, each term of $\theta(\kappa+\delta)$ decreases in $M$ and tends to $0$. So there is a smallest $M_1$ with $2C_d\delta^{-2}\theta\le\frac12$ for all $M\ge M_1$. $M_1$ increases in each multiplicative constant. That is the property used when the constants are doubled at complex references (Detail~\ref{d9:nl-tiers}, Step~8). $C_v$ enters through $\mathsf N_{\rm b}$ and $C_j$, and the note lists it among the arguments of $M_1$ (Theorem~\Nref{thm:weak-inverse}). At a level whose non-local parts come from Theorem~\Nref{thm:weak-inverse}(b) at earlier levels, the constants are free of $M$ only for $M$ at least the thresholds of those levels, with $C_{\rm cpt}$ taken at these thresholds. So the thresholds are fixed level by level, in the order of the levels, and $M_1$ denotes the largest of them (Section~\Nref{sec:nested-walks}, the paragraph before the levels; parameter ledger, Detail~\ref{det:ledger-M0}).
\end{step}
\begin{step}[Derivatives]
$|D\Phi_0\cdots\Phi_nD^T|_\gamma\le|D\Phi_0|_\gamma\prod_{0<i<n}|\Phi_i|_\gamma\,|\Phi_nD^T|_\gamma$.
\begin{itemize}
\item $D\mathcal F_0=\sum_js^{X(j)}(h_jDG_j+[D,h_j]G_j)$, with $|[D,h_j]|\le CM^{-1}$.
\item $\Phi_nD^T$ is $\mathcal AD^T$ or $\mathcal BD^T$ (Lemma~\Nref{lem:step-small}), or $\mathcal F_0D^T$ if $n=0$.
\end{itemize}
So $\|G^X\|^D_{\kappa,\gamma}$ obeys the same bound with a larger $C$.
\end{step}
\begin{step}[Identification at $s=1$]
At $s=1$ the representations of (W3) sum to $\mathfrak N^X$, and the regrouping is an identity. So the series is $\cK^X\sum_n\mathcal R^n$, with $\mathcal R=\sum_j\mathcal R^X_j=D^T\mathcal A+\mathcal B$ at $s=1$. Because of the left $D^T$, $\mathcal R$ need not be small in operator norm uniformly in $\eta$. We argue with the words of Detail~\ref{d9:wstep}. Fix $\eta$. Then $\|D^T\|\le C\eta^{-1}$, a finite number. By Lemma~\Nref{lem:parametrix-X}, $\fA^X\cK^X=I-\mathcal R$, so
\[
\fA^X\cK^X\sum_{n\le N}\mathcal R^n=I-\mathcal R^{N+1}.
\]
Expand $\mathcal R^{N+1}$ into its $2^{N+1}$ words and move each $D^T$ onto the factor on its left. This gives $\mathcal R^{N+1}=D^T\Omega_N+\Omega'_N$, where $\Omega_N$ and $\Omega'_N$ are sums of $2^N$ products of $N+1$ factors from $\{\mathcal A,\mathcal AD^T,\mathcal B,\mathcal BD^T\}$. By Lemma~\Nref{lem:weak-chain} and Lemma~\Nref{lem:step-small} at $\kappa+\delta$, each product is a weakenable family of norm at most $(N+1)(C_d\delta^{-2}\theta)^{N+1}\le(N+1)4^{-N-1}$. By Lemmas~\Nref{lem:weak-basic}(a) and~\Nref{lem:decay-norm}(a), its value at $s=1$ has operator norm at most $2^d$ times that. So
\[
\|\mathcal R^{N+1}\|\le(C\eta^{-1}+1)\,2^d(N+1)2^{-N-2}\longrightarrow0\qquad(N\to\infty).
\]
The partial sums $\cK^X\sum_{n\le N}\mathcal R^n$ converge in operator norm, by Step~1 and Lemma~\Nref{lem:weak-basic}(a). Their limit $G$ satisfies $\fA^XG=I$. The space $\ell^2(X)$ is finite-dimensional (the torus is finite), so a right inverse is the inverse. No property of $\fA^X$ at $s\ne1$ is used, and no coercivity theorem is needed for this step. At the levels with $\mathcal A=0$, such as (L1) and (L5), the shorter argument works: there $\|\mathcal R\|\le2^d\theta<1$.
\end{step}
\begin{step}[Locality and margin]
By Lemma~\Nref{lem:weak-chain}, the margin of a product is the largest margin of its factors. $\mathcal F_0$ has margin $0$, and the steps have margin at most $r+1$ (Detail~\ref{d9:wstep}). Kernel supports and data dependence hold term by term.
\end{step}
\status{Expanded.}
\end{detail}

\begin{detail}[Theorem~\Nref{thm:weak-inverse}(b),(c): compatibility]\label{d9:wi-b}\Ntarget{thm:weak-inverse}
\emph{The idea.} A term of the walk on $X$ and the matched term on $X'\subset X$ have the same polymer and the same monomial. A telescoping identity (Step~1) writes their difference as a sum of products in which exactly one factor compares the two domains. That factor has seam form (Step~2); the other factors are the bounded factors of the walk. Steps~3--4 collect the rates and sum over the words, Step~5 adds the derivatives, and Step~6 shows that only $\theta$ depends on $M$.
Let $X'\subset X$ and $\Sigma=\Sigma_w(X,X')$. On each field space let $P=\one_{X'}$ be the orthogonal projection onto the fields of $X'$ (on $2$-forms: the fields vanishing on the plaquettes with no corner in $X'$), and $Q=\one_X-P$. On gauge fields $Q\ne\one_{X\setminus X'}$: a bond with one endpoint in $X'$ and the other in $X\setminus X'$ lies in $\ell^2_1(X')$ and in $\ell^2_1(X\setminus X')$. On $\ell^2(\mathfrak B(X))$, at (L4), a coarse bond from $X'$ to $X\setminus X'$ lies in neither $\mathfrak B(X')$ nor $\mathfrak B(X\setminus X')$. In all cases, $2$-forms included, the fields in the range of $Q$ lie within one lattice step of $X\setminus X'$, hence in $\Sigma$. On scalar fields and on $\ell^2(\mathfrak F(X))$ (bonds starting in $X$), $Q=\one_{X\setminus X'}$. Write $X'_j=\Lam^v_j\cap X'$, with $G^{X'_j}:=0$ if $X'_j=\emptyset$.
\begin{itemize}
\item The index data of a term are: the grid points $j_0,\dots,j_m$, the kinds of the steps, and the terms $\omega_i$ of the non-local parts.
\item These data do not depend on the domain, since the monomials and the sets $X(j)$ are geometric.
\item So a term $T^X_\Omega=\Phi^X_0\cdots\Phi^X_m$ on $X$ is matched with $T^{X'}_\Omega=\Phi^{X'}_0\cdots\Phi^{X'}_m$ on $X'$, with the same polymer and monomial.
\end{itemize}
We prove:
\begin{itemize}
\item $\|G^X\|_{\kappa,\gamma_0/16}\le C_{\rm cpt}$;
\item $\|PG^XP-G^{X'}\|_{\kappa,\gamma_0/16,\Sigma}\le C_{\rm cpt}$, termwise;
\item the same for $DG^X$, $G^XD^T$ and $DG^XD^T$.
\end{itemize}
The note states the rate $\gamma_0/32$; that is weaker, so it follows. The first item is (a), since $\gamma_0/16\le\gamma_0/8$.
\begin{step}[A telescoping identity]
Let $A_0,\dots,A_m$ be bounded operators between field spaces of $X$, and $A'_0,\dots,A'_m$ the matched operators between field spaces of $X'$, extended by zero, so that $A'_i=PA'_iP$. Below the $A_k$ are the bounded factors of a word (Detail~\ref{d9:wstep}, Step~2), or the bounded pieces of such a factor: an external $D$ is part of $A_0$, an external $D^T$ part of $A_m$, and a left $D^T$ of a step is part of the factor on its left, whose last piece is a local propagator. So $P$ and $Q$ never stand next to a bare derivative. Put $T_k=A'_0\cdots A'_{k-1}PA_k\cdots A_mP$. Then $T_0=PA_0\cdots A_mP$, $T_{m+1}=A'_0\cdots A'_m$, and, inserting $P+Q=\one_X$ after $A_k$,
\[
T_k-T_{k+1}=A'_0\cdots A'_{k-1}(PA_kP-A'_k)\,PA_{k+1}\cdots A_mP+A'_0\cdots A'_{k-1}\,PA_kQ\,A_{k+1}\cdots A_mP
\]
for $k<m$. Also $T_m-T_{m+1}=A'_0\cdots A'_{m-1}(PA_mP-A'_m)$. Summing over $k$ writes
\[
PT^X_\Omega P-T^{X'}_\Omega
\]
as $2m+1$ products, each containing exactly one \emph{distinguished factor}:
\begin{itemize}
\item a difference $\Delta_k=PA_kP-A'_k$, or
\item a cross term $C_k=PA_kQ$.
\end{itemize}
The other factors are factors of the walk on $X'$, factors of the walk on $X$, or $P$.
\end{step}
\begin{step}[The distinguished factors are of seam form]
Let $\Delta_G=\one_{X'_j}G^{X_j}\one_{X'_j}-G^{X'_j}$. The inclusion that recurs below is
\[
X_j\setminus X'_j=\Lam^v_j\cap(X\setminus X')\subset X\setminus X'\quad\Longrightarrow\quad\Sigma_w(X_j,X'_j)\subset\Sigma .
\]
By (W2) for $X'_j\subset X_j$, $\Delta_G$ is of seam form for $\Sigma_w(X_j,X'_j)$. By Detail~\ref{d9:seam}, $|\Delta_G|_{\gamma_0/4,\Sigma}\le CC_GS^2_{\gamma_0/2}$. We go through the four kinds of factor.
\begin{enumerate}[label=(\alph*)]
\item \emph{First factor} $s^{X(j)}h_jG^{X_j}$. Since $h_j$ commutes with $P$ and $G^{X'_j}=PG^{X'_j}P$, we have $\Delta=s^{X(j)}h_jP(G^{X_j}-G^{X'_j})P$. On scalar fields $P\one_{X_j}=\one_{X'_j}$, so $\Delta=s^{X(j)}h_j\Delta_G$. On gauge fields $P\one_{X_j}=\one_{X'_j}+\one_{\rm b}$, where $\one_{\rm b}$ projects onto the bonds from $X_j\setminus X'$ to $X'\setminus\Lam^v_j$: they lie in $\ell^2_1(X')$ and in $\ell^2_1(X_j)$, but not in $\ell^2_1(X'_j)$. These bonds touch $X\setminus X'$ and lie at $\partial\Lam^v_j$, so $h_j\one_{\rm b}=0$ and $\Delta=s^{X(j)}h_j(\Delta_G+\one_{X'_j}G^{X_j}\one_{\rm b})$. The second term has its input in $\Sigma$ and is bounded like the cross term. The cross term $Ph_jG^{X_j}Q$ has its input within one lattice step of $X\setminus X'$, in $\Sigma$, so $|\cdot|_{\gamma,\Sigma}\le|\cdot|_{2\gamma}\le C_G$ for $2\gamma\le\gamma_0$. If $X'_j=\emptyset$, then $X'\subset\Tf\setminus\Lam^v_j$, which lies at distance $>(\rho-1)M$ from $K_j$ (Detail~\ref{d9:parametrix}, Step~1). So $Ph_j=0$, and $\Delta=0$ and the cross term vanish. The vanishing comes from $h_j$, not from $G^{X_j}$: bonds of $\ell^2_1(X_j)$ that cross $\partial X_j$ may touch $X'$, so $PG^{X_j}P$ need not vanish. For the other kinds of factor with $X'_j=\emptyset$, the matched factor on $X'$ is $0$, and the factor on $X$ contains $G^{X_j}$, which lives on $X_j\subset X\setminus X'$ and its crossing bonds. So it is supported on $\Sigma$, and its seam norm is bounded by its norm at the doubled rate (the last item of the proof of Theorem~\Nref{thm:weak-inverse}(b)).
\item \emph{Local commutator} $s^{X(j)}[\fA^X_{\rm loc},h_j]G^{X_j}$. Write $L=[\fA^X_{\rm loc},h_j]$ and $L'=[\fA^{X'}_{\rm loc},h_j]$. At (L1) and (L3) $L$ is a first-order differential operator, of norm of order $\eta^{-1}M^{-1}$; only its products with the local propagators are bounded. So the factor is compared as a whole, except at (L3), where it is two factors of the word (second item below).
\begin{itemize}
\item The difference $PLG^{X_j}P-L'G^{X'_j}$, and the same with $D^T$ on the right, is of seam form for $\Sigma$ with constants $(C_GC_{\rm loc},\gamma_0)$, by the comparison of the local commutators assumed in Theorem~\Nref{thm:weak-inverse}(b). It is checked at each level in Details~\ref{d9:L1}--\ref{d9:L5}, and at $k=0$ in Detail~\ref{xb:d-k0-walks}.
\item At (L3) Detail~\ref{d9:wstep}, Step~2, and Detail~\ref{d9:L3}, Step~4, write the part $\one_XD^T[D,h_j]G^{X_j}$ of $LG^{X_j}$ as a left $D^T$ in front of the $\mathcal A$-part $[D,h_j]G^{X_j}$. Then the two parts are separate factors of the word, and they are compared separately. $[D,h_j]$ is a multiplication by differences of $h_j$, of size $CM^{-1}$ and range one lattice step, and $[D,h_j]G^{X'_j}$ takes values on plaquettes with a corner in $X'$. So $P[D,h_j]G^{X_j}P-[D,h_j]G^{X'_j}=P[D,h_j](G^{X_j}-G^{X'_j})P$, of seam form by (W2). The rest, $LG^{X_j}-\one_XD^T[D,h_j]G^{X_j}$, is compared as in Detail~\ref{d9:L3}, Step~4, through the comparisons of $G^{X_j}$, $DG^{X_j}$, $\nabla G^{X_j}$ and of the torus divergence in (W2).
\item The cross term $PLG^{X_j}Q$ has its input in $\Sigma$, and $|LG^{X_j}|_{2\gamma}\le C_GC_{\rm loc}M^{-1}$ by the hypothesis of Lemma~\Nref{lem:step-small}. At (L3) the cross terms $P[D,h_j]G^{X_j}Q$ and $P(LG^{X_j}-\one_XD^T[D,h_j]G^{X_j})Q$ of the two parts are bounded in the same way, by Detail~\ref{d9:L3}, Step~4.
\end{itemize}
Apart from the split at (L3), whose left $D^T$ is part of the propagator on its left, the factor is not split by the projections: in such a split, the piece $PLQ\,G^{X_j}P$ differentiates the field $QG^{X_j}Pf$, which is cut sharply at $\partial X'$, and it is not bounded uniformly in $\eta$ when $\supp\partial h_j$ meets $\partial X'$.
\item \emph{Non-local commutator} $s^{X(j)}s^{n_\omega}[X_i^TM^X_{i,\omega}Y_i,h_j]G^{X_j}$, in the regrouped form. Its pieces are bounded: $M^X_{i,\omega}$ or $[M^X_{i,\omega},h_j]$, the multiplications $[X_i^T,h_j]$, $[Y_i,h_j]$, and $Y_iG^{X_j}$, that is $G^{X_j}$, $DG^{X_j}$ or (regrouped) $D^T_{X_j}G^{X_j}$, also with $D^T$ on the right. A left $D^T=X_i^T$ is not a piece: it belongs to the factor on the left (Detail~\ref{d9:wstep}, Step~2). Telescoping over the pieces gives three kinds of distinguished pieces:
\begin{itemize}
\item $\Delta_M=PM^X_{i,\omega}P-M^{X'}_{i,\omega}$, of seam form by (W3);
\item $\Delta_G$, and, with derivatives, $PD(G^{X_j}-G^{X'_j})P$ and the other comparisons of (W2), always of the zero extensions (Step~5);
\item cross terms $P\Phi Q$ of a piece $\Phi$, and differences of the multiplications $[X_i^T,h_j]$, $[Y_i,h_j]$; these are bounded and supported within one lattice step of $X\setminus X'$.
\end{itemize}
The divergences of the regrouping agree exactly: $PD^T_X=\one_{X'}D^T=D^T_{X'}$ and $PD^T_{X_j}=\one_{X'_j}D^T=D^T_{X'_j}$.
\item \emph{Seam term} $s^{X(j)}s^{n_\omega}h_jX_i^T(\one_{X_j}M^X_{i,\omega}\one_{X_j}-M^{X_j}_{i,\omega})Y_iG^{X_j}$. The inner differences on $X$ and on $X'$ give
\begin{multline*}
\one_{X'_j}(\one_{X_j}M^X\one_{X_j}-M^{X_j})\one_{X'_j}-(\one_{X'_j}M^{X'}\one_{X'_j}-M^{X'_j})\\
=\one_{X'_j}(PM^XP-M^{X'})\one_{X'_j}-(\one_{X'_j}M^{X_j}\one_{X'_j}-M^{X'_j}) .
\end{multline*}
The first part is of seam form for $\Sigma$ by (W3) for $X'\subset X$. The second is of seam form for $\Sigma_w(X_j,X'_j)\subset\Sigma$ by (W3) for $X'_j\subset X_j$. Four domains are involved; the inclusion displayed above is what makes all their seams lie in $\Sigma$. The factors $\Delta_G$ and the cross terms are treated as in (a) and (c). On gauge fields and $2$-forms the outer projections $\one_{X'_j}$ above stand for $P$; the two differ on the bonds of (a) and on the plaquettes with a corner in $X_j\setminus X'$ and one in $X'\setminus\Lam^v_j$, which have an index in $\Sigma$.
\end{enumerate}
\end{step}
\begin{step}[Rates]
Inside a distinguished factor there is one seam-form piece, at rate $\gamma_0/4$ (from (W2) or from the comparison of the local commutators) or $\gamma_0$ (from (W3)). The other pieces are $h_j$, the multiplications $[X_i^T,h_j]$, $[Y_i,h_j]$, the local commutators with their propagators, $G^{X_j}$ with its derivatives at the rate $\gamma_0$, and $M$ at $2\gamma_0$. By \Neqref{eq:seam-product}, each distinguished family has
\[
\|\cdot\|_{\kappa+\delta,\gamma_0/4,\Sigma}\le C_1:=Ce^{(\kappa+\delta)\mathsf N_{\rm b}}C_j\,\delta^{-1}C_G(1+C_{\rm loc}+C_{\mathfrak N}) .
\]
The $\delta^{-1}$ and the anchoring are as in Steps~5--6 of Detail~\ref{d9:step-small}. The non-distinguished factors are bounded at the rate $\gamma_0/8$, by $\theta$ (steps) or by $C_je^{(\kappa+\delta)\mathsf N_{\rm b}}C_G$ (first factors). The seam product at the rate $\gamma_s=\gamma_0/16$ needs the neighbours at $2\gamma_s=\gamma_0/8$, which holds.
\end{step}
\begin{step}[Summation]
Each of the $2m+1$ products of a word of length $m$ is bounded by the seam version of Lemma~\Nref{lem:weak-chain} (Detail~\ref{d9:seam}, Step~4). Its bound is the product of the distinguished factor $C_d\delta^{-2}C_1$, the first factor $C_d\delta^{-2}C_je^{(\kappa+\delta)\mathsf N_{\rm b}}C_G$, and at least $m-1$ step factors $C_d\delta^{-2}\theta\le\frac14$. With $2^m$ words, $(m+1)$ from the product lemma and the $2m+1$ positions, the sum over $m$ is at most
\[
C_d^2\delta^{-4}C_1C_je^{(\kappa+\delta)\mathsf N_{\rm b}}C_G\sum_{m\ge0}(2m+1)(m+1)\,2^m\,4^{-\max(m-1,0)}<\infty .
\]
This defines $C_{\rm cpt}$. It depends on $\kappa,\delta,\gamma_0,C_G,C_{\rm loc},C_{\mathfrak N},\rho,C_v$, and not on $X$, $k$, $n$, $\eta$.
\end{step}
\begin{step}[Derivatives]
$D$ is part of the first factor and $D^T$ part of the last: $DA_0$ and $A_mD^T$ are bounded factors of the word, and Step~1 applies to them. The distinguished factor at the first position is $P(DA_0)P-DA'_0$, the comparison of the already differentiated factor. It is not $D(PA_0P-A'_0)$, which contains $[D,P]A_0P$: on a bond crossing $\partial X'$ this is $\eta^{-1}$ times the value of $A_0P$ at the outer endpoint, of order $\eta^{-1/2}$ in $L^2$ for smooth data, and its support on $\Sigma$ does not bound it. For $A_0=s^{X(j)}h_jG^{X_j}$, $DA_0=s^{X(j)}(h_jDG^{X_j}+[D,h_j]G^{X_j})$ up to shifts by $\eta$, with $|[D,h_j]|\le CM^{-1}$. The regrouped divergences agree exactly (Step~2(c)). The derivative terms are handled by the following clauses of (W2), applied to the domains $X'_j\subset X_j$:
\begin{center}\small
\begin{tabularx}{\textwidth}{@{}>{\raggedright\arraybackslash}p{0.36\textwidth}>{\raggedright\arraybackslash}X@{}}
\toprule
term & clause of (W2)\\
\midrule
$PD^T_{X_j}G^{X_j}P-D^T_{X'_j}G^{X'_j}$ (regrouped divergence) & comparison of domains for the torus divergence, with output points within $\eta$ of $X'_j$ (Theorem~\Nref{thm:reg-locality}(a)); at the repetition with $W_2$, Proposition~\Nref{prop:dirichlet-averaged}(d) (Detail~\ref{d9:L4}, Step~4)\\
$PD(G^{X_j}-G^{X'_j})P$, $P\nabla(G^{X_j}-G^{X'_j})P$, derivatives of the zero extensions & comparison of domains for $P=D$ and $P=\nabla$\\
$PG^{X_j}D^TP-G^{X'_j}D^T$ & comparison of domains for $Q=D^T$\\
\bottomrule
\end{tabularx}
\end{center}
For the first row, $PD^T_{X_j}=\one_{X'_j}D^T$ and $G^{X'_j}=G^{X'_j}P$, so the row equals $\one_{X'_j}D^T(G^{X_j}-G^{X'_j})P$. This is the torus-divergence comparison with outputs in $X'_j$. At the scalar levels the same table holds with $D^T_X$ in place of the torus divergence (Theorem~\Nref{thm:reg-scalar}(e)).
\end{step}
\begin{step}[Dependence on $M$ and on $C_G$]
In $C_{\rm cpt}$ only $\theta$ depends on $M$, and it decreases in $M$. So the value at $M=M_1$ serves for all $M\ge M_1$.

The dependence on $C_G$ cannot be removed. For $\fA^X=tI_X$ on unit-lattice fields and $\mathfrak N=0$, the family is $t^{-1}\sum_js^{X(j)}h_j$. By Lemma~\Nref{lem:weak-basic}(a), $t^{-1}=\|G^X(1)\|\le2^d\|G^X\|$, so $\|G^X\|\ge2^{-d}t^{-1}$. The note carries this factor in the bound of Theorem~\Nref{thm:weak-inverse}(a) and lists $C_G$ among the arguments of $C_{\rm cpt}$ in (b).
\end{step}
\begin{step}[(c)]
Real gauge transformations act by conjugation with $R(u)$, which commutes with the multiplications by $h_j$, with $P$, $Q$ and with the local operators. So products of covariant factors are covariant.
\end{step}
\status{Expanded. The telescoping of Step~1 and the inclusion of the four seams in Step~2(d) are written out here; the note compresses both into one sentence.}
\end{detail}

% IX.4
\subsection{The five levels}\label{d9:levels}

The five levels apply the walk to the operators of the step, from the scalar operators of Part~II to the fluctuation operator of Part~VII: (L1) the scalar operators, (L2) the scalar averaged operators, (L3) the gauge-field operators with their non-local parts, (L4) the averaged gauge operators and the Hessian, whose local propagators are the Dirichlet operators of Part~VIII, and (L5) the fluctuation operator (Section~\Nref{ss:41.3}). At each level the hypotheses (W1)--(W3) follow from the earlier levels and from Parts~III, VII and~VIII (Section~\Nref{ss:41.3}); each nesting of walks and products costs $\delta$ or $2\delta$ in the weakening rate (Detail~\ref{d9:budget}). The order of the levels is (L1), (L2), then (L3) and (L4) with $W=0$, then $\omega_0$ and $W_2$, then (L3) and (L4) with $W_2$, then (L5) (Section~\Nref{ss:41.3}). At each level we check (W1)--(W3), the hypothesis of Lemma~\Nref{lem:step-small} and the comparison of the local commutators assumed in Theorem~\Nref{thm:weak-inverse}(b). The compatibility of products is used at every level; the note states it in Section~\Nref{ss:41.3} without proof. We prove it first.

\begin{detail}[Compatibility of products and local factors (New argument), Section~\Nref{ss:41.3}]\label{d9:prod-compat}\Ntarget{ss:41.3}
\emph{Claim.} Let $O_1,\dots,O_n$ have compatible representations with constants $(C_i,\gamma_i)$ and seam width $w$. Let $\mathcal L^X$ be bounded local operators of range $r_0$, with $|\mathcal L^X|_{2\gamma}$ bounded uniformly in $X$ and $\eta$, that depend on $X$ only through cut terms within $r_0$ of $\Tf\setminus X$, such as $\one_X$, $\bQ_X$, $Q'$, $\mathfrak F$, $W$ and the additive terms split box by box. A fine derivative $D$, $D^T$ or $D^T_X$ is not such an operator: its norm is of order $\eta^{-1}$. In the applications it is part of the adjacent scalar family ($DG'$, $G'D^T_X$, $DG'D^T_X$, $DP_\cV$, $P_\cV D^T_X$), whose compatibility with derivatives is Theorem~\Nref{thm:weak-inverse}(b) at (L1), (L2). Then the products $O_1^X\cdots O_n^X$, indexed as in Lemma~\Nref{lem:weak-chain}, and $\mathcal L^XO^X$, $O^X\mathcal L^X$ have compatible representations at $\kappa$, with:
\begin{itemize}
\item the factors evaluated at $\kappa+\delta$;
\item the rate $\gamma=\min_i\gamma_i$;
\item the constant $2n^2\prod_iC_d\delta^{-2}C_i$ ($2n-1$ positions of the distinguished factor, times the bound of Lemma~\Nref{lem:weak-chain}), times $(1+|\mathcal L|_{2\gamma})$ for each local factor;
\item the seam width $w$.
\end{itemize}
\begin{step}[Index sets]
The product index set consists of the tuples whose consecutive $r$-neighbourhoods meet. It is determined by the geometry, so it is the same on all domains, and matched tuples carry the same polymer and monomial. A tuple vanishes on $X$ if one of its terms does.
\end{step}
\begin{step}[Differences]
Apply the telescoping identity of Detail~\ref{d9:wi-b}, Step~1, with $A_i=O_i^X$ and $A'_i=O_i^{X'}$. The distinguished factor is of one of two kinds:
\begin{itemize}
\item $PO^X_iP-O^{X'}_i$, which is of seam form at the rate $\gamma_i$ by compatibility;
\item $PO_i^XQ$, with $Q=\one_X-P$ as in Detail~\ref{d9:wi-b}, whose input fields lie within one lattice step of $X\setminus X'$, hence in $\Sigma$, so its seam norm at $\gamma$ is at most its norm at $2\gamma$.
\end{itemize}
The neighbours are bounded at $2\gamma_i\ge2\gamma$; this is exactly the second condition of Definition~\Nref{def:compatible}. The seam product lemma (Detail~\ref{d9:seam}, Step~4) and the $2n-1$ positions give the bound. For a local factor, $P\mathcal L^XP-\mathcal L^{X'}$ is supported within $r_0$ of $X\setminus X'$, and so is $P\mathcal L^XQ$. Since $r_0\le w$, both are supported in $\Sigma=\Sigma_w(X,X')$, and their seam norms are bounded by their norms at $2\gamma$. So the seam width does not grow.
\end{step}
\begin{step}[Additive terms]
For $T\one_\Box$ split box by box (Section~\Nref{ss:40.2}), $P(\one_XT\one_X)\one_\Box P-\one_{X'}T\one_{X'}\one_\Box=0$ when $T$ is the identity. For a local $T$ of range $r_0\le w$ the difference is a cut term in $\Sigma_w$.
\end{step}
\status{New argument.}
\end{detail}

\begin{detail}[(L1), Section~\Nref{ss:41.3}]\label{d9:L1}\Ntarget{ss:41.3}
$\fA'^X=\one_X(D^TD+aQ'^TQ')\one_X$, $\mathfrak N=0$, no mass term. (W2) is Theorem~\Nref{thm:reg-scalar}(b),(e) on the domains $X_j$; these hold for every domain. The hypothesis of Lemma~\Nref{lem:step-small} is the commutator bound of the proof of Theorem~\Nref{thm:random-walk}(a). Its proof ``applies verbatim'' (Section~\Nref{ss:41.3}) with three changes:
\begin{enumerate}[label=(\roman*)]
\item $\Lam^w_j$ is replaced by $X_j$. The decay of $G'_{X_j}$, $DG'_{X_j}$, $G'_{X_j}D^T_{X_j}$ and $DG'_{X_j}D^T_{X_j}$ is Theorem~\Nref{thm:reg-scalar}(b) for the domain $X_j$.
\item $[\one_X\fA'\one_X,h_j]=\one_X[\fA',h_j]\one_X$, because $h_j$ and $\one_X$ are both multiplication operators.
\item $[D^TD,h_j]=D^T[D,h_j]+[D^T,h_j]D$. Here $[D,h_j]$ is multiplication by $\eta^{-1}(h_j(x+\eta e_\mu)-h_j(x))$, a function $\psi_\mu$ with $|\partial^k\psi_\mu|\le CM^{-1-k}$, composed with the parallel transport. By the discrete Leibniz rule, $D^T(\psi\,G'_j)$ is a sum of terms $\varphi'G'_j$ with $|\varphi'|\le CM^{-2}$ and $\varphi DG'_j$ with $|\varphi|\le CM^{-1}$, up to shifts by $\eta$. So no left $D^T$ remains ($\mathcal A=0$), and $|[\fA',h_j]G'_j|_{\gamma_0}+|[\fA',h_j]G'_jD^T|_{\gamma_0}\le CC_GM^{-1}$.
\end{enumerate}
The term $a[Q'^TQ',h_j]$ is block-local with norm $CM^{-1}$. The constants depend on $a$, $\alpha$ and $\dim\mathfrak g$ only.

The comparison of the local commutators assumed in Theorem~\Nref{thm:weak-inverse}(b) follows in the same way. Let $X'\subset X$ and $X'_j=\Lam^v_j\cap X'$. Since $h_j$ commutes with $\one_X$ and $\one_X\one_{X'}=\one_{X'}$,
\[
[\fA'^X,h_j]G'_{X_j}-[\fA'^{X'},h_j]G'_{X'_j}=\one_X[\fA',h_j](G'_{X_j}-G'_{X'_j})+(\one_X-\one_{X'})[\fA',h_j]G'_{X'_j},
\]
with $\fA'=D^TD+aQ'^TQ'$. By (iii) the first term consists of multiplications of size $CM^{-1}$ applied to $G'_{X_j}-G'_{X'_j}$ and to $D(G'_{X_j}-G'_{X'_j})$, up to shifts, and of $a[Q'^TQ',h_j]$ applied to the difference. Theorem~\Nref{thm:reg-scalar}(e) gives the seam form, for the zero extensions. The second term lives on the points of $X\setminus X'$, where $[\fA',h_j]G'_{X'_j}$ is bounded by (iii). The same holds with $D^T_{X_j}$, $D^T_{X'_j}$ on the right. So the comparison holds with constants $(CC_GM^{-1},\gamma_0)$.
\status{Literal repetition, with the changes (i)--(iii) checked.}
\end{detail}

\begin{detail}[(L2), Section~\Nref{ss:41.3}]\label{d9:L2}\Ntarget{ss:41.3}
The operators are $\fA^X=Q'G'^XQ'^T$ and $\fA^X=M'^X=Q'(G'^X)^2Q'^T$ on unit-lattice fields on $X_1$, with $\fA_{\rm loc}=0$, $\mathfrak N^X=\fA^X$ and no mass term.
\begin{itemize}
\item \emph{(W3).} The representations are products of the (L1) family with the local factors $Q'$, $Q'^T$ (Lemmas~\Nref{lem:weak-chain} and~\Nref{lem:weak-local-ops}(a)). They are compatible by Detail~\ref{d9:prod-compat}.
\item \emph{(W2).} The local propagators are $\Gamma'_{X_j}$ and $(M'_{X_j})^{-1}$. Their lower bounds, decay and comparison of domains are Theorem~\Nref{thm:reg-scalar}(c),(e).
\item \emph{Lemma~\Nref{lem:parametrix-X}.} It applies with $h_j$ restricted to unit points. The commutator bound of Lemma~\Nref{lem:decay-norm}(c) holds on the unit lattice. The local commutator is $0$, so the step hypothesis and the comparison of the local commutators of Theorem~\Nref{thm:weak-inverse}(b) hold with $C_{\rm loc}=0$.
\end{itemize}
The products are:
\begin{align*}
H'^X&=G'^XQ'^T\Gamma'^X, &P^X_\cV&=G'^XQ'^T(M'^X)^{-1}Q'G'^X, &R_X&=\one_X-P^X_\cV,\\
S_X&=G'^XR_XD^T_X, &\Pi_X&=\one_X-DS_X .
\end{align*}
The derivative versions follow from (L1) with derivatives. The identities in $R_X$ and $\Pi_X$ are additive terms, split box by box. All are compatible by Detail~\ref{d9:prod-compat}.
\status{Expanded.}
\end{detail}

\begin{detail}[(L3): the non-local part and the step hypothesis, Section~\Nref{ss:41.3}]\label{d9:L3}\Ntarget{ss:41.3}
\begin{step}[The formula for $\mathfrak N^X$]
By Section~\Nref{ss:19.1}, $\fA^X=\Pi_X^T\mathfrak B_X\Pi_X+DR_XD^T_X+a\Pi_X^T\bQ^T\bQ\Pi_X$, with $\mathfrak B_X=\one_X(D^TD+W)\one_X$ and $\Pi_X=\one_X-DS_X$. Then:
\begin{itemize}
\item Since $E=0$, $\bQ\Pi_X=\bQ$, because $\bQ DS_X=D_VQ'S_X=0$ by Theorem~\Nref{thm:reg-scalar}(d).
\item $\Pi_X^T=\one_X-S_X^TD^T_X$, and $S_X^TD^T_X\mathfrak B_X=(\mathfrak B_XDS_X)^T$ because $\mathfrak B_X^T=\mathfrak B_X$.
\item $DR_XD^T_X=DD^T_X-DP_{\cV,X}D^T_X$.
\end{itemize}
Hence $\fA^X$ is the sum of
\begin{gather*}
\one_X(D^TD+W)\one_X+DD^T_X+a\bQ^T\bQ\qquad\text{and}\\
\mathfrak N^X=-DP_{\cV,X}D^T_X-\mathfrak B_XDS_X-(\mathfrak B_XDS_X)^T+S_X^TD^T_X\mathfrak B_XDS_X .
\end{gather*}
With $\fA^{\rm ext}_X=D\one_{\Tf\setminus X}D^T$, the local part is $\one_X(D^TD+DD^T+W+a\bQ^T\bQ)\one_X$. On scalar fields in $\ell^2_0(X)$, $DD\lambda=\mathfrak F\lambda$, so $\mathfrak B_XD=\one_X(D^T\mathfrak F+WD)$. At points of $X$, $D^T_XD^T=\one_X\mathfrak F^T$ on $2$-forms (proof of Lemma~\Nref{lem:local-system}). This gives the four terms (i)--(iv) listed in the proof of Lemma~\Nref{lem:N-tier} (Detail~\ref{d9:N-tier}).
\end{step}
\begin{step}[(W3)]
The middle factors are products of the families of (L1) and (L2), with their derivatives $DG'$, $G'D^T_X$, $DG'D^T_X$, $DP_\cV$, $P_\cV D^T_X$, and of the bounded local factors $\one_X$, $\mathfrak F$, $\mathfrak F^T$ and $W$. A fine derivative is never a factor by itself. The local factors enter the plain norms through Lemma~\Nref{lem:weak-local-ops}(a). That lemma needs the block Schur norm $|\mathfrak F|_0\le C\sup|c_{\mu\nu}|\le C\eps$, not the operator norm. Compatibility follows from Theorem~\Nref{thm:weak-inverse}(b) at (L1), (L2) and Detail~\ref{d9:prod-compat}. A left gradient falls on $G'$ or $P_\cV$, never on $\mathfrak F$ or $W$.
\end{step}
\begin{step}[(W2)]
$\bG^{X_j}=(\fA^{X_j,\kappa_m})^{-1}$. Theorem~\Nref{thm:reg-gauge}(a) gives the decay of $P\bG^{X_j}Q$ for $P\in\{1,D,D^T,\nabla\}$ and $Q\in\{1,D^T,D\}$, where $D^T$ on the left is the torus divergence of the zero extension. Theorem~\Nref{thm:reg-locality}(a) gives the comparison of domains for the same list. The constants are uniform in $X$ and free of $M\ge M_0$. The operators depend on $U$ on $\hat X_j$.
\end{step}
\begin{step}[The step hypothesis]
$[\fA^X_{\rm loc},h_j]=\one_X[D^TD+DD^T+W+a\bQ^T\bQ,h_j]\one_X$. Applied to $G_j$:
\begin{itemize}
\item $D^T([D,h_j]G_j)$ is a left $D^T$ in front of an $\mathcal A$-part $[D,h_j]G_j$ of size $CM^{-1}C_G$.
\item $[D^T,h_j]DG_j$ has size $CM^{-1}C_G$.
\item $D([D^T,h_j]G_j)=D(\psi\,G_j)$. By the Leibniz rule this is a sum of $\varphi'G_j$ and $\varphi\nabla G_j$, with $|\varphi'|\le CM^{-2}$ and $|\varphi|\le CM^{-1}$. The full gradient is needed on $K_j^{(1)}$, which may meet $\partial X$. It is bounded up to $\partial X$ by Lemma~\Nref{lem:full-gradient}, through the third item of Theorem~\Nref{thm:reg-gauge}(a). This is where the exterior divergence term is essential (Remark~\Nref{rem:L3-domains}).
\item $[D,h_j]D^TG_j$: bounded by the decay of $D^T\bG^{X_j}$.
\item $[a\bQ^T\bQ,h_j]G_j$: $\bQ^T\bQ$ has range one block, so the size is $CM^{-1}$.
\item $[W,h_j]G_j$: $W$ has range one plaquette, so the size is $CM^{-1}\eps$.
\end{itemize}
With $D^T$ on the right, $G_j$ is replaced by $G_jD^T$; the same list of the theorem bounds it. At complex perturbations, Proposition~\Nref{prop:reg-complex} gives the same bounds with the constants of $U$, up to factors close to $1$. So $C_{\rm loc}=C(d,\dim\mathfrak g,a)$.

The comparison of the local commutators assumed in Theorem~\Nref{thm:weak-inverse}(b) holds with the same constant, enlarged by a fixed factor. Let $X'\subset X$, $X'_j=\Lam^v_j\cap X'$ and $\fA^{\Tf}_{\rm loc}=D^TD+DD^T+W+a\bQ^T\bQ$, the torus operator. Since $h_j$ commutes with $\one_X$ and $\one_X\one_{X'}=\one_{X'}$,
\[
[\fA^X_{\rm loc},h_j]G^{X_j}-[\fA^{X'}_{\rm loc},h_j]G^{X'_j}=\one_X[\fA^{\Tf}_{\rm loc},h_j]\big(G^{X_j}-G^{X'_j}\big)+(\one_X-\one_{X'})[\fA^{\Tf}_{\rm loc},h_j]G^{X'_j} .
\]
In the first term, by the Leibniz rule of the third item, applied also to $D^T[D,h_j]$, $[\fA^{\Tf}_{\rm loc},h_j]$ acts through $D$, the torus divergence $D^T$, $\nabla$, $\bQ$, $W$ and multiplications of size $CM^{-1}$. Theorem~\Nref{thm:reg-locality}(a) gives the comparison of the zero extensions of $\bG$, $D\bG$, $D^T\bG$, $\nabla\bG$ and of their versions with $D^T$ on the right, of seam form for $\Sigma_{10M}(X_j,X'_j)\subset\Sigma_w(X,X')$. In the second term $\one_X-\one_{X'}$ projects onto bonds within one lattice step of $X\setminus X'$, and $[\fA^{\Tf}_{\rm loc},h_j]G^{X'_j}$ is bounded by the items above, with $D^T[D,h_j]G^{X'_j}$ expanded by the Leibniz rule of the third item and bounded through the full gradient of $G^{X'_j}$. At complex perturbations Proposition~\Nref{prop:reg-complex} gives the same.
\end{step}
\status{Expanded.}
\end{detail}

\begin{detail}[(L4), and the repetitions of (L3), (L4) with $W_2$, Section~\Nref{ss:41.3}]\label{d9:L4}\Ntarget{ss:41.3}
\emph{The idea.} Level~(L4) applies the walk to $\bQ_X\bG^X\bQ_X^T$, whose local propagators are the Dirichlet operators of Part~VIII (Step~1). The repetitions add the potential $W_2$ of the coarse action, which depends on $X$ globally and is therefore put into the non-local part (Steps~2 and~3). The work is (W2) for the propagators with $W_2$, including the full gradient (Step~4); Steps~5--7 check the step hypothesis, the coarse-action parts and the complex references.
\begin{step}[(L4) with $W=0$]
$\fA^X=\bQ_X\bG^X\bQ_X^T$ on $\ell^2(\mathfrak B(X))$, with $\fA_{\rm loc}=0$ and $\mathfrak N^X=\fA^X$, so the step hypothesis and the comparison of the local commutators hold trivially. (W3) holds by (L3) and Detail~\ref{d9:prod-compat}, since $\bQ_X$ is a local factor of range one block. (W2) for $\boldsymbol\Gamma^{X_j}$ consists of:
\begin{itemize}
\item lower bound and decay: Proposition~\Nref{prop:dirichlet-averaged}(a);
\item comparison of domains: Proposition~\Nref{prop:dirichlet-averaged}(d), for coarse indices in $\mathfrak B(X'_j)$, which is the compressed form needed;
\item complex perturbations: Proposition~\Nref{prop:dirichlet-averaged}(e);
\item data dependence: on $\hat X_j=X_j^{(4)}$ (Proposition~\Nref{prop:dirichlet-averaged}(c); Part~VIII's hat, Section~\Nref{sec:dirichlet-averaged}, paragraph ``Conventions for Part VIII'').
\end{itemize}
This gives $\boldsymbol\Gamma^X(s)$, then $\bH^X(s)=\bG^X(s)\bQ_X^T\boldsymbol\Gamma^X(s)$, $\bC^X(s)$ by Section~\Nref{ss:35.1} with every factor weakened, and $\omega_0^X(s)=\bH_0^X(s)^T\one_XJ_U$.
\end{step}
\begin{step}[The family $W_2^X(s)$]
By Definition~\Nref{def:dirichlet-hessian}, $\ip A{W_2^XA}=d^2\mathcal N(0)[\one_XA,\one_XA]-\ip{\omega_0^X}{d^2\Cav_k(0)[A,A]}$.
\begin{itemize}
\item The first part is a sum of plaquette forms, an additive local term, split box by box.
\item For a term $\omega$ of the family $\omega_0^X(s)$, the operator $A\mapsto\sum_c\omega_{0,\omega}(c)\,\partial^2_A\Cav_k(0)(c)$ has a kernel supported in $\bigcup_c B(c_-)\cup B(c_+)$, by Proposition~\Nref{prop:Q-expansion}(a),(c). Its block Schur norm is at most $C\sup_c|\omega_{0,\omega}(c)|$.
\item At real references, $|\omega_{0,\omega}(c)|\le\sum_y\|\one_c\bH_{0,\omega}^T\one_y\|\,\|J_U\|_{B(y)}\le C\eps|\bH_{0,\omega}|_0$.
\item At complex references, $J_U=D^TF^s_U$ (Lemma~\Nref{lem:current}(a)), so $\omega_{0,\omega}(c)=\ip{D\bH_{0,\omega}\delta_c}{F^s_U}$, bounded by $C|D\bH_{0,\omega}|_0|F^s_U|_\infty$. The bound on $F^s$ at complex references is step~(v) of the proof of Proposition~\Nref{prop:complex-ext}.
\end{itemize}
So $\|W_2^X(s)\|_{\kappa,\gamma}\le Ce^\kappa\eps+C\eps\|\bH_0\|^D_{\kappa,\gamma}$.

For compatibility, two facts. First, $P\,d^2\mathcal N(0)|_X\,P=d^2\mathcal N(0)|_{X'}$ exactly, also box by box, because $\one_X\one_{X'}=\one_{X'}$ and $\one_\Box$ commutes with both. Second, the $\omega_0$-part on $\mathfrak B(X')$ is
\[
(P\bH_0^X\one_{\mathfrak B(X')}-\bH_0^{X'})^TJ_U+(Q\bH_0^X\one_{\mathfrak B(X')})^TJ_U .
\]
The first summand is of seam form on the fine side, which Proposition~\Nref{prop:dirichlet-averaged}(d) provides for all unit points; in weakened form this is Theorem~\Nref{thm:weak-inverse}(b) with Detail~\ref{d9:prod-compat}. The second is a cross term supported on the seam. This is Proposition~\Nref{prop:dirichlet-hessian}(d), term by term.
\end{step}
\begin{step}[The operator of the repetition]
With $W=W_2^X$ the local part keeps $W=0$; $W_2^X$ goes into the non-local part. By Step~1 of Detail~\ref{d9:L3}, the non-local part is $\mathfrak N^X[0]+\mathfrak N^X_{W}$ with
\[
\mathfrak N^X_W=\one_XW_2^X\one_X-W_2^XDS_X-(W_2^XDS_X)^T+S_X^TD_X^TW_2^XDS_X .
\]
These are factorized terms with $X_i^T=1$. The middle factors $M$ and right factors $Y$ are:
\begin{center}
\begin{tabular}{@{}ll@{}}
\toprule
$M$ & $Y$\\
\midrule
$W_2^X(s)$ & $1$\\
$W_2^X(s)\,(DG'R)(s)$ & $D^T_X$\\
$(DRG'D^T_X)(s)\,W_2^X(s)$ & $1$\\
$(DRG'D^T_X)(s)\,W_2^X(s)\,(DG'R)(s)$ & $D^T_X$\\
\bottomrule
\end{tabular}
\end{center}
Their plain norms and compatibility follow from Step~2 and Detail~\ref{d9:prod-compat}. (W1) holds, since $W_2^X$, which depends on $X$ globally, is not in the local part.
\end{step}
\begin{step}[(W2) for the propagators with $W_2^{X_j}$, real references]
$G_j=\bG^{X_j}[W_2^{X_j}]$. The potential $W_2^{X_j}$ has the following properties (Proposition~\Nref{prop:dirichlet-hessian}(b),(c)):
\begin{itemize}
\item it is real symmetric at real $U$;
\item it has range $\le3$ in $A$;
\item $\|W_2^{X_j}\|\le C\eps$;
\item it depends on $U$ on $\hat X_j$.
\end{itemize}
By Section~\Nref{sec:dirichlet-averaged}, paragraph ``Conventions for Part VIII'', the $L^2$ statements of Section~\Nref{sec:reg-gauge} hold verbatim for such potentials: coercivity in $\tn\cdot_{X,+}$ for $C\eps\le\frac12c_*$, decay of $P\bG Q$, and comparison of domains with the same potential.
\begin{itemize}
\item \emph{The full gradient.} In the proof of Theorem~\Nref{thm:reg-gauge}(a), the decay of $\nabla\bG^X$, $\nabla\bG^XD^T$ and $\nabla\bG^XD$ comes from Lemma~\Nref{lem:full-gradient} applied to $\varphi v$. That proof uses only the $\tn\cdot_{X,+}$-decay of $\bG^X$, $D\bG^X$ and $D^T\bG^X$, not the potential. So it holds with $W_2^{X_j}$, with the same constants.
\item \emph{Comparison of domains.} For $X'_j\subset X_j$, Proposition~\Nref{prop:dirichlet-averaged}(d) with $W\ne W'$ and $P\in\{1,D,D^T,\nabla\}$ gives the seam sum, plus
\[
\frac C\eps\sum_{w\in X'_{j,1}}e^{-\gamma(|y-w|+|w-y'|)}\|W_2^{X_j}-W_2^{X'_j}\|_w .
\]
The proof of Proposition~\Nref{prop:dirichlet-hessian}(d) gives $\|W_2^{X_j}-W_2^{X'_j}\|_w\le C\eps\sum_{v\in\Sigma}e^{-\gamma|w-v|}$. Since
\[
|y-v|+|v-y'|\le2\big(|y-w|+|w-y'|+|w-v|\big),
\]
we have $\sum_we^{-\gamma(|y-w|+|w-y'|)}e^{-\gamma|w-v|}\le S_{\gamma/2}\,e^{-\frac\gamma4(|y-v|+|v-y'|)}$. So the difference is of seam form at the rate $\gamma/4$, with a constant free of $\eps$ and of $M\ge M_0$. The seam $\Sigma_{10M}(X_j,X'_j)$ lies in $\Sigma_w(X,X')$.
\item \emph{Data dependence.} The operators depend on $U$ on $\hat X_j=X_j^{(4)}\subset(\Lam^v_j)^{(R_v)}$.
\end{itemize}
\end{step}
\begin{step}[The step hypothesis at the repetition]
The local part is that of (L3) with $W=0$. So the local commutator is that of Detail~\ref{d9:L3}, Step~4, now applied to the propagators with $W_2^{X_j}$. It needs $|G_j|^D$ and $\nabla G_j$, which Step~4 provides. $C_{\rm loc}$ is unchanged. The comparison of the local commutators follows as in Detail~\ref{d9:L3}, Step~4, with the comparison of domains of Step~4 above ($P\in\{1,D,D^T,\nabla\}$, $W\ne W'$) in place of Theorem~\Nref{thm:reg-locality}(a).
\end{step}
\begin{step}[The repetition of (L4) and the coarse-action parts]
$\boldsymbol\Gamma^X[W_2^X](s)$ comes from a walk with (W2) from Proposition~\Nref{prop:dirichlet-averaged}(a),(d) with the potentials $W_2^{X_j}$. Then $\Delta^{0,X}(s)=\boldsymbol\Gamma^X[W_2^X](s)-a$, with $-a$ split box by box.

By Section~\Nref{ss:37.4}, $\Delta^{\rm c,X}(b,b')$ is the sum of $\Delta^{0,X}(\hat Z^1\delta_b,\hat Z^1\delta_{b'})$ and the polarization of $2\ip{\omega_0^X}{\hat Z^2}$, when $\hat Z^1\delta_b,\hat Z^1\delta_{b'}$ are supported in $\mathfrak B(X)$. Here $\hat Z^1$ is linear of range $2L$ and $\hat Z^2$ is quadratic and local (Lemma~\Nref{lem:eliminate}). So $\Delta^{\rm c,X}(s)$ is obtained from the families $\Delta^{0,X}(s)$ and $\omega_0^X(s)$ by local factors of range $2L$ (Lemma~\Nref{lem:weak-local-ops}(a)), followed by a cut. Finally $\Delta^{\rm c}_Y(s)=\one_{\mathfrak F(Y)}\Delta^{\rm c,Y^+}(s)\one_{\mathfrak F(Y)}$.
\end{step}
\begin{step}[Complex references]
At complex perturbations $U_\delta$, (W2) for the propagators with $W_2^{X_j}$ uses the range-$3$ clause of Proposition~\Nref{prop:reg-complex} and Proposition~\Nref{prop:dirichlet-averaged}(e). Both require that $W_2^{X_j}$ extends holomorphically, with an $L^2$ bound uniform in $\eta$ and $\|W_2^{X_j}(U_\delta)-W_2^{X_j}(U)\|\le C\delta$.

This is proved in the proof of Proposition~\Nref{prop:dirichlet-hessian}(c), in the paragraph ``Complex perturbations in (c)''. The bound on $W_0$ of Lemma~\Nref{lem:div-curvature} needs only $|H_p-1|\le C\eta^2(\eps+\delta)$, transports bounded by $2$, and $|F^s|\le C(\eps+\delta)$; step~(v) of the proof of Proposition~\Nref{prop:complex-ext} gives these at $U_\delta$, using $|\partial\delta\!A|\le\delta$. A Schwarz lemma on the ball of radius $2\delta_*$, with $\delta_*$ fixed once there, gives $\|W_2^X(U_\delta)-W_2^X(U)\|\le C_W\delta$ with a fixed constant $C_W$. Its ingredients, listed in Detail~\ref{d8:A1}, are the coefficient bound of Lemma~\Nref{lem:div-curvature} (the computation of the proof of Proposition~\Nref{prop:action}(c) at $A=0$), Proposition~\Nref{prop:complex-ext}(c) with step~(v) of its proof, and Proposition~\Nref{prop:AB}. The chapter on Part~VIII expands it (Detail~\ref{d8:A1}), including the parameter condition of the doubled ball.

With that argument, the $\omega_0$-part changes by $C\delta$ by Step~2, and Steps~4--6 hold at $U_\delta$ with the constants of $U$ doubled.
\end{step}
\status{Expanded at real references; the full gradient is closed in Step~4. At complex references the step is the paragraph ``Complex perturbations in (c)'' of the proof of Proposition~\Nref{prop:dirichlet-hessian}(c), expanded in Detail~\ref{d8:A1}.}
\end{detail}

\begin{detail}[(L5), Section~\Nref{ss:41.3}]\label{d9:L5}\Ntarget{ss:41.3}
$\fA^X=\Delta_X+r$ on $\ell^2(\mathfrak F(X))$, with $\fA^X_{\rm loc}=r+K^{\rm gf}_X$ and $\mathfrak N^X=\Delta^{\rm c}_X(s)$.
\begin{step}[(W1)]
$K^{\rm gf}_X=\one_{\mathfrak F(X)}\mathfrak G_W^T\mathfrak G_W\one_{\mathfrak F(X)}$. Its kernel vanishes unless $b$ and $b'$ are inner bonds of a common block $B'(y)$, so its range is at most $L$ (Section~\Nref{ss:37.4}), and its entries are at most $C_L$ (Lemma~\Nref{lem:avg-transport}(c)). It depends on $X$ only through the compression, i.e.\ through cut terms. At complex references it changes by a block-local operator of size $C_L\delta_c$ (Lemma~\Nref{lem:avg-transport}(f)).
\end{step}
\begin{step}[(W3): compatibility at seam width $R_0+10M+2L$]
For $X'\subset X$,
\[
\one_{\mathfrak F(X')}\Delta^{\rm c}_X\one_{\mathfrak F(X')}-\Delta^{\rm c}_{X'}=\one_{\mathfrak F(X')}\big(\Delta^{\rm c,X^+}-\Delta^{\rm c,X'^+}\big)\one_{\mathfrak F(X')} .
\]
For $b\in\mathfrak F(X')$, $\hat Z^1\delta_b$ is supported within $2L$ of $b$, hence in $\mathfrak B(X'^+)$, since $R_0\ge2L+2$. So the cut in the definition of $\Delta^{\rm c}$ is inactive on both sides, and the difference is
\[
\hat Z^{1T}\big(\Delta^{0,X^+}-\Delta^{0,X'^+}\big)\hat Z^1+\text{(the $\omega_0$-difference contracted with $\hat Z^2$)} .
\]
\begin{itemize}
\item By Detail~\ref{d9:L4}, these differences are of seam form for $\Sigma_{10M}(X^+,X'^+)$.
\item A point within $R_0$ of $X$ but not within $R_0$ of $X'$ lies within $R_0$ of $X\setminus X'$. So $X^+\setminus X'^+\subset(X\setminus X')^{(R_0)}$, and $\Sigma_{10M}(X^+,X'^+)\subset\Sigma_{R_0+10M}(X,X')$.
\item The local maps $\hat Z^1$, $\hat Z^2$ enlarge the seam by $2L$.
\end{itemize}
So the seam width is $w=R_0+10M+2L$. Matched terms come from the matched index data on $X^+$ and $X'^+$.
\end{step}
\begin{step}[(W2)]
Proposition~\Nref{prop:local-fluct}:
\begin{itemize}
\item (d): $\Rea(\Delta_{X_j}+r)\ge\frac1{16}c_L+r$, and $(\Delta_{X_j}+r)^{-1}$ has decay $(C_L(\frac1{16}c_L+r)^{-1},\gamma_C)$ uniformly in $r$;
\item (c): the comparison of domains with width $R_0+10M+2L$;
\item (b): data on $\widehat{X_j^+}=(X_j^+)^{(4)}$. This lies within $R_0+4$ of $\Lam^v_j$, inside $(\Lam^v_j)^{(R_v)}$ because $R_v\ge R_0+4+2r_{\rm loc}$ (Detail~\ref{d9:fam-setting}, Step~4).
\end{itemize}
The lower bound is stated for the operator, not for its inverse (Section~\Nref{ss:41.3}). That is the form the walk needs: it gives the norm and the decay of the local propagators.
\end{step}
\begin{step}[The step hypothesis]
$[r,h_j]=0$. With $h_j$ acting through the initial points,
\[
[K^{\rm gf}_X,h_j](b,b')=K^{\rm gf}_X(b,b')\big(h_j(b'_-)-h_j(b_-)\big),\qquad|b_--b'_-|_\infty\le L .
\]
So $|h_j(b'_-)-h_j(b_-)|\le CLM^{-1}$. A row has at most $CL^d$ non-zero entries, so $|[K^{\rm gf}_X,h_j]|_{\gamma_0}\le C_LLM^{-1}$, with $e^{\gamma_0L}$ absorbed into $C_L$. Hence $C_{\rm loc}=C_LL$. No derivatives occur.

We check the comparison of the local commutators assumed in Theorem~\Nref{thm:weak-inverse}(b). Let $X'\subset X$, $X'_j=\Lam^v_j\cap X'$, $K^{\rm gf}_{\Tf}=\mathfrak G_W^T\mathfrak G_W$ and $G^Y=(\Delta_Y+r)^{-1}$. Since $h_j$ acts through the initial points, it commutes with $\one_{\mathfrak F(X)}$, and
\[
[K^{\rm gf}_X,h_j]G^{X_j}-[K^{\rm gf}_{X'},h_j]G^{X'_j}=\one_{\mathfrak F(X)}[K^{\rm gf}_{\Tf},h_j]\big(G^{X_j}-G^{X'_j}\big)+\big(\one_{\mathfrak F(X)}-\one_{\mathfrak F(X')}\big)[K^{\rm gf}_{\Tf},h_j]G^{X'_j} .
\]
$[K^{\rm gf}_{\Tf},h_j]$ has range $L$ and $|[K^{\rm gf}_{\Tf},h_j]|_{\gamma_0}\le C_LLM^{-1}$. By Step~3 (Proposition~\Nref{prop:local-fluct}(c),(d)) the difference $G^{X_j}-G^{X'_j}$ is of seam form on $\mathfrak F(X'_j)$, with the factor $(\frac1{16}c_L+r)^{-1}$; its other entries have an index on a bond starting in $X_j\setminus X'_j\subset X\setminus X'$, in the seam. The second term lives on the bonds starting in $X\setminus X'$. So the comparison holds with constants $C\,C_G(r)\,C_LLM^{-1}$, uniformly in $r$, and seam width $R_0+10M+2L$.
\end{step}
\begin{step}[Uniformity in $r$]
$C_G(r)=C_L(\frac1{16}c_L+r)^{-1}\le16C_L/c_L$, and $C_{\rm loc}$, $C_{\mathfrak N}$ do not depend on $r$. Choose $M_1$ for the worst case $C_G=16C_L/c_L$. The bound of Theorem~\Nref{thm:weak-inverse}(a) is proportional to $C_G(r)$, so $\|R_r\|_{\kappa,\gamma}\le C\,C_G(r)\le C'(c_L+r)^{-1}$. Since $\theta$ is linear in $C_G$ (Lemma~\Nref{lem:step-small}), also $\theta(r)\le C\,C_G(r)$, which is used in Detail~\ref{d9:cov-cor}.
\end{step}
\status{Expanded.}
\end{detail}

% IX.5
\subsection{Covariance, determinant and the one-loop source}\label{d9:cov}

\begin{detail}[Corollary~\Nref{cor:weak-covariance}]\label{d9:cov-cor}\Ntarget{cor:weak-covariance}
\begin{step}[The square root]
Let $A$ be a matrix with $\Rea\ip{\bar v}{Av}\ge c\|v\|^2$, $c>0$. Then:
\begin{itemize}
\item $\|(A+r)^{-1}\|\le(c+r)^{-1}$ for $r\ge0$;
\item $\frac1\pi\int_0^\infty r^{-1/2}(A+r)^{-1}dr$ converges absolutely and equals $A^{-1/2}$, the holomorphic functional calculus of the principal branch of $z^{-1/2}$ on $\{\Rea z>0\}$.
\end{itemize}
For the second point: the scalar identity $\frac1\pi\int_0^\infty r^{-1/2}(\lambda+r)^{-1}dr=\lambda^{-1/2}$ holds for $\lambda>0$, and both sides are holomorphic on $\Rea\lambda>0$. The matrix statement follows from the Riesz--Dunford integral, which commutes with the $r$-integral by Fubini. (This is the formula of Dimock's paper, (192)--(193), cited in Corollary~\Nref{cor:weak-covariance}.) We apply it with $A=\Delta_k$, using $\Rea\Delta_k\ge\frac18c_L$ (Lemma~\Nref{lem:coercive-general}).

The terms of $R_r(s)$ are indexed by an index set that does not depend on $r$: the data are geometric, and $r$ enters only the local part. Each term is continuous in $r$. By (L5) (Detail~\ref{d9:L5}, Step~5),
\[
\sum_\omega e^{\kappa'\deg\omega}\int_0^\infty r^{-1/2}|O_\omega(r)|_{\gamma'}\,dr\le C\int_0^\infty r^{-1/2}(c_L+r)^{-1}dr=C\pi c_L^{-1/2},
\]
anchored at any box. By Tonelli, the terms $\frac1\pi\int r^{-1/2}O_\omega(r)\,dr$, with the polymers and monomials of $R_r(s)$, form a weakenable family. By Fubini its value at $s=1$ is $\frac1\pi\int r^{-1/2}(\Delta_k+r)^{-1}dr=C_k^{1/2}$.
\end{step}
\begin{step}[Rate, margin, uniformity]
The loss $\kappa'=\kappa-C\delta$, $C\le20$, is the depth count of Detail~\ref{d9:budget}.

The margins grow only through local factors (Lemma~\Nref{lem:weak-local-ops}(a)), by their ranges: $1$ for $D$, $D^T$; one block for $\bQ$, $Q'$; $2L$ for $\hat Z^1$, $\hat Z^2$; $L$ for $K^{\rm gf}$. Products and walks keep the largest margin of their factors (Lemma~\Nref{lem:weak-chain}, Detail~\ref{d9:wi-a}, Step~5), and the propagators sit inside $X(j)$ with margin $0$. So the margin is at most $C(L)$, independent of $r$, $M$ and $R_0$, and at most $M/16$ once $M\ge16C(L)$.

The bounds are uniform over the class of references of Section~\Nref{sec:weakenable}, since each level's constants are.

Analyticity: each term is a finite combination of products and inverses of operators that are analytic in the perturbation (Propositions~\Nref{prop:reg-complex}, \Nref{prop:dirichlet-averaged}(e), \Nref{prop:local-fluct}(b)). An $r$-integral of analytic functions with an integrable uniform bound is analytic.

Reality: at real $U$ and real $s$, every term is built from real operators by sums, products and inverses at real $r$.
\end{step}
\begin{step}[The determinant densities]
By Proposition~\Nref{prop:det-densities}, $\ell_k(b)=\int_0^\infty\operatorname{tr}_{\mathfrak g}\big[\frac1{1+t}-(\Delta_k+t)^{-1}(b,b)\big]dt$. By (L5) at $r=t$ and $s=1$, $(\Delta_k+t)^{-1}=\sum_jh_j(\Delta_{\Lam^v_j}+t)^{-1}+R^{(\ge1)}_t$. The terms without steps have diagonal $(h_jG_j)(b,b)=h_j(b)G_j(b,b)$. Split $\frac1{1+t}=\sum_jh_j(b)\frac1{1+t}$, a finite sum over the $j$ with $b\in K_j$. Then
\[
\ell_k(b)=\sum_jh_j(b)\,\ell_{\Lam^v_j}(b)-\int_0^\infty\operatorname{tr}_{\mathfrak g}\big[R^{(\ge1)}_t(b,b)\big]dt .
\]
Each piece is integrable on its own. For $\ell_{\Lam^v_j}$ this is the proof of Proposition~\Nref{prop:det-local}; for $b$ with $h_j(b)\ne0$, $b$ is based in $K_j\subset\Lam^v_j$, so $b\in\mathfrak F(\Lam^v_j)$.
\end{step}
\begin{step}[The bound $\|R^{(\ge1)}_t\|\le C(c_L+t)^{-2}$]
Write $C_G(t)=C_L(\frac1{16}c_L+t)^{-1}\le C(c_L+t)^{-1}$. The terms with $n\ge1$ steps are bounded as in Detail~\ref{d9:wi-a}, Step~1, by
\[
\sum_{n\ge1}(n+1)C_d\delta^{-2}C_je^{(\kappa+\delta)\mathsf N_{\rm b}}C_G(t)\big(2C_d\delta^{-2}\theta(t)\big)^n\le C\,C_G(t)\,\theta(t) .
\]
Here we used $2C_d\delta^{-2}\theta(t)\le\frac12$. With $\theta(t)\le C\,C_G(t)$ (Lemma~\Nref{lem:step-small}), this is $\le C(c_L+t)^{-2}$. An earlier form of the lemma with $C_G^2$ would give the same, since $C_G(t)^2\le(16C_L/c_L)\,C_G(t)$ (Detail~\ref{d9:step-small}, Step~6).
\end{step}
\begin{step}[The family of numbers]
The terms are:
\begin{itemize}
\item $s^{X(j)}h_j(b)\ell_{\Lam^v_j}(b)$, with polymer $X(j)$;
\item $-\int_0^\infty\operatorname{tr}_{\mathfrak g}O_\omega(t)(b,b)\,dt$ for the terms $O_\omega(t)$ of $R^{(\ge1)}_t(s)$, with their polymers.
\end{itemize}
A term vanishes unless $b\in X_\omega^{(r)}$ (kernel support), and $|\operatorname{tr}_{\mathfrak g}O(b,b)|\le\dim\mathfrak g\,\|\one_bO\one_b\|\le C|O|_{\gamma'}$. So
\[
\|\ell_k(b)\|_{\kappa'}\le C_je^{\kappa'\mathsf N_{\rm b}}C_L+C\int_0^\infty(c_L+t)^{-2}dt\le C ,
\]
using $|\ell_{\Lam^v_j}(b)|\le C_L$. The bound is uniform in $b$ and in the reference. By Fubini, at $s=1$ the terms sum to $\ell_k(b)$.
\end{step}
\status{Expanded. It uses $\theta\le C\,C_G$ of Lemma~\Nref{lem:step-small}.}
\end{detail}

\begin{detail}[Scope, the first cube, the order of choices, Section~\Nref{ss:41.3}]\label{d9:scope}\Ntarget{ss:41.3}
Every term of a walk starts with $s^{X(j_0)}h_{j_0}G_{j_0}$, so $G^X(0)=0$, and the background dependence is cut together with the field. The rate conditions on $\kappa'=\kappa-C\delta$ (Remark~\Nref{rem:parameters}, Proposition~\Nref{prop:localized-pieces}) involve constants depending on $d$ only. So $\kappa$ may be chosen after $L$ and before $M$. This order is checked in the parameter ledger (Detail~\ref{det:ledger-rates} and Section~\ref{subsec:ledger-order}). Here we note one requirement: $\kappa'>3^d$ is needed for $\kappa_2-1>0$ (Remark~\Nref{rem:one-loop-polymer}).
\status{Expanded (the global order is in the parameter ledger).}
\end{detail}

\begin{detail}[Proposition~\Nref{prop:one-loop-polymer}]\label{d9:one-loop}\Ntarget{prop:one-loop-polymer}
\begin{step}[(a), real references]
By Proposition~\Nref{prop:det-densities}, $\log\det\Delta_k=\sum_{b\in\mathfrak F}\ell_k(b)$, at $U$ and at $1$. By Detail~\ref{d9:cov-cor}, $\ell_k(b)=\sum_{\omega\in\Omega_b}O_\omega$. Only terms with $b\in X^{(r)}_\omega$ occur, so $X_\omega$ contains one of the at most $2^d$ boxes within $r$ of $b$, and $\sum_\omega|O_\omega|\le2^dC$. The torus is finite, so $\sum_b\sum_\omega|O_\omega|<\infty$, and the terms may be grouped by $X^\sharp_\omega$. The index sets $\Omega_b$ are the same at $U$ and at $1$ (geometric), so the normalized differences $O_\omega(U)-O_\omega(1)$ make sense term by term.
\end{step}
\begin{step}[(a), complex references and the branch of $\log\det$]
On the convex ball $\mathcal B$, $\Rea\Delta_k\ge\frac18c_L$ (Lemma~\Nref{lem:coercive-general}). So ``$\log\det\Delta_k$'' is read as $\operatorname{Tr}\int_0^\infty\big(\frac1{1+t}-(\Delta_k+t)^{-1}\big)dt$ (Proposition~\Nref{prop:det-densities}(c)). This is the trace of the holomorphic functional calculus of the principal logarithm, $\sum_i\operatorname{Log}\lambda_i$. It is holomorphic on $\mathcal B$, has exponential $\det\Delta_k$, and is real at the real points. With this branch the identity holds on $\mathcal B$ in two ways:
\begin{itemize}
\item directly, since Detail~\ref{d9:cov-cor} holds at the complex references and Fubini applies;
\item or by the identity theorem, as the note says (step~(iii) of Proposition~\Nref{prop:complex-ext}): $\mathcal B\cap\R^N$ is a non-empty open subset of a totally real subspace of maximal dimension, and $\mathcal B$ is connected.
\end{itemize}
The same branch is used in Detail~\ref{xa:d-identity}.
\end{step}
\begin{step}[(b), the activity]
\begin{itemize}
\item \emph{Support.} $O_\omega$ depends on the reference only on $X^{(r)}_\omega\subset X^\sharp_\omega$, so $\mathfrak e(X,\cdot)$ depends only on the bonds of $X$.
\item \emph{Invariance.} $O_\omega$ is the trace of a diagonal block of a covariant kernel (conjugated by $R(u(b_-))$), or $h_j(b)\ell_{\Lam^v_j}(b)$ with $\ell$ invariant.
\item \emph{Analyticity.} The configurations of Definition~\Nref{def:bg-activity} are $e^{\ii\eta\delta\!A}U_r$ with $|\delta\!A|,|\partial\delta\!A|,[\partial\delta\!A]_\alpha\le\delta_a$. For $\eps_a\le\eps_{\max}$ they are complex references in the sense of Section~\Nref{sec:complex-references}, of size at most $C_{\rm cmp}\delta_a$: the size is the sum of the three seminorms, with covariant gradients, and the first item of that section compares it with the flat seminorms up to the factor $C_{\rm cmp}=4$ (Detail~\ref{xa:d-cmp}). The families are analytic there once this size is $<\delta_P$ (Proposition~\Nref{prop:complex-ext}). So the note requires $\eps_a\le\eps_{\max}$ and $C_{\rm cmp}\delta_a<\delta_P$; since $\eps_a$ and $\delta_a$ are existential, this costs nothing. The condition $2C_{\rm ext}c_a<\delta_a$ compares the cutoff extension with the separate seminorms, and needs no factor.
\item \emph{Normalization.} $\mathfrak e(X,1)=0$ by construction.
\end{itemize}
So $\mathfrak e(X,\cdot)$ is a normalized background activity at $X$. Remark~\Nref{rem:bg-to-graded} turns it into a graded activity on $\mathcal U(X;c_a,c_a,c_a)$ under three conditions: $c_a\le a_2$, $\max(1,B_5)c_a\le\eps_a$ and $2C_{\rm ext}c_a<\delta_a$. The note leaves $\eps_a$, $\delta_a$ existential (the proof of Proposition~\Nref{prop:one-loop-polymer}). At $k=0$ the pieces come from Proposition~\Nref{prop:first-step}(c), and the proof through the coercivity of Proposition~\Nref{prop:first-step}(b) forces $\eps_a$ below the first-step thresholds. So no text has to change. The choice of $\eps_a$, $\delta_a$ and $c_a$ in the global order is in the parameter ledger (Detail~\ref{det:ledger-eps-M}).
\end{step}
\begin{step}[(b), the bound]
Let $X^\sharp_\omega=X$. Every box of $X$ meets $X^{(r)}_\omega$ and touches a box of $X_\omega$, so $|X|\le3^d|X_\omega|\le3^d\deg\omega$. Fix a box $\Box\subset X_\omega\subset X$. Then
\[
\sum_{\omega\in\Omega_b:X^\sharp_\omega=X}|O_\omega|\le\sum_{\Box\subset X}\sum_{\omega\ni\Box}e^{\kappa'\deg\omega}|O_\omega|\,e^{-\kappa'3^{-d}|X|}\le|X|\,\|\ell_k(b)\|_{\kappa'}e^{-\kappa_2|X|} .
\]
The bonds $b\in\mathfrak F$ that contribute are based in $X$. There are at most $4M^d|X|$ of them (unit lattice, boxes of volume $\le M^d$). The two terms $O_\omega(U)$ and $O_\omega(1)$ give the factor $2$. So $\sup|\mathfrak e(X,\cdot)|\le CM^d|X|^2e^{-\kappa_2|X|}$.
\end{step}
\begin{step}[(c), the hypotheses of Lemma~\Nref{lem:quadratic-bound}]
Take $P=X$, $\theta_r=\delta_1=\delta_0=c_a\le\frac12c_{\rm q}$ and $\Lam=\Lam_R(z)$ with $z\in X$ and $R=\operatorname{diam}X+4$. Since $\operatorname{diam}X\le M|X|$ (a $*$-connected chain of boxes of side $\le M$), $R/(4M+2)\le\frac14|X|+1$. Hence:
\begin{itemize}
\item $\Lam\supset X^{(3)}$, and $2R=2\operatorname{diam}X+8\le n-8<n$ for $\operatorname{diam}X\le\frac n2-8$; also $n\ge n_0$ (torus hypotheses).
\item $X$ is a $*$-connected union of boxes, and $X^{(3)}$ is connected. Corner contacts are thickened to cubes of side $6$, so chains of parallel bonds exist inside $X^{(3)}$, and $D_N\le CM(1+|X|)$. This is the connectedness the lemma uses.
\item $B_P=3C_EB_5\max(1,R/(4M+2))^p\le C(1+|X|)^p$.
\item $c_*=\delta_0+(1+D_N)(|A'|+B_P\theta)\le\frac12c_{\rm q}+C(1+|X|)^{p+1}(|A'|+\theta)$.
\item $RB_P\theta\le C(1+|X|)^{p+1}\theta$.
\item $\mathfrak m_1\le C\big(\mathfrak m+|A'|^2+(1+|X|)^{p+1}|A'|\theta\big)$, since $c_*\le1$.
\item With $\delta_1=\delta_0=c_a$ and $|A'|\le1$,
\[
\mathfrak s\le C(1+D_N)\frac{\mathfrak m_1+B_P\theta}{c_a}\le C(1+|X|)^{p+2}\mathfrak t .
\]
\end{itemize}
Choose $C$ in the condition $C(1+|X|)^{q_1}(\mathfrak t+|A'|)\le1$ so large that it implies:
\begin{itemize}
\item $c_*\le c_{\rm q}$ and $RB_P\theta\le1$;
\item $|A'|<\frac12c_a=\frac12\delta_0$ and $\theta\le\mathfrak t<c_a=\theta_r$;
\item $\mathfrak s\le\frac12$ (note $q_1=2p+6\ge p+2$).
\end{itemize}
The lemma gives $|\mathfrak e(X,U)|\le2\mathfrak s^2\sup_{\mathcal U(X;c_a,c_a,c_a)}|\mathfrak e(X,\cdot)|$. By Step~4 this is at most $C(1+|X|)^{2p+4}|X|^2\mathfrak t^2e^{-\kappa_2|X|}\le C(1+|X|)^{q_1}\mathfrak t^2e^{-\kappa_2|X|}$.
\end{step}
\begin{step}[(d)]
The conditions of Definition~\Nref{def:graded} are monotone in $(\theta_r,\delta_1,\delta_0)$, so $\mathcal U(X;\theta_r,\delta_1,\delta_0)\subset\mathcal U(X;c_a,c_a,c_a)$ for $\theta_r,\delta_1,\delta_0\le c_a$. For a representation with $|F_{U_r}|,|J_{U_r}|<\theta_r$, we have $\mathfrak t\le\delta_1+\delta_0^2+\theta_r=\mathfrak t_0$ and $|A'|<\delta_0$. Two cases:
\begin{itemize}
\item If $C(1+|X|)^{q_1}(\mathfrak t_0+\delta_0)\le1$, then (c) gives $C\mathfrak t_0^2(1+|X|)^{q_1}e^{-\kappa_2|X|}\le C\mathfrak t_0^2e^{-(\kappa_2-1)|X|}$.
\item Otherwise $1+|X|>(C(\mathfrak t_0+\delta_0))^{-1/q_1}$. Since $|X|^2e^{-|X|/2}\le C$, (b) gives
\[
C|X|^2e^{-\kappa_2|X|}\le Ce^{-(\kappa_2-1)|X|}e^{-(1+|X|)/2}\le Ce^{-(\kappa_2-1)|X|}\exp\big(-c(\mathfrak t_0+\delta_0)^{-1/q_1}\big).
\]
\end{itemize}
For $\operatorname{diam}X>\frac n2-8$: $|X|\ge(\frac n2-8)/M\ge n/(4M)$ for $n\ge32$, and $|X|^2e^{-\kappa_2|X|}\le Ce^{-(\kappa_2-1)|X|}e^{-n/(8M)}$.
\end{step}
\status{Expanded. The branch of $\log\det$ at complex references is written in Step~2. The choice of $\eps_a,\delta_a,c_a$, also at $k=0$, is in the parameter ledger.}
\end{detail}

% IX.6
\subsection{Torus independence}\label{d9:torus}

These results are main line for Theorem~\Nref{thm:ol-uv} and, through Section~\Nref{sec:gen-locality}, for Proposition~\Nref{prop:beta-identification}.

\begin{detail}[The setting and Lemma~\Nref{lem:grid-transfer}]\label{d9:torus-setting}\Ntarget{lem:grid-transfer}
\begin{itemize}
\item \emph{Lattices.} In the units of step $k$, $\T^{(i)}=L^{i-k}\Z^4$ for $i\le k+1$. $M$ is a power of $L$ and $M>L$, so $M\Z^4\subset L^{i-k}\Z^4$. $LM\mid n$ by (T1), so all these lattices project to $\Tf[n]$.
\item \emph{$4R_{\rm g}<n$.} With $R_{\rm g}=(\rho+C_v+8)M$ and $n\ge8(\rho+C_v)M$, this is $\rho+C_v>8$, which holds for $\rho\ge30$.
\item \emph{Grid translations.} $\tau_a$, $a\in M\Z^4$, is an isometry that maps $L^{i-k}\Z^4$, $M\Z^4$ and lattice steps onto themselves. So (a) and (b) are immediate, with $\iota_x=\tau_a$ (``clear'', the proof of Lemma~\Nref{lem:grid-transfer}).
\item \emph{Transfers of lifts.} The proof uses only three inequalities, each checked: $|\tilde p-\tilde q|<\frac n2$ on $\tilde P$; $\frac n2+R_{\rm g}<n$; and $4R_{\rm g}<\min(n,n')$ for the isometry of the balls.
\item \emph{The consequence.} A rule of range $\le R_{\rm g}$ at a site based at $y$ sees only the ball about $y$. There $\iota_y$ carries the lattices, the grid, the distances and the domain: for $Y\subset P$ by (b), and for $Y=\Tf[n]$ because $\iota_y$ is onto.
\end{itemize}
\status{Expanded.}
\end{detail}

\begin{detail}[Lemma~\Nref{lem:local-inventory}]\label{d9:inventory}\Ntarget{lem:local-inventory}
We check that each ingredient is given by a local rule of range $\le R_{\rm g}$ (in the sense of Lemma~\Nref{lem:grid-transfer}).
\begin{step}[(a) Geometry]
\begin{itemize}
\item \emph{Blocks.} $L$ is odd, and $B_i(y)$ is the cube of $L^d$ points of $\T^{(i-1)}$ centred at $y\in\T^{(i)}$. So the blocks nest exactly: $\bigcup_{x\in B_i(y)}B_{i-1}(x)$ is the cube of $L^{2d}$ points of $\T^{(i-2)}$ centred at $y$. Iterating down to the fine lattice gives $B(y)$, the fine points within $\frac12(1-\eta)$ of $y$.
\item \emph{Contours and bonds.} The contours of $G(y,x)$ are coordinate paths in the box spanned by $y$ and $x$. The crossing bonds and the bonds $b_0(c)$ are defined by membership in blocks $B'$ and by the rule $b_--\frac12(L-1)e_\mu\in\T'$. All have range $\le L$.
\item \emph{Paving.} The boxes are the grid cubes, the unions of the blocks $B(y)$ with $|y-c|_\infty\le\frac12(M-1)$, $c$ in the grid. Polymers, $|X|$, $X^{(r)}$ and $X^\sharp$ are defined from boxes and distances. The additive split uses boxes and base points.
\item \emph{Partition of unity.} $M\mid n$ gives $\check M=M$, so the $z_j$ are the grid points. For $n\ge8M$ the periodization of $\phi_j$ does not overlap. At $y$ only the $z_i$ with $|y-z_i|_\infty<M$ have $\phi_i(y)\ne0$, so $\chi_j(y)=\phi_j(y)(\sum_i\phi_i(y)^2)^{-1/2}$ is a rule of range $M$.
\item \emph{Sets attached to $j$.} $K_j$, $\Lam_j$ and $\Lam^v_j$ lie within $\rho M+1$ of $z_j$, and $X(j)$ within $\rho M+1+R_v+M\le(\rho+C_v+2)M$.
\item \emph{Mass profile.} $m_Y(y)=\sum_{i\notin J(Y)}\chi_i(y)^2$ uses only the $i$ with $|z_i-y|<M$. Their $\Lam_i$ lie within $4M+1$ of $y$, so this is a rule of range $4M+1$ in $Y$.
\item \emph{Domain sets.} $Y_1$, $\mathfrak B(Y)$, $\hat Y$ (either hat), $\T'(Y)$, $\mathfrak F(Y)$ and $Y^+$ use $Y$ within distance $R_0+L\le R_{\rm g}$.
\end{itemize}
\end{step}
\begin{step}[(b) Constants]
None of the choices in Remarks~\Nref{rem:parameters}, \Nref{rem:domain-constants} and Sections~\Nref{sec:weakenable}, \Nref{sec:nested-walks} involves $n$; the order of choices is that of the parameter ledger (Section~\ref{subsec:ledger-order}). The IMS scale $M'$ and the local gauges occur only in proofs of lower bounds. The operators are defined by Section~\Nref{sec:reg-gauge}'s formulas, which contain neither.
\end{step}
\begin{step}[(c) Local operators, with ranges]
\begin{itemize}
\item $D$, $D^T$, $D^T_Y$, $\one_Y$, $\one_\Box$, $\fA^{\rm ext}_Y$ and $m_Y$ are explicit in $U$ at the site and in $Y$ near it.
\item $W_j$ on a bond $c$ of $\T^{(j)}$. By Lemma~\Nref{lem:one-step-basic}(c), the one-step average of $c$ depends on $W_{j-1}$ on the bonds of $B_j(c_-)\cup B_j(c_+)$. By the nesting of Step~1, iterating gives dependence on $U$ only on the bonds of $B(c_-)\cup B(c_+)$, in the units of the current scale. This is the claim in the proof of Lemma~\Nref{lem:local-inventory}.
\item The conditions of Definition~\Nref{def:avg-domain} are imposed contour by contour and bond by bond inside $B_j(y)$ and $B_j(c_\pm)$, so they are local.
\item At complex references the same recursion is used (Proposition~\Nref{prop:AB}(i)), so the continued $W_j$ have the same dependence region. The same holds for every local operator of this list, by a general fact. Let $f$ be holomorphic on a convex domain of complex references whose real points form a non-empty open subset of a real form, and let $f$ depend, at the real points, only on the variables of a set $S$. Its partial derivatives in the other variables are holomorphic and vanish at the real points, a totally real set of full dimension. So they vanish everywhere, by the identity theorem. The fibres of a convex domain are convex, hence connected, so $f$ depends only on the variables of $S$ everywhere.
\item $\bQ_k$, $Q'_k$ (Proposition~\Nref{prop:linearized-averages}(d)): range one block.
\item $\mathfrak F$ (Lemma~\Nref{lem:curvature}): multiplication with a shift by $2\eta$.
\item $F_U$, $F^s_U$, $J_U$: plaquette-local.
\item $d^2\mathcal N(0)$: plaquette forms (proof of Lemma~\Nref{lem:hessian-gamma}).
\item $d^2\Cav_k(0)[A,A](c)$: uses $B(c_-)\cup B(c_+)$ (Proposition~\Nref{prop:Q-expansion}(a)).
\item $\ell_c$, $\zeta_c$: use $P_c$. $\hat Z^1\delta_b$, $\hat Z^2$: support and dependence within $2L$ (Lemma~\Nref{lem:eliminate}).
\item $\mathfrak G_WZ(y,x)$: uses the box spanned by $y$ and $x$, inside $B'(y)$ (Lemma~\Nref{lem:avg-transport}(a)). $K^Y$, $K^{\rm gf}_Y$: range $L$.
\item At $k=0$: $\Delta^0$ is plaquette-local, and $\Delta_0$ has range $2L+2$ (Proposition~\Nref{prop:first-step}(b)).
\end{itemize}
The maximum range is $4L+2\le R_{\rm g}$.
\end{step}
\begin{step}[(d) Jacobians]
$\mathfrak r_c(0)$ is the derivative of $\ell_c$ at $Z=0$ with respect to $Z(b_0(c))$, at the base (Proposition~\Nref{prop:jacobians}(b), Lemma~\Nref{lem:eliminate}(b)). By (c) it is a formula in $U$ on the bonds of $P_c$, and $\log\det(1+\mathfrak r_c(0))$ is a finite-dimensional function of it.
\end{step}
\begin{step}[The consequence]
The kernel of $\mathcal L\one_S$ is supported on $S^{(r_{\mathcal L})}\times S$. With the convention of (c), its entry at $(x,x')$ uses the configuration within $r_{\mathcal L}$ of the input site $x'\in S$, so it needs $S^{(r_{\mathcal L})}\subset P$. For $\one_S\mathcal L$ the input sites lie in $S^{(r_{\mathcal L})}$ and their balls in $S^{(2r_{\mathcal L})}$. This is why the margin in Lemma~\Nref{lem:weak-local-ops}(a) is $r+2r_0$ (Detail~\ref{d9:local-ops}).
\end{step}
\status{Expanded.}
\end{detail}

\begin{detail}[Remark~\Nref{rem:global-objects}]\label{d9:global}\Ntarget{rem:global-objects}
\begin{itemize}
\item The torus inverses enter
\[
\mathfrak e(X,U)=-\tfrac12\sum\big(O_\omega(U)-O_\omega(1)\big)
\]
only through their terms. The normalizations $\mathcal N_k$, $N_\chi$ and $\sum_b\ell_k(b;1)$ are not of this form.
\item The value at a point of a graded domain does not depend on the cutoff extension, since $\mathfrak e(X,\cdot)$ depends only on the bonds of $X$ (Detail~\ref{d9:one-loop}, Step~3).
\end{itemize}
Remark~\Nref{rem:one-loop-local-use} is treated in Detail~\ref{d9:ol-use}.
\status{Expanded.}
\end{detail}

\begin{detail}[Lemma~\Nref{lem:torus-transfer}]\label{d9:transfer}\Ntarget{lem:torus-transfer}
$\iota_*$ is conjugation by the relabelling of the sites based in $P$ (fibres identified). So it commutes with:
\begin{itemize}
\item sums, products and transposes;
\item inverses on the spaces of sites based in $P$ (finite-dimensional);
\item absolutely convergent integrals in a parameter;
\item diagonal traces.
\end{itemize}
\begin{step}[Domains]
Let $X(j)\subset P$, so $z_j\in P$. Apply Lemma~\Nref{lem:grid-transfer}(b) at $x=z_j$. The sets $\Lam^v_j$, $(\Lam^v_j)^{(R_v)}$ and $X(j)$ are defined by rules of range $\le R_{\rm g}$ about $z_j$ and lie in $P$. So $\iota_{z_j}=\iota$ on them, and $\iota\Lam^v_j=\Lam^v_{\iota j}$, $\iota X(j)=X(\iota j)$. For related $X,X'$ we have $\iota(\Lam^v_j\cap X)=\Lam^v_{\iota j}\cap X'$: either $X\subset P$ and $X'=\iota X$, or both are the whole tori. If $Y,Y^{(R_0)}\subset P$, then $\iota(Y^+)=(\iota Y)^+$ by Lemma~\Nref{lem:grid-transfer}(a), since $R_0\le R_{\rm g}$.
\end{step}
\begin{step}[Dirichlet operators]
The operators on $Y$ and $Y^+$ are built from the local operators of Detail~\ref{d9:inventory}, (c), by the finite algebra listed in Step~2 of the proof of Lemma~\Nref{lem:torus-transfer}, in that order. By the consequence of Lemma~\Nref{lem:local-inventory}, they are carried as soon as $(Y^+)^{(2r_{\mathcal L})}=Y^{(R_0+2r_{\mathcal L})}\subset P$. This holds since $R_v\ge R_0+4+2r_{\rm loc}$ and $Y^{(R_v)}\subset P$. Their dependence regions $\hat Y$ and $\widehat{Y^+}=Y^{(R_0+4)}$ (Part~VIII hat) lie in $Y^{(R_v)}$ as well (Detail~\ref{d9:fam-setting}, Step~4). The inverses exist on both tori, so they are carried. So is the $t$-integral defining $\log\Delta_Y$, together with its diagonal trace $\ell_Y(b)$.
\end{step}
\begin{step}[One walk]
A term of $G^X(s)$ is determined by its index data: $j_0,\dots,j_m$, the kinds, and the terms $\omega_i$ of $\mathfrak N$. The condition that consecutive neighbourhoods meet is geometric. Let $X_\Omega^{(r)}\subset P$.
\begin{itemize}
\item All factor polymers lie in $X_\Omega$, and their margins are at most $r$, since margins only grow under products, local factors and walks. So $X(j_i)\subset P$ and $X_{\omega_i}^{(r)}\subset P$.
\item The image term $\iota(\Omega)$ has data $\iota j_i$, the same kinds, and $\iota(\omega_i)$, which exist by the induction hypothesis on the levels.
\item Its factors are carried: $h_{\iota j}$ and the local parts by (c), with ranges $<R_v$; $G^{X'}_{\iota j}$ by Step~2, since $(\Lam^v_j)^{(R_v)}\subset X(j)\subset P$; the $\mathfrak N$-terms by hypothesis.
\end{itemize}
So $O'_{\iota(\Omega)}=\iota_*O_\Omega$. The index sets on the two tori differ; the bijection is only between terms with $X^{(r)}\subset P$ and $X^{(r)}\subset P'$. It is a bijection because the same construction for $\iota^{-1}$ (a grid transfer with vector $-a$) inverts it.
\end{step}
\begin{step}[Products, additive terms, levels, Corollary~\Nref{cor:weak-covariance}, margin $0$]
Products are indexed geometrically and carried factor by factor. Additive terms $T\one_\Box$ with $\Box^{(r)}\subset P$ are carried by (c). The $r$- and $t$-integrals have index sets independent of the parameter.

The induction runs over the levels in the order of Section~\Nref{sec:nested-walks}. At (L5) the seam terms use $\Delta^{\rm c}_{\Lam^v_j}(s)$ on $(\Lam^v_j)^+$, which are related domains by Step~1.

Margin $0$: $X^\sharp_\omega\subset P$ implies $X^{(r)}_\omega\subset P$. A box of $\Tf[n']$ meeting $X^{(r)}_{\iota(\omega)}=\iota(X^{(r)}_\omega)$ is the image of the box containing the corresponding point. So $X^\sharp_{\iota(\omega)}=\iota(X^\sharp_\omega)$.

$k=0$: one walk with $\fA^X=\one_X(\Delta_0+r)\one_X$ and $\mathfrak N=0$, with $\Delta_0$ local by (c).
\end{step}
\begin{step}[Complex flat references]
Section~\Nref{sec:gen-locality} applies the lemma at the references $U'=e^{\ii\eta\chi B}$ of Lemma~\Nref{lem:gl-flat}(b). Their real part is $1$, which is $(M,0)$-regular with $F=J=0$. They lie in the class of Section~\Nref{sec:weakenable} under two conditions:
\begin{itemize}
\item \emph{The two sizes are comparable.} At $U_r=1$ the size of a complex reference (Section~\Nref{sec:complex-references}, paragraph ``Complex references'') and the $\delta$ of Proposition~\Nref{prop:reg-complex} (Section~\Nref{subsec:reg-complex}) measure the same field $\delta\!A=A'=\chi B$, with ordinary derivatives in both, since the local gauge of $1$ is trivial. The H\"older quotient of Section~\Nref{sec:complex-references}, paragraph ``Complex references'' is taken on the cubes $\Lam_1(c_y)$, while Proposition~\Nref{prop:reg-complex} uses the cubes $\Lam_3(z')$ of side $6$ of the local gauges (Definition~\Nref{def:regular}). Pairs of points at distance at most $1$ lie in a common cube $\Lam_1(c_{y'})$. Pairs at distance in $(1,6]$ have H\"older quotient at most $2|\partial A'|$. So $\delta_{\rm reg}<2\delta_c$: at $U_r=1$ the comparison constant is at most $2$, within the constant $C_{\rm cmp}=4$ of Section~\Nref{sec:complex-references}, paragraph ``Complex references'' (Detail~\ref{xa:d-cmp}).
\item \emph{The size is small.} We need $\delta_c<\delta_P$. This follows from (DC4) and Remark~\Nref{rem:domain-constants}(ii), which bound the size of $\chi B$ by $2C_{\rm ext}(L^{-1}\delta_0+L^{-2}\delta_1)$ (Remark~\Nref{rem:domain-constants}(ii)).
\end{itemize}
Under these conditions every operator of the construction exists at $U'$ (Propositions~\Nref{prop:reg-complex}, \Nref{prop:complex-ext}). The local operators are the holomorphic continuations of the same formulas (Lemma~\Nref{lem:local-inventory}(c), Proposition~\Nref{prop:AB}(i)). So their dependence regions are those at real references (Detail~\ref{d9:inventory}, Step~3), and Steps~1--4 use nothing else. So the lemma holds literally at the complex flat references. For a general complex reference, the comparison constant is $C_{\rm cmp}=4$ (Section~\Nref{sec:complex-references}, paragraph ``Complex references'', Detail~\ref{xa:d-cmp}). With the smallness of $\delta_P$ (Proposition~\Nref{prop:complex-ext}; Detail~\ref{xa:d-deltaP}), the perturbations of size $<C_{\rm cmp}\delta_P$ lie within the radii of Proposition~\Nref{prop:reg-complex}, and the same argument applies.
\end{step}
\status{Expanded. At the complex flat references the lemma holds literally; the comparison constant there is at most $2$, and the size bound is that of Remark~\Nref{rem:domain-constants}(ii).}
\end{detail}

\begin{detail}[Proposition~\Nref{prop:one-loop-local}]\label{d9:one-loop-local}\Ntarget{prop:one-loop-local}
\begin{step}[(a)]
Since $\operatorname{diam}\tilde X<\frac n2$, $\pi_n$ is injective and isometric on $\tilde X$. Since $M\mid n$, it maps grid cubes to distinct grid cubes and $*$-adjacency to $*$-adjacency. So $X$ is a polymer with $|X|=|\tilde X|$ that does not wrap. Conversely, a non-wrapping $X$ has a lift of $X^{(3)}$ (``Lifts and linear fields'', in Section~\Nref{sec:generations}), and the preimage $\tilde X$ is $*$-connected.

The cutoff extensions $\tilde U,\tilde U'$ (Remark~\Nref{rem:bg-to-graded}) are references of the setting (proof of Proposition~\Nref{prop:one-loop-polymer}(b)). They agree with $U,U'$ on the bonds of $X$, $X'$. So $\mathfrak e[n](X,U)=\mathfrak e[n](X,\tilde U)$ (Detail~\ref{d9:global}). A term with $X^\sharp_\omega=X$ has $b\in X^{(r)}_\omega\subset X$. Lemma~\Nref{lem:torus-transfer}(b),(c) with $P=X$, at $(\tilde U,\tilde U')$ and at $(1,1)$, gives a bijection $(b,\omega)\mapsto(\iota b,\iota(\omega))$ onto the pairs with $X^\sharp=X'$, with equal values. Here $b\mapsto\iota b$ maps $\mathfrak F$ to $\mathfrak F'$, since the rule for $b_0(c)$ is local. The sums converge absolutely (Detail~\ref{d9:one-loop}, Step~1), so they are equal.
\end{step}
\begin{step}[(b), (c)]
(b) is (a) for $\tau_a$:
\begin{itemize}
\item $\tau_a$ maps $\mathcal U(X)$ onto $\mathcal U(X+a)$;
\item translates of cutoff extensions are cutoff extensions;
\item $X^\sharp_{\tau_a(\omega)}=X^\sharp_\omega+a$.
\end{itemize}
(c): $\log\mathfrak j_c$ is the formula of Lemma~\Nref{lem:local-inventory}(d) in the configuration on $P_c$. The lift of $P_c$ has diameter $<2L+1<\frac12\min(n,n')$, and the consequence of Lemma~\Nref{lem:grid-transfer} applies.
\end{step}
\begin{step}[(d)]
\begin{itemize}
\item \emph{$n$-independence.} $V_B[n]$ and $V_B[n']$ agree on $\pi_n\tilde X$ under the transfer with $a=0$. Apply (a).
\item \emph{Holomorphy.} $B\mapsto V_B[n]$ is holomorphic, and $\mathfrak e(X,\cdot)$ is holomorphic on a neighbourhood of $\mathcal U(X)|_X$.
\item \emph{Vanishing.} $\mathfrak e(X,1)=0$ gives vanishing at $B=0$.
\item \emph{Reality.} This is the argument of Lemma~\Nref{lem:raw-new}(a): the terms are real at real references (Detail~\ref{d9:cov-cor}, Step~2).
\item \emph{Gauge invariance.} Let $e^{\ii\eta B'}=e^{\ii\lambda(b_-)}e^{\ii\eta B}e^{-\ii\lambda(b_+)}$ on the bonds of $\tilde X$. Extend $\lambda\circ\pi_n^{-1}$ to a real function on $\Tf[n]$ and put $u=e^{\ii\lambda}$. Then $V_B^u$ agrees with $V_{B'}$ on the bonds of $\pi_n\tilde X$, and $V_B^u\in\mathcal U(\pi_n\tilde X)$ by the real gauge invariance of the graded domains. Then use the invariance of $\mathfrak e$ and its dependence on the bonds of $X$ only.
\item \emph{Covariance} follows from (b), since $V_{\tau_aB}=\tau_aV_B$.
\item \emph{The bound} is Lemma~\Nref{lem:raw-new}(b).
\end{itemize}
\emph{The Jacobians.} Let $\tilde c$ be a positively oriented bond of $L\Z^4$, $\tilde P$ the lift of $P_c$ that contains it, and $B$ a field on the bonds of $\tilde P^{(3)}$ with the bounds of $\mathcal U_{\rm fl}$. The note says that the same arguments apply, with (c) in place of (a), (b). Written out:
\begin{itemize}
\item \emph{Domain.} As for $\tilde X$, $V_B[n]\in\mathcal U(P_c;c_a,c_a,c_a)$, since $\operatorname{diam}\tilde P^{(3)}<\frac12n$ and the real part is $1$.
\item \emph{$n$-independence.} $V_B[n]$ and $V_B[n']$ agree on $P_c$ under the transfer of $\tilde P$ with $a=0$. Apply (c).
\item \emph{Holomorphy and dependence.} $\log\mathfrak j_c$ is a normalized graded activity at $B'(c_-)\cup B'(c_+)$ on $\mathcal U(\cdot;c_a,c_a,c_a)$ (Proposition~\Nref{prop:jacobians}(c)), and a formula in the configuration on $P_c$ (Lemma~\Nref{lem:local-inventory}(d)). The map $B\mapsto V_B[n]$ is holomorphic.
\item \emph{Vanishing.} At $B=0$ the configuration is $1$, the base is trivial, and $\mathfrak j_c=1$ (Proposition~\Nref{prop:jacobians}(b)).
\item \emph{Reality.} At real $B$ the reference is real, and $\log\mathfrak j_c$ is real there (Proposition~\Nref{prop:jacobians}(a),(b)).
\item \emph{Gauge invariance.} As for $\mathfrak e$: extend $\lambda$ from $\pi_n\tilde P$ to a real function on $\Tf[n]$, and use the invariance of $\mathfrak j_c$ under real gauge transformations (Proposition~\Nref{prop:jacobians}(b)) and its dependence on $P_c$ only.
\item \emph{Covariance.} $J^\infty_k(\tilde c+a,\tau_aB)=J^\infty_k(\tilde c,B)$ is the last claim of (c), with $V_{\tau_aB}=\tau_aV_B$.
\item \emph{The bound.} $|J^\infty_k|\le C_{\rm jac}$ is Lemma~\Nref{lem:raw-new}(b), i.e.\ the bound $|\log\mathfrak j_c|\le C$ of Proposition~\Nref{prop:jacobians}(c).
\end{itemize}
\end{step}
\status{Expanded.}
\end{detail}

\begin{detail}[Remark~\Nref{rem:one-loop-local-use}(a)]\label{d9:ol-use}\Ntarget{rem:one-loop-local-use}
\begin{step}[The size check]
Let $Y'$ be a polymer of scale $k+1$ that is not wrapping, and $V_B\in\mathcal U^\flat_1(Y')$. So, in the units of scale $k+1$, $|B|<\delta_0$ and $|\nabla B|+[\nabla B]_\alpha<\delta_1^\flat=\frac14\delta_1$ on $Y'^{(3)}$, with the real part $1$ (Section~\Nref{subsec:ol-period-vacpol}, paragraph ``Flat domains and the infinite lattice''; $\delta_1^\flat=\frac14\delta_1$ by the Setting of Section~\Nref{sec:flow4}). Read $B$ in the units of step $k$. By Remark~\Nref{rem:domain-constants}(i) with one scale, $B$ is multiplied by $L^{-1}$, $\nabla B$ by $L^{-2}$, and $[\nabla B]_\alpha$ by $L^{-2-\alpha}$. Let $X$ be a polymer of scale $k$ with $\bar X=Y'$, or $X=P_c$ with $Y'_c=Y'$. Then $X\subset Y'$ and $X^{(3)}\subset Y'^{(3)}$. By the cube argument of Remark~\Nref{rem:domain-constants}(i), each cube $\Lam_3(z')$ of scale $k$ in $X^{(3)}$ lies in a cube $\Lam_3(z'')\subset Y'^{(3)}$ of scale $k+1$. So, on $X^{(3)}$ and in the units of step $k$,
\[
|B|<L^{-1}\delta_0\le c_a,\qquad|\nabla B|+[\nabla B]_\alpha<\tfrac14L^{-2}\delta_1\le c_a,
\]
by (DC1), $\delta_1\le\delta_0\le c_a$. The polymer $X$ does not wrap, since $Y'$ does not. Inside the lift of $Y'^{(3)}$ on which $B$ is given, the preimage of $X$ is a lift $\tilde X$, and $\tilde X^{(3)}$ lies in that lift. So the restriction of $B$ to the bonds of $\tilde X^{(3)}$ lies in $\mathcal U_{\rm fl}(\tilde X)$, and $V_B$ agrees on $X$ with the configuration $V_{B|}[n]$ of Proposition~\Nref{prop:one-loop-local}(d). Moreover $V_B\in\mathcal U(X;L^{-2}\theta_r,L^{-2}\delta_1,L^{-1}\delta_0)\subset\mathcal U(X;c_a,c_a,c_a)$, by Remark~\Nref{rem:domain-constants}(i) and (DC1). The same holds for $P_c$ and $\tilde P_{\tilde c}$.
\end{step}
\begin{step}[The conclusion]
By the last claim of Proposition~\Nref{prop:one-loop-local}(d), $\mathfrak e(X,V_B)=E^\infty_k(\tilde X,B|_{\tilde X^{(3)}})$, and in the same way $\log\mathfrak j_c(V_B)=J^\infty_k(\tilde c,B|_{\tilde P^{(3)}_{\tilde c}})$. Summing over the $X$ with $\bar X=Y'$ and the $c$ with $Y'_c=Y'$, $E'_{\rm ol}(Y',V_B)$ is a sum of values of $E^\infty_k$ and $J^\infty_k$ at lifts inside the lift of $Y'$. These values do not depend on the torus. So $E'_{\rm ol}(Y',\cdot)$ is the same function on every torus on which $Y'$ does not wrap, with no error term.

Grid covariance (G2) at scale $k+1$: a translation by $M\Z^4$ in the units of scale $k+1$ is a translation by $LM\Z^4\subset M\Z^4$ in those of step $k$. It maps the grids of both scales onto themselves, so $\overline{X+a}=\bar X+a$ and $Y'_{c+a}=Y'_c+a$. Proposition~\Nref{prop:one-loop-local}(b),(c) then give (G2).
\end{step}
\status{Expanded.}
\end{detail}

% IX.7
\subsection{Interior estimates with \texorpdfstring{$L^p$}{Lp} data}\label{d9:lp}

Throughout, $d=4$, $\|h\|_{p,Q}=(\eta^d\sum_Q|h|^p)^{1/p}$, $p^*=dp/(d-p)$ and $\beta_0=2-d/p$, so $1-d/p^*=\beta_0$ and $d/p-d/p^*=1$.

\begin{detail}[The exponent chain of Section~\Nref{sec:pointwise-weak}]\label{d9:chain42}\Ntarget{sec:pointwise-weak}
The chain has three steps.
\begin{enumerate}[label=(\roman*)]
\item $L^2$ data give $\nabla u\in L^4$, by Theorem~\Nref{thm:lp-local}(c) at $p=2$, $p^*=4$. They give values in $L^4$ by Lemma~\Nref{lem:local-sobolev}. They do not give bounded values: $p=2=\frac d2$ is the endpoint of (I2).
\item On unit cubes $L^4\subset L^{8/3}$. At $p=\frac83$, $p^*=8$ and $\beta_0=2-\frac32=\frac12$, so Proposition~\Nref{prop:scalar-lp}(b) gives $C^{1/2}$ values and $\nabla u\in L^8$.
\item At $p=8>d$, Corollary~\Nref{cor:lp-dirichlet}(a) gives $C^{1,\beta}$ with $\beta\le1-\frac48=\frac12$.
\end{enumerate}
So three interior factors take block-$L^2$ data to $C^{1,\beta}$, $\beta=\min(\alpha,\frac12)$. Route~1 (Section~\Nref{sec:pointwise-weak}, paragraph ``Route 1: boundary regularity'') is a sketch the note does not use; it is not expanded.
\status{Expanded.}
\end{detail}

\begin{detail}[Lemma~\Nref{lem:lp-kernels}, Theorem~\Nref{thm:lp-local}, Lemma~\Nref{lem:rescaled-lp}]\label{d9:lp-kernels}\Ntarget{lem:lp-kernels,thm:lp-local,lem:rescaled-lp}
\begin{step}[Lemma~\Nref{lem:lp-kernels}]
\begin{itemize}
\item (a), (b): H\"older with $(p,p')$ and Lemma~\Nref{lem:lattice-sums}(a), (c) with the exponent $sp'<d$, respectively $>d$. For example, $(\sum_{|y|\le R}\eta^d(|y|+\eta)^{-sp'})^{1/p'}\le CR^{d/p'-s}$.
\item (c): $H=\sum h(x')\one_{W(x')}$ has $\|H\|_{L^p}=\|h\|_p$, and $k_s(x-x')\le C|w-z|^{-s}$ on $W(x)\times W(x')$. So $\|k_s\star h\|_q\le C\|I_s|H|\|_{L^q}$. The Hardy--Littlewood--Sobolev theorem (Stein, \emph{Singular integrals}, Ch.~V, \S1.2) gives $\|I_sF\|_q\le C\|F\|_p$ for $0<s<d$ and $\frac1q=\frac1p-\frac{d-s}d$.
\item (d): $\eta^dK(x)=K_1(x/\eta)$. The multiplier has modulus $\le1$. The third differences of $g_1$ (Proposition~\Nref{prop:green-eta}) give H\"ormander's condition for $K_1$ and $K_1(-\cdot)$. Then the Calder\'on--Zygmund theorem on spaces of homogeneous type (Stein, \emph{Harmonic analysis}, Ch.~I), Marcinkiewicz, and duality.
\end{itemize}
\end{step}
\begin{step}[Theorem~\Nref{thm:lp-local}]
Import the representation \Neqref{eq:rep} of the proof of Theorem~\Nref{thm:local-regularity}: $u=\cG\Phi_1+\cG\Phi_2+\sum(\partial^*_\mu g_\eta)\star(\chi g_\mu)+\sum(\partial^*_\mu g_\eta)\star m_\mu$ on $\Lam_{1+2\eta}$, with $\Phi_1=\chi F-\sum(\partial^*_\mu\chi)S^{-1}_\mu g_\mu$. The terms with $\Phi_2$, $m_\mu$ live at distance $\ge\frac14$ and are bounded by $\|u\|_{\Lam_2}$ (Step~2). Step~0 handles $\eta>\frac1{16}$. The exponents:
\begin{itemize}
\item (a), $p>d$: $d-1<d/p'$. The near part ($R\sim\delta$) and the far part ($s=d$) of the H\"older quotient give $\delta^{1-d/p}\le\delta^\beta$.
\item (b), $\frac d2<p<d<q$: $d-2<d/p'$ and $d-1<d/q'$ give the sup bounds. The quotients are $\delta^{2-d/p}$ (far part with $s=d-1>d/p'$) and $\delta^{1-d/q}$.
\item (c): HLS with $s=d-1$ gives $\frac1q=\frac1p-\frac1d$, and (d) handles the $g$-term.
\end{itemize}
In every case $\|\Phi_1\|_p\le\|F\|_p+C\|g\|_p$, with $\|g\|_p$ bounded by the norm used (bounded volume).
\end{step}
\begin{step}[Lemma~\Nref{lem:rescaled-lp}]
Here $x\in\eta\Z^d$, so that $x'\mapsto x+sx'$ maps the lattice $(\eta/s)\Z^d$ onto lattice points. Put $u_s(x')=u(x+sx')$. Then $-\Lap u_s=s^2F_s+\sum\partial^*(sg_s)$, $\|h_s\|_{p,\Lam_r(0)}=s^{-d/p}\|h\|_{p,\Lam_{rs}}$ and $\partial u_s=s(\partial u)_s$. In (c) the left side scales as $s^{1-d/p^*}$, like both data terms, since $2-d/p=1-d/p^*$.
\end{step}
\status{Expanded; textbook theorems stated, not proved; the steps of Theorem~\Nref{thm:local-regularity} imported.}
\end{detail}

\begin{detail}[Proposition~\Nref{prop:scalar-lp}]\label{d9:lp-scalar}\Ntarget{prop:scalar-lp}
\emph{(a) is a literal repetition} of the first bound of Proposition~\Nref{prop:scalar-interior}, with these changes:
\begin{center}\small
\begin{tabularx}{\textwidth}{@{}>{\raggedright\arraybackslash}p{0.4\textwidth}>{\raggedright\arraybackslash}X@{}}
\toprule
in Proposition~\Nref{prop:scalar-interior} & here\\
\midrule
Lemma~\Nref{lem:rescaled}, first inequality & Lemma~\Nref{lem:rescaled-lp}, first inequality\\
$s^2|bv|$ & $s^{2-d/p}\|bv\|_{p,\Lam_{2s}}\le C\eps s^2|v|_{\Lam_{t'}}$\\
source $s^2|F|/s$ & $s^{1-d/p}\|F\|_p\le\|F\|_p$ ($s\le1$, $p>d$)\\
exponent $\alpha$ & $\beta\le\alpha$, using $[a_\mu]_\beta\le C|\partial a_\mu|\le C\eps$\\
\bottomrule
\end{tabularx}
\end{center}
The iteration (Lemma~\Nref{lem:iteration}) and the case $s<3\eta$ are unchanged.

\emph{(b).} With $A(t)=|v|_{\Lam_t}+\|\partial v\|_{p^*,\Lam_t}$ and $s=(t'-t)/6\ge3\eta$:
\begin{itemize}
\item \emph{Values.} The second inequality of Lemma~\Nref{lem:rescaled-lp} with $q=p^*$ gives the factor $s^{\beta_0}$ on the data. Also $\|bv\|_p+\|av\|_{p^*}\le C\eps|v|_{\Lam_{t'}}$ on cubes of side $\le6s\le1$.
\item \emph{Gradients.} The centres $x_i$ are the points of $\Lam_{t+s}$ on a grid of lattice points of spacing $h=\eta\lfloor s/\eta\rfloor\in[\frac12s,s]$, so that Lemma~\Nref{lem:rescaled-lp} applies at them. The cubes $\Lam_s(x_i)$ cover $\Lam_t$, and the cubes $\Lam_{3s}(x_i)$ lie in $\Lam_{t'}$ and overlap at most $(6s/h+1)^d\le13^d$ times. The third inequality, raised to the power $p^*$ and summed, gives the data terms. Here $\sum a_i^{p^*}\le(\sum a_i^p)^{p^*/p}$ and $\sum a_i^{p^*}\le(\sum a_i^2)^{p^*/2}$, since $p^*>p$ and $p^*>2$. The $L^2$ term is $s^{-\theta}\|v\|$, $\theta=1+\frac d2-\frac d{p^*}$.
\item \emph{Small $s$.} If $s<3\eta$, then $A(t)\le C\eta^{-1-d/2}\|v\|\le C(t'-t)^{-1-d/2}\|v\|$.
\end{itemize}
So $A(t)\le C\eps A(t')+C(t'-t)^{-1-d/2}\|v\|+C\mathfrak m$, using $s^{-d/2},s^{-\theta}\le s^{-1-d/2}$. Lemma~\Nref{lem:iteration} applies for $C\eps\le\frac12$, and the H\"older quotient follows at the fixed scale $s_1=(r-\rho)/8$.

\emph{Proposition~\Nref{prop:schauder} and Remark~\Nref{rem:schauder-averaging}} are not used outside Section~\Nref{sec:interior-lp}. The torus identity \Neqref{eq:C-identity} replaces them. They are not expanded.
\status{Literal repetition for (a), with the changes listed; Expanded for (b).}
\end{detail}

\begin{detail}[Proposition~\Nref{prop:system-lp}]\label{d9:lp-system}\Ntarget{prop:system-lp}
\emph{(a)} repeats the proof of Proposition~\Nref{prop:system-interior}, with $\beta$ for $\alpha$ and Lemma~\Nref{lem:rescaled-lp} for Lemma~\Nref{lem:rescaled}. Step~1 holds with $\beta$, since $\beta\le\alpha$ and the H\"older quotients of $a_\mu,b,T,c$ and $W$ are $\le C\eps$. The pointwise perturbations on $\Lam_{2s}$ have $L^p$ norm $\le Cs^{d/p}|\cdot|$, so they keep their powers of $s$. The sources enter as $s^{1-d/p}\|F\|_p\le\|F\|_p$. Steps~3--5 are unchanged.

\emph{(b), Step~2 rewritten.} Let $A_v=|v|_{\Lam_t}+\|\partial v\|_{p^*,\Lam_t}$, $v\in\{u,\sigma,\tau\}$. For the terms of Step~1 of Proposition~\Nref{prop:system-interior}:
\begin{center}\small
\begin{tabularx}{\textwidth}{@{}>{\raggedright\arraybackslash}p{0.32\textwidth}>{\raggedright\arraybackslash}X>{\raggedright\arraybackslash}p{0.2\textwidth}@{}}
\toprule
type of term & bound on a cube of side $\sim s$ & factor (values; gradients)\\
\midrule
pointwise in $F_\bullet$: $bu$, $T^T(\mathfrak F\sigma)$, $\kappa mu$, $T^Tu$, $b\sigma$, $b\tau$ & $\|f\|_{p}\le Cs^{d/p}|f|$; the $\ell^{p^*}$-sum over $Cs^{-d}$ cubes gives $s^{d/p-d/p^*}$ & $s^{\beta_0+d/p}=s^2$; $s$\\
gradient in $F_\bullet$: $D\tau$, $\eps\partial u$, $\eps\partial\sigma$ & H\"older: $\|\partial v\|_p\le Cs\|\partial v\|_{p^*}$ & $s^{1+\beta_0}$; $s$\\
in $g_\bullet$: $a_\mu v$, $\mathfrak F\sigma$, $W(u-D\sigma)$ & $\le C\eps(|v|\,s^{d/p^*}$ or $\|\partial v\|_{p^*})$ & $C\eps$\\
$u_\mu$ in $g_\sigma$ (coefficient $1$) & $\|u\|_{p^*,\Lam_{3s}}\le Cs^{d/p^*}|u|$ & $s^{\beta_0+d/p^*}=s$; $C$\\
\bottomrule
\end{tabularx}
\end{center}
With these factors, the three bounds of Step~2 of Proposition~\Nref{prop:system-interior} take the same form:
\begin{itemize}
\item $\sigma$: the $u_\mu$ term gives $C_4A_u$, and the rest $C\eps$;
\item $\tau$: every coefficient carries $\eps$ (from $\mathfrak F$, $W$, $T$, $b$, $a$);
\item $u$: $D\tau$ gives $C_3sA_\tau$, $\kappa mu$ gives $C_3\kappa sA_u$, and the rest $C\eps$.
\end{itemize}
The remainder is $R=C(\sum\mathfrak m^{p,p^*}_r(\mathcal S_\bullet)+(t'-t)^{-\theta'}\mathfrak m)$ with $\theta'=1+\frac d2-\frac d{p^*}$. The H\"older quotient of $\partial\sigma$ (needed in (a) for $W D\sigma$ in $g_\tau$) does not enter, because only $\|\partial\sigma\|_{p^*}$ appears. Step~3 (the weights $w_\sigma,w_\tau$ and the small cubes $s\le s_*$), Lemma~\Nref{lem:iteration}, and Step~5 (covering, $\ell^{p^*}$-sums) apply without change. The H\"older quotients of the values follow at a fixed scale from the second inequality of Lemma~\Nref{lem:rescaled-lp}.
\status{Literal repetition, with the changes checked.}
\end{detail}

\begin{detail}[Corollary~\Nref{cor:lp-dirichlet}]\label{d9:lp-cor}\Ntarget{cor:lp-dirichlet}
For $O=\bG^X$, follow the basic estimate of the proof of Theorem~\Nref{thm:ID}. Here $u=\bG^Xf$, and $(u,\sigma,\tau)$ satisfies \Neqref{eq:system} on $\Lam_3(c_y)$ with
\[
\mathcal S_u=f+D\pi_1-D\pi_2-a\bQ^T\bQ u,\qquad\mathcal S_\sigma=-\pi_1,\qquad\mathcal S_\tau=-aQ'^TQ'\tau .
\]
\begin{itemize}
\item $f=F+D^T_Xg$ with $D^T_Xg=\partial^*g+T^Tg$ in the local gauge. $T^Tg$ goes into $F$, with $\|T^Tg\|_p\le C\eps\|g\|_{p^*}$ on bounded cubes.
\item $D\pi_i$, $a\bQ^T\bQ u$ and $aQ'^TQ'\tau$ are bounded pointwise by $C\mathfrak L(u)$ (Proposition~\Nref{prop:scalar-interior} for $\pi_i$, and the sup bounds of $\bQ^T$, $Q'^T$).
\item $\mathfrak L(u)\le C(\mathfrak E^p_y(F)+\mathfrak E^q_y(g))$ for $p,q\ge2$. This uses the $L^2$ decay of $\bG^X$, $D\bG^X$, $\bG^XD^T_X$, $D\bG^XD^T_X$ (Theorem~\Nref{thm:reg-gauge}(a)) and $\|h\|_{B}\le\|h\|_{p,B}$ for unit volume; $\mathfrak E^q_y(g)\le\mathfrak E^\beta_y(g)$.
\end{itemize}
Proposition~\Nref{prop:system-lp}(a) or (b) concludes. For $O=G'_X$, $D^TDv=f-aQ'^TQ'v$ in the interior with $|Q'v(y')|\le C\|v\|_{B(y')}$, and Proposition~\Nref{prop:scalar-lp} applies.

For the $L^q\to L^q$ claim: $p=dq/(d+q)$ maps $q\in(d,\infty)$ onto $p\in(\frac d2,d)$, with $p^*=q$. The endpoint $q=d$ gives $p=\frac d2$, which is excluded, as the note says.
\status{Literal repetition, with the changes checked.}
\end{detail}

% IX.8
\subsection{Tiers, interior factors, trace and cut term}\label{d9:tiers}

\begin{detail}[Tiers and the inequalities between them, Definition~\Nref{def:tier}]\label{d9:tier-def}\Ntarget{def:tier}
Blocks have volume $\eta^d\cdot\eta^{-d}=1$, so H\"older's inequality on $B(y)$ gives, with constant $1$,
\[
\mathfrak n^0_y\le\mathfrak n^{1/2}_y\le\mathfrak n^1_y\le|\cdot|_{B(y)}\le\mathfrak n^2_y,\qquad\mathfrak n^0_y\le\mathfrak n^3_y .
\]
Hence in a global tier norm the far part (in $\mathfrak n^0$) can be absorbed in the near part without a constant, and Lemma~\Nref{lem:tier-calculus}(b) holds with constant $1$.
\begin{itemize}
\item $\mathfrak n^2_y(\nabla h)\le C\mathfrak n^3_y(h)$, because $\widetilde{\nabla h}=\nabla\tilde h$ in the local gauge (covariance of $\nabla$).
\item $\mathfrak n^2_y(g)\le C\max_{|y'-y|_\infty\le1}\mathfrak n^3_{y'}(g)$. The H\"older quotient of $\tilde g_y$ on $\Lam_1(c_y)$ is bounded by $|\nabla\tilde g|$ there, which is covered by the blocks within $1$ of $y$, up to the gauge correction $C\eps|g|$.
\item $K_j$ and $\operatorname{supp}\nabla h_j$ lie at block distance $\ge(\rho-1)M-1$ from $\partial\Lam^v_j$ (Detail~\ref{d9:parametrix}, Step~1).
\item The inequalities $d_\Sigma(y)\le|y-w|_\infty\le|y-w|$ give the seam rules with Euclidean weights and $\infty$-distance seams.
\end{itemize}
\status{Expanded.}
\end{detail}

\begin{detail}[Lemma~\Nref{lem:tier-calculus} and tier norms of families]\label{d9:tier-calc}\Ntarget{lem:tier-calculus}
\begin{enumerate}[label=(\alph*)]
\item $\mathfrak n^0_y(Fg)=\|\one_yFg\|$, and $\|\one_yFg\|\le\sum_{y'}\|\one_yF\one_{y'}\|\mathfrak n^0_{y'}(g)$, with kernel $k=\|\one_yF\one_{y'}\|$ and $k'=0$.
\item Insert the bound of $G$ into that of $F$. With the absorption above, the kernels compose as matrices, and $e^{\gamma|y-y''|}\le e^{\gamma|y-y'|}e^{\gamma|y'-y''|}$ gives submultiplicativity of the weighted row sums.
\item For $y\in Y_F$:
\[
\mathfrak n^c_y(FGg)\le\sum_{|y'-y|\le R_F}k_F\,\mathfrak n^b_{y'}(Gg)+\sum_{y'}k'_F\,\mathfrak n^0_{y'}(Gg) .
\]
In the first sum $y'\in Y_G$, so $\mathfrak n^b_{y'}(Gg)$ splits into a near part of radius $R_G$ (total radius $R_F+R_G$) and a far part. In the second, $\mathfrak n^0_{y'}(Gg)\le\sum_{y''}\|\one_{y'}G\one_{y''}\|\mathfrak n^0_{y''}(g)$. The weighted row sum is at most
\[
\Big(\sup_y\sum_{y'}e^{\gamma|y-y'|}k_F\Big)\,c_G+\Big(\sup_y\sum_{y'}e^{\gamma|y-y'|}k'_F\Big)\,|G|_\gamma\le c_F(c_G+|G|_\gamma) .
\]
Only row sums of $G$ enter.
\item Let $\varphi=h_j-h_j(c_y)$, which depends on the output block $y$ only through a constant. Then $[F,h_j]g=F(\varphi g)-\varphi Fg$. By Lemma~\Nref{lem:partition}, $|\partial^k\varphi|\le CM^{-k}$ and $|\varphi|\le CM^{-1}(1+|y-y'|)$ on $\Lam_1(c_{y'})$. For each tier:
\begin{itemize}
\item $t\le1$: immediate.
\item $t=2$: $\widetilde{\varphi g}=\varphi\tilde g$, and
\[
[\varphi\tilde g]_\beta\le|\varphi|_{\Lam_1}[\tilde g]_\beta+[\varphi]_\beta|\tilde g|_{\Lam_1},\qquad|\tilde g|_{\Lam_1}\le|g|_{B(y')}+C[\tilde g]_\beta .
\]
\item $t=3$: the discrete Leibniz rule gives $\nabla_\mu(\varphi h)=(S_\mu\varphi)\nabla_\mu h+(\partial_\mu\varphi)h$. The H\"older quotient needs $|\tilde h|$ and $|\nabla\tilde h|$ on all of $\Lam_1(c_{y'})$:
\begin{itemize}
\item $|\nabla\tilde h|_{\Lam_1}\le|\nabla h|_{B(y')}+C[\nabla\tilde h]_\beta$;
\item $|\tilde h|_{\Lam_1}\le|h|_{B(y')}+C|\partial\tilde h|_{\Lam_1}$, with $\partial\tilde h=\nabla\tilde h-T\tilde h$ and $|T|\le C\eps$;
\item the term $C\eps|\tilde h|_{\Lam_1}$ is absorbed for $C\eps\le\frac12$.
\end{itemize}
\end{itemize}
In all cases $\mathfrak n^t_{y'}(\varphi g)\le CM^{-1}(1+|y-y'|)\mathfrak n^t_{y'}(g)$ and $\mathfrak n^t_y(\varphi Fg)\le CM^{-1}\mathfrak n^t_y(Fg)$. Then $(1+t)e^{-\gamma t/2}\le C\gamma^{-1}$ gives the rate $\gamma/2$.
\item The seam rule \Neqref{eq:seam-product} holds for tier kernels by the same inequality. In the local composition (c), $d_\Sigma(y)\le|y-y'|+d_\Sigma(y')$ moves the seam weight from the seam factor to its neighbours, in the near and in the far parts, at the cost of the factor $2$ in their rates.
\item \emph{Families} (Section~\Nref{ss:44.1}, paragraph ``Weakenable families''). The proofs of Lemma~\Nref{lem:weak-basic}(a) and Lemma~\Nref{lem:weak-chain} use three facts: subadditivity of the seminorms; submultiplicativity, which is (b); and the supports. The H\"older quotient at $y$ is taken on $\Lam_1(c_y)$, which can reach one more block, so the margins are increased by $1$. Tier norms are weighted row sums; the column sums of Lemma~\Nref{lem:weak-basic}(a) are needed only for the decay norm $|\cdot|_\gamma$, which is kept separately in $|\cdot|^{\rm sp}$.
\end{enumerate}
\status{Expanded.}
\end{detail}

\begin{detail}[Lemma~\Nref{lem:split-chain}]\label{d9:split}\Ntarget{lem:split-chain}
Let $c_i=|F_i|^{a_i\to b_i}_{\gamma;R_i,Y_i}$. We prove by downward induction on $k$ that
\[
|F_k\cdots F_n|^{a_m\to b_k}\le2c_k\prod_{k<i\le m}(c_i+|F_i|_\gamma)\prod_{i>m}|F_i|_\gamma .
\]
\begin{itemize}
\item \emph{$k=m<n$.} Detail~\ref{d9:tier-calc}(c) with $G=F_{m+1}\cdots F_n$, which satisfies $0\to0$ at every block with constant $|G|_\gamma$, gives the bound $c_m(|G|_\gamma+|G|_\gamma)$. Since $a_m=0$, the near part is of tier $0$ and becomes a far part.
\item \emph{$k=m=n$.} Trivial.
\item \emph{Inductive step $k<m$.} Detail~\ref{d9:tier-calc}(c) gives at most $c_k(c_G+|G|_\gamma)$, with $c_G$ the bound for $k+1$. Using $|F_i|_\gamma\le c_i+|F_i|_\gamma$,
\[
c_G+|G|_\gamma\le(2c_{k+1}+|F_{k+1}|_\gamma)\prod_{k+1<i\le m}(c_i+|F_i|)\prod_{i>m}|F_i|\le2\prod_{k<i\le m}(c_i+|F_i|)\prod_{i>m}|F_i| .
\]
\end{itemize}
Adding $|F_1\cdots F_n|_\gamma\le\prod|F_i|_\gamma$ gives $|F_1\cdots F_n|^{\rm sp}\le3\prod|F_i|^{\rm sp}$. The summation is that of Lemma~\Nref{lem:weak-chain}, since its proof uses only the product form of the one-term bound. The inclusion hypothesis is checked term by term. A $j$-term of $\sum_js^{X(j)}\varphi_jG_j$ vanishing outside $K_j^{(2)}$ satisfies its bound on $Y_j\cup\{y\notin K_j^{(2)}\}$.
\status{Expanded.}
\end{detail}

\begin{detail}[Lemma~\Nref{lem:local-sobolev}]\label{d9:sobolev}\Ntarget{lem:local-sobolev}
$w=\chi v$ is finitely supported on $\eta\Z^4$, and $d=4\ge3$, so $w=g_\eta\star(-\Lap w)=\sum_\mu(\partial^*_\mu g_\eta)\star(\partial_\mu w)$ by summation by parts. Proposition~\Nref{prop:green-eta} gives $|\partial^*g_\eta|\le Ck_{d-1}$. Lemma~\Nref{lem:lp-kernels}(c) with $s=3$, $p=2$, $\frac1q=\frac12-\frac14$ gives $\|w\|_4\le C\|\partial w\|_2$. Finally $\partial(\chi v)=(S\chi)\partial v+(\partial\chi)v$.
\status{Expanded.}
\end{detail}

\begin{detail}[Lemma~\Nref{lem:int-factors}]\label{d9:int-factors}\Ntarget{lem:int-factors}
\begin{step}[Geometry]
For $y\in Y_j$ all blocks within $8$ of $y$ lie in $(Z_j)_{(4)}$. So $y\in(Z_j)_{(4)}$, and $\Lam_3(c_y)$ lies in $Z_j$ at distance $>\frac12\ge8\eta$ from the complement (for $\eta\le\frac1{16}$; coarser lattices are Step~0 of Theorem~\Nref{thm:local-regularity}). So the equations of Lemma~\Nref{lem:local-system} (gauge case) or $D^TDv=f-aQ'^TQ'v$ (scalar case) hold on $\Lam_3(c_y)$. Corollary~\Nref{cor:lp-dirichlet} applies at $y$.
\end{step}
\begin{step}[The split form from the proof of Corollary~\Nref{cor:lp-dirichlet}]
In that proof the bound at $y$ is the bound of Proposition~\Nref{prop:system-lp} (resp.\ \Nref{prop:scalar-lp}) on $\Lam_3(c_y)\to\Lam_1(c_y)$. It is of the form $C\big(\mathfrak n_3(\mathcal S)+\mathfrak m\big)$. The source is $\mathcal S=f+(\text{terms in }\pi_i,\ \bQ^T\bQ u,\ Q'^TQ'\tau)$.
\begin{itemize}
\item \emph{The data $f$} enter only through their norms on $\Lam_3(c_y)$. These are covered by the blocks within $4$ of $y$:
\begin{itemize}
\item $\|f\|_{p,\Lam_3}\le\big(\sum_{|y'-y|\le4}\|f\|^p_{p,B(y')}\big)^{1/p}\le\sum_{|y'-y|\le4}\mathfrak n^a_{y'}(f)$;
\item for H\"older data, two points of $\Lam_3$ at distance $\le\frac12$ lie in a common $\Lam_1(c_{y'})$, larger distances use the sup norm, and the gauge conversion costs $C\eps|f|$ (proof of Theorem~\Nref{thm:ID}(b)).
\end{itemize}
This is the near part, with radius $4$.
\item \emph{The other terms} and $\mathfrak m$ are bounded by $C\mathfrak L(u)$, $\mathfrak L(u)=\sum_{y'}e^{-\gamma|y-y'|}(\|u\|_{B(y')}+\|Du\|_{B(y')})$ (Lemma~\Nref{lem:local-system}). By the $L^2$ decay of $G_j$ and $DG_j$ (or with $D^T$ on the right), $\mathfrak L(u)\le\sum_{y''}k'(y,y'')\mathfrak n^0_{y''}(f)$ with $\sum_{y''}e^{\gamma_1|y-y''|}k'\le CC_GS_{\gamma_0-\gamma_1}$. This is the far part, at rate $\gamma_1=\gamma_0/2$.
\end{itemize}
\end{step}
\begin{step}[The entries]
\begin{center}
\begin{tabularx}{\textwidth}{@{}>{\raggedright\arraybackslash}p{0.34\textwidth}>{\raggedright\arraybackslash}X@{}}
\toprule
entry & source\\
\midrule
$\varphi G_j\colon1\to3$ & Corollary~\Nref{cor:lp-dirichlet}(a), $p=8>d$, $\beta\le1-\frac48=\frac12$, data $F=g$\\
$\varphi G_jD^T\colon1\to2$ (values) & Corollary~\Nref{cor:lp-dirichlet}(b), $p=\frac83$, $p^*=8$, $f=D^T_Xg$, $\beta\le2-\frac32=\frac12$\\
$\varphi DG_jD^T\colon1\to1$ & the same, the bound $\|\nabla u\|_{8,B(y)}$ (the Calder\'on--Zygmund input)\\
$\varphi G_jD^T\colon2\to3$ & Corollary~\Nref{cor:lp-dirichlet}(a) with $F=0$ and H\"older $g$ ($\mathfrak E^\beta_y$), i.e.\ Proposition~\Nref{prop:system-interior}\\
$\varphi DG_j\colon1\to2$, $\varphi DG_jD^T\colon2\to2$ & from $1\to3$ and $2\to3$, by $\mathfrak n^2(\nabla h)\le C\mathfrak n^3(h)$\\
$\varphi\bG^{Z_j}D\colon2\to3$ (scalar $\lambda$, supported in $Z_j$ at distance $\ge2\eta$ from the complement) & see below\\
\bottomrule
\end{tabularx}
\end{center}
For the last entry: the restriction of $D\lambda$ to $\ell^2_1(Z_j)$ is $D\one_{Z_j}\lambda$, with no cut term. In the local gauge, $D\lambda=\partial\lambda+T\lambda$ and $\partial_\mu\lambda=-\partial^*_\mu(S_\mu\lambda)$. So $D\lambda$ is a source in local form with H\"older $g=-S\lambda$ (tier $2$ of $\lambda$) and bounded $F=T\lambda$. The far part uses the decay of $\bG^{Z_j}D\one_{Z_j}=(D^T_{Z_j}\bG^{Z_j})^T$. Multiplication by $\varphi$ preserves all tiers, by the Leibniz rule of Detail~\ref{d9:tier-calc}(d).
\end{step}
\begin{step}[$L^2$ inputs, and the full gradient]
The decay of $G_j$, $DG_j$, $G_jD^T$, $DG_jD^T$ is Theorems~\Nref{thm:reg-scalar}(b) and~\Nref{thm:reg-gauge}(a). The decay of $D^T_{Z_j}\bG^{Z_j}$ and $D^T_{Z_j}\bG^{Z_j}D^T$ holds because $\tn\cdot_X$ contains $\|D^T_XA\|$; it is in the list of Theorem~\Nref{thm:reg-gauge}(a). Where the full gradient of $v$ is needed in $L^2$ near $y$, take a cutoff $\chi$ supported inside $Z_j$ (so $D^T=D^T_{Z_j}$ on $\chi v$). Lemma~\Nref{lem:hodge} in the local gauge gives
\[
\|\partial(\chi v)\|^2=\|d(\chi v)\|^2+\|d^*(\chi v)\|^2,
\]
and $D=d+T$, $|T|\le C\eps$. So $\|\partial(\chi v)\|\le\|D(\chi v)\|+\|D^T_{Z_j}(\chi v)\|+C\eps\|\chi v\|$, and the commutators with $\chi$ are bounded.
\end{step}
\begin{step}[$\varphi\equiv1$, and the torus]
With $\varphi\equiv1$ the proof uses only $\Lam_3(c_y)$, so the bounds hold at every $y\in Y_j$. For $X=\Tf$ (Section~\Nref{ss:44.2}): a block within $2$ of $K_j$, together with the blocks within $8$ of it, is at distance $\ge(\rho-1)M-1-2-8>4$ from $\Tf\setminus\Lam^v_j$. So $K_j^{(2)}\subset Y_j$, and the bounds of Lemmas~\Nref{lem:int-factors}, \Nref{lem:endpoint-factors} hold at every block where the factors do not vanish. They are global tier bounds.
\end{step}
\status{Expanded.}
\end{detail}

\begin{detail}[Lemma~\Nref{lem:endpoint-factors}]\label{d9:endpoint}\Ntarget{lem:endpoint-factors}
The proof in the note is written out. We checked:
\begin{itemize}
\item the covering of $\Lam_3$ by the $5^d$ cubes $\Lam_1(x_i)$, $x_i\in\Lam_2$ on a grid of spacing $1$, with $\Lam_2(x_i)\subset\Lam_4$;
\item $\|aQ'^TQ'v\|_{p,\Lam_r}\le C\|v\|_{\Lam_{r+1}}$ for $p\le4$, from $|Q'w(y')|\le C\|w\|_{B(y')}$, the boundedness of $Q'^T$ in sup norm, and block-locality;
\item $|\tau_\mu|\le2e^2\eps$ in $D=\partial+T$;
\item H\"older on $\Lam_3$ (volume $6^d$): $\|f\|_{8/3}\le C\|f\|_4$;
\item the averaging commutator $[aQ'^TQ',h_j]v=aQ'^TQ'(\psi v)-\psi aQ'^TQ'v$ on $B(y)$, with $|\psi|\le CM^{-1}$ on $B(y)$.
\end{itemize}
The entries from tier $0$ use only $\eps\le1$, and the entries from tier $\frac12$ use $\eps\le\eps'_S$.
\status{Expanded.}
\end{detail}

\begin{detail}[Scalar gradients at the cut, and rectangles, Section~\Nref{ss:44.2}]\label{d9:cut-par}\Ntarget{ss:44.2}
For a scalar $\lambda$ on $\Tf$ and a domain $Z$, the restriction of $D\lambda$ to $\ell^2_1(Z)$ is $D(\one_Z\lambda)+c_Z(\lambda)$. On a bond $b$ crossing $\partial Z$, with exterior endpoint $x\notin Z$, $c_Z(\lambda)(b)=\pm\eta^{-1}(\text{transport of }\lambda(x))$. So $|c_Z(\lambda)|\le\eta^{-1}|\lambda(x)|$, and $\|c_Z(\lambda)\|$ is not bounded uniformly in $\eta$.

$\Lam^v_j$ is a product of intervals of length at most $2\rho M+2$. Its complement in each direction has length $\ge n-2\rho M-2\ge2$, since $n\ge8(\rho+C_v)M$. So $\Lam^v_j$ is a rectangle in the sense of Lemma~\Nref{lem:trace} at every $\eta\le1$. The note argues via $X(j)$ (Section~\Nref{sec:tiers}, \emph{Scalar gradients at the cut}): $(\Lam^v_j)^{(R_v)}\subset X(j)$ and $X(j)$ does not wrap, so each complement has at least $2R_v\ge2$ points. This argument does not depend on the choice of the hat.
\status{Expanded.}
\end{detail}

\begin{detail}[Lemma~\Nref{lem:trace}]\label{d9:trace}\Ntarget{lem:trace}
\begin{step}[The identity at $U=1$]
Let $A$ vanish off $\mathfrak B(Z)$. Lemma~\Nref{lem:hodge} on the torus gives $\|dA\|^2+\|d^*A\|^2=\|\partial A\|^2$. Let $x\notin Z$ have a bond into $Z$.
\begin{itemize}
\item \emph{At most one neighbour in $Z$.} If $x+\eta e_\mu\in Z$, then $x_\lambda\in I_\lambda$ for all $\lambda\ne\mu$. A second neighbour $x\pm\eta e_\nu\in Z$ with $\nu\ne\mu$ would give $x_\mu\in I_\mu$, hence $x\in Z$. And $x-\eta e_\mu\in Z$ would need the complement of $I_\mu$ to be the single point $x_\mu$.
\item So $(d^*A)(x)=\pm\eta^{-1}A(b_x)$, with $b_x$ the unique bond from $x$ into $Z$.
\item The parallel bond $b'_x$ continuing outward has its far endpoint at distance $2\eta$ from $I_\mu$ in direction $\mu$. It is outside $Z$ because the complement has $\ge2$ points, so $A(b'_x)=0$.
\item The pair $(b_x,b'_x)$ contributes $\eta^d|\eta^{-1}(A(b_x)-0)|^2=\eta^{d-2}|A(b_x)|^2=\eta^d|(d^*A)(x)|^2$ to $\|\partial A\|^2$, and distinct $x$ give distinct pairs.
\end{itemize}
Hence $\|dA\|^2+\|d^*_ZA\|^2=\|\partial A\|^2-\|\one_{Z^c}d^*A\|^2\ge\|\partial A\|^2_{\mathfrak B(Z)}$, since all omitted terms are non-negative.
\end{step}
\begin{step}[One-dimensional trace]
Let $b\in\mathfrak B_{\rm cr}(y')$ have direction $\pm e_\nu$. Since $Z$ is a union of blocks with side $\ge1$, the bonds $b_m=b\mp m\eta e_\nu$, $0\le m<\eta^{-1}$, lie in $\mathfrak B(Z)$. Then $|A(b)|\le|A(b_m)|+\sum_{i<m}|A(b_i)-A(b_{i+1})|$. Averaging over $m$ and summing over the crossing bonds of $B(y')$ with weight $\eta^{d-1}$:
\[
\sum_b\eta^{d-1}|A(b)|\le\sum\eta^d|A(b_m)|+\sum\eta^d|\partial_\nu A_\nu|\le C\big(\|A\|+\|\partial A\|_{\mathfrak B(Z)}\big)_{B(y')^{(1)}}
\]
by Cauchy--Schwarz on a region of volume $\le C$.
\end{step}
\begin{step}[Localization and background]
For $(\chi A)(b)=\chi(b_-)A(b)$ we have $\chi A\in\ell^2_1(Z)$. The commutators of $d$, $d^*_Z$ with $\chi$ are differences of $\chi$ times shifted values, bounded by $C\|A\|_{B(y')^{(2)}}$. The quantities are gauge-invariant. In a gauge with (B) on a cube $\supset B(y')^{(3)}$ (regularity cubes have half-side $4M+2\ge6$), $D=d+T$ with $|T|\le C\eps$. At complex perturbations, $D_{U_\delta}=D_U+T_\delta$ with $|T_\delta|\le C\delta_c$.
\end{step}
\begin{step}[The slab example]
For a slab of one fine layer, with $A_1=a$ on the normal bonds through both faces: $dA=0$ and $d^*_ZA=0$, the localized norm is $O(\sqrt\eta)|a|$, and the trace is $\sim|a|$. So the side length $\ge1$ is needed.
\end{step}
\status{Expanded.}
\end{detail}

\begin{detail}[Lemma~\Nref{lem:cut-term}, and the $\mathrm b$-symmetry of $\bG^Z$]\label{d9:cut-term}\Ntarget{lem:cut-term}
\begin{step}[The $\mathrm b$-symmetry]
By Section~\Nref{ss:19.1},
\[
\fA^{Z,\kappa}=\Pi_Z^T\mathfrak B_Z\Pi_Z+DR_ZD^T_Z+a\Pi_Z^T\bQ^T\bQ\Pi_Z+D\one_{\Tf\setminus Z}D^T+\kappa m_Z .
\]
Here $\mathfrak B_Z^T=\mathfrak B_Z$ (since $W^T=W$), and $R_Z^T=R_Z$. The latter holds because $P_{\cV,Z}=G'_ZQ'^T(M'_Z)^{-1}Q'G'_Z$ with $(G'_Z)^T=G'_Z$ (from $(\fA'_Z)^T=\fA'_Z$) and $(M'_Z)^T=M'_Z$. Also $(D^T)^T=D$ for the bilinear transpose. So $(\fA^{Z,\kappa})^T=\fA^{Z,\kappa}$ as an algebraic identity, at real and at complex references (Section~\Nref{ss:19.1}, paragraph ``Dirichlet conditions for all components''), and therefore $(\bG^Z)^T=\bG^Z$. The identity theorem invoked in the proof of Lemma~\Nref{lem:cut-term} is not needed. At real references the bilinear transpose is the adjoint on real fields.
\end{step}
\begin{step}[$L^2$ input]
Theorem~\Nref{thm:reg-gauge}(a) gives decay $(C,\gamma_0)$ for $P\bG^XQ$, $P\in\{1,D,D^T_X\}$ and $Q\in\{1,D^T,D\one_X\}$. For $Q=D\one_X$ use $\ip u{D\one_Xg}=\ip{D^T_Xu}g$.
\end{step}
\begin{step}[The bound]
$\|u\|_{B(y)}=\sup|\mathrm b(w^*,u)|$ over $w$ supported in $B(y)$ with $\|w\|\le1$. By Step~1,
\[
\mathrm b(w^*,\bG^Zc_Z(\lambda))=\mathrm b(\bG^Zw^*,c_Z(\lambda)),\qquad|\mathrm b(v,c_Z(\lambda))|\le C\sum_{b\ \rm crossing}\eta^{d-1}|v(b)|\,|\lambda(x_b)| .
\]
Lemma~\Nref{lem:trace} applies to $v=\bG^Zw^*$, $\bG^ZD^Tw^*$, $\bG^ZD\one_Zw^*$ (for $D\bG^Z$, $D^T_Z\bG^Z$). It needs $\tn v_Z$ near the crossing blocks $y'$, which is bounded by $Ce^{-\gamma_0|y-y'|}$ by Step~2. Summing over $y'$ gives the stated bound with $|\lambda|_{B(y')^{(1)}}$.
\end{step}
\begin{step}[The consequence]
Write $\lambda=\lambda_{\rm in}+\lambda_{\rm cut}$, with $\lambda_{\rm in}=0$ within $2$ of $\Tf\setminus Z$ and $\lambda_{\rm cut}$ supported within $3$. For $\lambda_{\rm in}$, use the entry $\varphi\bG D\colon2\to3$ of Detail~\ref{d9:int-factors}. At $y\in Y_j$, at distance $>12$ from the cut, the source $D(\one_Z\lambda_{\rm cut})+c_Z(\lambda)$ vanishes on $\Lam_3(c_y)$. So it enters only through $\mathfrak L(u)$:
\begin{itemize}
\item through the decay of $\{1,D,D^T_Z\}\bG^ZD\one_Z$, giving tier $0$ in $\lambda_{\rm cut}$;
\item through Step~3, giving $|\lambda|$, i.e.\ tier $2$.
\end{itemize}
So the far part is in $\mathfrak n^2$ at the blocks meeting $\partial Z$.
\end{step}
\status{Expanded. The ``taken as stated'' $\mathrm b$-symmetry is closed algebraically.}
\end{detail}

% IX.9
\subsection{Lift, non-local parts and the tier walks}\label{d9:liftsec}

In this subsection $\gamma_0$ is the rate of level~(L3). By the paragraph \emph{Rates} (Section~\Nref{ss:44.4}, paragraph ``Rates'', Detail~\ref{d9:rates}), the walks of (L1) and (L2) have output and compatibility rates at least $32\gamma_0$. Their tier factors (Lemmas~\Nref{lem:int-factors}, \Nref{lem:endpoint-factors}) therefore hold at the rate $16\gamma_0\ge2\gamma_0$. $C_T\ge1$ and $C_G\ge1$ are the constants of these lemmas and of (W2). $Z_{(r)}$ and $Y_j$ are as in Lemma~\Nref{lem:int-factors}.

\begin{detail}[The rates, Section~\Nref{ss:44.4}]\label{d9:rates}\Ntarget{ss:44.4}
At (L1) and (L3) let $\gamma_0$ be at most three numbers:
\begin{itemize}
\item the rate of (W2);
\item the rate $\gamma$ of Corollary~\Nref{cor:lp-dirichlet};
\item at (L3), $\frac1{32}$ of the smallest output and compatibility rate of (L1), (L2).
\end{itemize}
These output and compatibility rates are $\gamma_0^{(1)}/8$ and $\gamma_0^{(1)}/16$ by Detail~\ref{d9:wi-b}, or $\gamma_0^{(1)}/32$ as stated in the note. Then:
\begin{itemize}
\item Lemmas~\Nref{lem:int-factors} and~\Nref{lem:endpoint-factors} hold at $\gamma_0/2$: their far kernels decay at least at $\gamma_0$, and $\sum_{y'}e^{(\gamma_0/2-\gamma_0)|y-y'|}=S_{\gamma_0/2}$.
\item The (L1), (L2) inputs are available at $32\gamma_0$, hence at $2\gamma_0$ in decay and at $\gamma_0$ in seam form, as (W3) at (L3) requires.
\end{itemize}
Each level divides the rate by a fixed number ($8$ for outputs, $32$ for compatibility), and there are six levels. So all final rates are positive and depend only on the rates of Parts~III, VII and VIII.
\status{Expanded.}
\end{detail}

\begin{detail}[Lemma~\Nref{lem:lift}]\label{d9:lift}\Ntarget{lem:lift}
\emph{The idea.} At the scalar levels the walk has only local steps. So each term is a product of the first factor $\sum_js^{X(j)}h_jG'_j$ and of step factors built from the $G'_j$ and smooth cut-offs of size $M^{-1}$ or $M^{-2}$ (Step~1). Near a block of the support of $h_j$, each factor is a local propagator evaluated in the interior of its own domain (Step~2), where Lemmas~\Nref{lem:int-factors} and~\Nref{lem:endpoint-factors} give tier bounds (Step~3). This gives (Z1)--(Z3) (Steps~4--6), and the telescoping identity of Detail~\ref{d9:wi-b} gives the seam form of the differences (Step~7).
Let $Z$ be a domain, $Z_j=\Lam^v_j\cap Z$ and $G'_j=G'_{Z_j}$.
\begin{step}[The terms of $G'^Z(s)$]
At (L1), $\mathfrak N=0$, and the walk has only local steps. A term with $n$ steps is $\Phi_0\Phi_1\cdots\Phi_n$, with $\Phi_0=\mathcal F_0=\sum_js^{X(j)}h_jG'_j$ and $\Phi_i=\mathcal B$. By Detail~\ref{d9:L1}(iii),
\[
\mathcal B=\sum_js^{X(j)}\big(\varphi'_jG'_j+\varphi_jDG'_j+[aQ'^TQ',h_j]G'_j\big)\quad\text{(up to shifts by at most $\eta$)},
\]
with $\varphi'_j,\varphi_j$ supported in $K_j^{(1)}$, $|\partial^k\varphi'_j|\le CM^{-2}$ and $|\partial^k\varphi_j|\le CM^{-1}$ for $k\le2$. This uses $|\partial^kh_j|\le C_kM^{-k}$ up to $k=4$ (Lemma~\Nref{lem:partition}(c)). Also $D\mathcal F_0=\sum_js^{X(j)}(h_jDG'_j+[D,h_j]G'_j)$, where $[D,h_j]$ is a multiplication by $\psi_j$ with $|\partial^k\psi_j|\le CM^{-1-k}$, up to a shift.
\end{step}
\begin{step}[Interior blocks]
Let $y\in Z_{(16)}$ be within distance $2$ of $K_j$. A block within distance $8$ of $y$:
\begin{itemize}
\item is at distance $>16-8>4$ from $\Tf\setminus Z$;
\item lies within $10$ of $K_j$, hence at distance $\ge(\rho-1)M-1-10\ge5>4$ from $\Tf\setminus\Lam^v_j$, using $(\rho-1)M\ge16$.
\end{itemize}
Since $\Tf\setminus Z_j=(\Tf\setminus Z)\cup(\Tf\setminus\Lam^v_j)$, such blocks lie in $(Z_j)_{(4)}$, so $y\in Y_j$. At blocks farther than $2$ from $K_j$ the $j$-terms of $\mathcal F_0$, $D\mathcal F_0$ and $\mathcal B$ vanish: $h_j$, $\psi_j$, $\varphi_j$, $\varphi'_j$ vanish there, and $[aQ'^TQ',h_j]$ vanishes on blocks where $h_j=0$, because $Q'^TQ'$ is block-local.
\end{step}
\begin{step}[The three factor families]
By Step~2 and Lemmas~\Nref{lem:int-factors}, \Nref{lem:endpoint-factors}, applied with $\varphi=h_j$, $M\varphi_j$ and $M^2\varphi'_j$, every term satisfies, on $Z_{(16)}$, with radius $4+1=5$ (the shifts cost $1$):
\begin{center}
\begin{tabular}{@{}lll@{}}
\toprule
family & tier bound & source\\
\midrule
$\mathcal F_0$ & $\frac12\to2$ (values) & $\varphi G'_j$ of Lemma~\Nref{lem:endpoint-factors}\\
$D\mathcal F_0$ & $\frac12\to1$ & $\varphi DG'_j$ of Lemma~\Nref{lem:endpoint-factors}, and $\psi_jG'_j$: $\frac12\to2\subset\frac12\to1$\\
$\mathcal B$ & $0\to\frac12$ & entries from tier $0$ of Lemma~\Nref{lem:endpoint-factors}\\
\bottomrule
\end{tabular}
\end{center}
The averaging commutator maps $0\to1$ at every block (Lemma~\Nref{lem:endpoint-factors}), hence $0\to\frac12$, since $\mathfrak n^{1/2}\le\mathfrak n^1$. In the anchored sums over $j$ the monomial costs $e^{(\kappa+\delta)\mathsf N_{\rm b}}$ and $\#\{j:\Box\subset X(j)\}\le C_j$. With the decay norms of Lemma~\Nref{lem:step-small} at (L1),
\[
\|\mathcal F_0\|^{\rm sp}_{\kappa+\delta,2\gamma_0}+\|D\mathcal F_0\|^{\rm sp}_{\kappa+\delta,2\gamma_0}\le C_je^{(\kappa+\delta)\mathsf N_{\rm b}}(C_T+C_G),\qquad\|\mathcal B\|^{\rm sp}_{\kappa+\delta,2\gamma_0}\le C_je^{(\kappa+\delta)\mathsf N_{\rm b}}C_TCM^{-1}+\theta .
\]
\end{step}
\begin{step}[(Z1)]
The term with $n=0$ is $\mathcal F_0$, respectively $D\mathcal F_0$. It maps $\frac12\to2$ (values), respectively $\frac12\to1$, and a fortiori $1\to2$, respectively $1\to1$.

For $n\ge1$ apply Lemma~\Nref{lem:split-chain} with $m=2$:
\begin{itemize}
\item $F_1=\Phi_0$ (or $D\Phi_0$), $a_1=\frac12$, on $Y_1=Z_{(24)}$, radius $5$;
\item $F_2=\Phi_1$, $a_2=0$, $b_2=\frac12$, on $Y_2=Z_{(16)}$;
\item $F_3,\dots,F_{n+1}=\Phi_2,\dots,\Phi_n$, in the decay norm.
\end{itemize}
The hypotheses of the lemma hold:
\begin{itemize}
\item $a_1=\frac12\le b_2=\frac12$;
\item $a_2=0$;
\item every block within $5$ of $Z_{(24)}$ lies in $Z_{(16)}$, since $24-5>16$.
\end{itemize}
The result is $0\to2$ (values), respectively $0\to1$. Its near part is of tier~$0$, so it is a far part.

Summing over $n$ with the bound of Lemma~\Nref{lem:split-chain}, and with $C_d\delta^{-2}\theta\le\frac14$ (Theorem~\Nref{thm:weak-inverse}(a) at (L1) at the rate $\kappa+\delta$), gives
\[
\|G'^Z(s)\|^{\rm sp}_{\kappa,2\gamma_0}\le\|\mathcal F_0\|^{\rm sp}+3\sum_{n\ge1}(n+1)(C_d\delta^{-2})^2\|\mathcal F_0\|^{\rm sp}\|\mathcal B\|^{\rm sp}4^{-(n-1)}\le C(\kappa,\delta) .
\]
So the near part enters only through the term without steps, with radius $5\le8$.
\end{step}
\begin{step}[(Z2)]
$P^Z_\cV(s)=G'^Z(s)\,F_2$, with $F_2=Q'^T(M'^Z)^{-1}(s)Q'G'^Z(s)$.
\begin{itemize}
\item $|(Q'G'^Z(s)g)(y)|\le C\sum_{y'}\|\one_yQ'G'^Z(s)\one_{y'}\|\,\mathfrak n^0_{y'}(g)$, since $|Q'\lambda(y)|\le C\|\lambda\|_{B(y)}$.
\item $(M'^Z)^{-1}(s)$ is a unit-lattice family with a decay norm, from (L2).
\item $Q'^T$ is block-local, and $|Q'^T\omega|_{B(y)}\le C|\omega(y)|$ (Proposition~\Nref{prop:linearized-averages}(e)). Since blocks have volume $1$, $\mathfrak n^1_y\le|\cdot|_{B(y)}$.
\end{itemize}
So $F_2$ maps $0\to1$ at every block, with only a far part, and its split norm is its decay norm. Lemma~\Nref{lem:split-chain} with $m=2$ applies: $F_1=G'^Z(s)$ is $1\to2$ on $Z_{(24)}$ by (Z1), $F_2$ is $0\to1$ with radius $0$, and $a_1=1\le b_2=1$, $a_2=0$. It gives $0\to2$ (values). The same holds for $DP^Z_\cV(s)$, with $0\to1$.
\end{step}
\begin{step}[(Z3)]
Inside products $R_Z(s)=\one_Z-P^Z_\cV(s)$, with the literal identity. At a block $y\in Z_{(24)}$, $\one_Zg=g$ on $\Lam_1(c_y)$, so the identity preserves every tier with radius $0$. $P^Z_\cV(s)$ maps $0\to2\subset0\to t$ for $t\le2$, using $\mathfrak n^t\le\mathfrak n^2$. Standing alone, $R_Z(s)$ has its identity split box by box (Section~\Nref{ss:40.2}). That split does not preserve the H\"older tiers, and no tier bound is claimed for it. Every use of $R(s)$ in a tier bound is inside a product: Detail~\ref{d9:N-tier} and Theorem~\Nref{thm:HC-pointwise}. So no indicator split has to preserve a tier.
\end{step}
\begin{step}[Seams]
Let $X'$ be a domain, $P=\one_{X'}$ and $Q=\one_{\Tf\setminus X'}$. Write the difference of matched terms on $\Tf$ and on $X'$ by the telescoping identity of Detail~\ref{d9:wi-b}, Step~1:
\begin{itemize}
\item the factors to the left of the distinguished factor are walk factors on $X'$; Steps~2--4 apply to them with $Z=X'$ on $X'_{(24)}$ and $X'_{(16)}$;
\item the factors to the right are walk factors on $\Tf$, measured in the decay norm.
\end{itemize}
Let $y\in X'_{(24)}$. We go through the kinds of distinguished factor.
\begin{enumerate}[label=(\roman*)]
\item \emph{A difference of local propagators at position $k\le1$.} It is $\varphi(G'_{Z_1}-G'_{Z_2})$, $\varphi D(G'_{Z_1}-G'_{Z_2})$, or the averaging commutator of the difference, with $Z_1=\Lam^v_{j'}$ and $Z_2=\Lam^v_{j'}\cap X'$, between factors $P$. The argument $f$ is supported in $X'$, so $\one_{Z_1}f=\one_{Z_2}f$. Put $\delta v=(G'_{Z_1}-G'_{Z_2})f$.

At points $x$ of $Z_2$ at least one block away from $\Tf\setminus Z_2$, both Dirichlet operators act on $\delta v$ as the torus operator $D^TD+aQ'^TQ'$. Here ``one block'' is needed because $Q'^TQ'$ is block-local. Since $\fA'_{Z_i}G'_{Z_i}f=f$ there,
\[
D^TD\,\delta v=-aQ'^TQ'\,\delta v\qquad\text{at such points.}
\]
Let $y''$ be a block within distance $5$ of $y$ where $\varphi\ne0$. Then $y''$ lies within $1$ of $K_{j'}$ and at distance $>19$ from $\Tf\setminus X'$. So $\Lam_4(c_{y''})$ consists of such points.

The first bound of Proposition~\Nref{prop:scalar-interior}, with the bounded source $-aQ'^TQ'\delta v$ and $|Q'^TQ'h|_{B}\le C\|h\|_B$, gives
\[
\mathfrak n^3_{y''}(\delta v)\le C\sum_{|y'''-y''|_\infty\le4}\|\delta v\|_{B(y''')}.
\]
By Theorem~\Nref{thm:reg-scalar}(e), these norms are of seam form for $\Sigma(Z_1,Z_2)$, the unit points within $2$ of $Z_1\setminus Z_2\subset\Tf\setminus X'$. That set lies in $\Sigma_w(\Tf,X')$.

So this factor maps $0\to3\subset0\to2$ in seam form, with no near part. The same holds for $\varphi D\delta v$, giving tier~$2$, and for the averaging commutator, giving tier~$1$. If the distinguished factor is $\Phi_1$, the factor $\Phi'_0$ to its left maps $\frac12\to2$, and the seam version of Lemma~\Nref{lem:split-chain} (Lemma~\Nref{lem:tier-calculus}(e)) composes, with $\Phi'_0$ at $2\gamma_0$ and the distinguished factor at $\gamma_0$.
\item \emph{A difference at position $k\ge2$.} It is of seam form in $L^2$, as in Detail~\ref{d9:wi-b}, Step~2. The factors $\Phi'_1$ ($0\to\frac12$) and $\Phi'_0$ ($\frac12\to2$, respectively $\frac12\to1$) act as in (Z1). The seam version of Lemma~\Nref{lem:split-chain} moves the seam weight onto the neighbours at twice the rate.
\item \emph{A cross term $P\Phi_kQ$ or a factor whose local domain $\Lam^v_{j'}\cap X'$ is empty.} Such a factor is supported, on its input or its output side, within $r_0+1\le3$ of $\Tf\setminus X'$, where $r_0\le2$ is the range of the local part at (L1). On the seam, the seam norm at $\gamma_0$ is bounded by the norm at $2\gamma_0$.
\begin{itemize}
\item For $k\ge2$ it is treated as in (ii).
\item For $k\le1$, the blocks within distance $5$ of $y$ are at distance $>24-5-3=16$ from this support. So the near part of $\Phi'_0$ at $y$ does not see it.
\item The bound $0\to\frac12$ of $\Phi_1$ involves only $\mathfrak n^0$ (Lemma~\Nref{lem:endpoint-factors}: in the entries from tier $0$ both sums involve only $\mathfrak n^0$).
\end{itemize}
So only $L^2$ kernels see the support, and the seam weight follows from \Neqref{eq:seam-product}.
\end{enumerate}
For (Z2) and (Z3), telescope the product $G'(s)\,Q'^T\,(M')^{-1}(s)\,Q'G'(s)$. The distinguished factor is of one of two kinds:
\begin{itemize}
\item a difference of the left $G'(s)$, bounded by the seam part of (Z1) with an input of tier~$1$ at every block;
\item a difference of $Q'G'(s)$ or of $(M')^{-1}(s)$, of seam form in $L^2$ by Theorem~\Nref{thm:weak-inverse}(b) at (L1), (L2) and Detail~\ref{d9:prod-compat}. It is followed by $Q'^T$ (block-local, giving tier~$1$) and by $G'^{X'}(s)$ ($1\to2$ on $X'_{(24)}$).
\end{itemize}
\end{step}
\begin{step}[Constants]
They are products of $C_T$, $C_G$, the walk constants of (L1), (L2) (Theorem~\Nref{thm:weak-inverse} at $M\ge M_1$, where the value at $M_1$ bounds all larger $M$), $e^{(\kappa+\delta)\mathsf N_{\rm b}}$, $C_j$ and powers of $\delta^{-1}$. None depends on $M$. The factors $M^{-1}$ in $\|\mathcal B\|^{\rm sp}$ only help.
\end{step}
\status{Expanded. The seam bullet (i) needs the interior points of $Z_2$ to be at least one block from the cut, because $Q'^TQ'$ is block-local; this holds with room to spare.}
\end{detail}

\begin{detail}[Lemma~\Nref{lem:N-tier}]\label{d9:N-tier}\Ntarget{lem:N-tier}
\begin{step}[The terms]
Detail~\ref{d9:L3}, Step~1, gives the four terms; here we take $W=0$. In each middle factor the derivatives stand next to the scalar families, as $DR=D-DP_\cV$, $DG'$, $G'D^T$ and $DG'D^T$.
\begin{center}
\begin{tabular}{@{}llll@{}}
\toprule
term & $X^T$ & $M$ & $Y$\\
\midrule
(i) $-DP_\cV D^T$ & $1$ & $DP_\cV$ & $D^T_Z$\\
(ii) $-D^T\mathfrak FS$ & $D^T$ & $\mathfrak FG'R$ & $D^T_Z$\\
(iii) $-DRG'\mathfrak F^TD$ & $1$ & $DRG'\mathfrak F^T$ & $D$\\
(iv) $DRG'\mathfrak F^T\mathfrak FG'RD^T$ & $1$ & $DRG'\mathfrak F^T\mathfrak FG'R$ & $D^T_Z$\\
\bottomrule
\end{tabular}
\end{center}
The terms containing $W$, namely $-WDS$, $-DRG'D^TW$ and the $W$-part of (iv), vanish. We used $(\mathfrak BDS)^T=S^T(\mathfrak F^TD+D^TW)$, $S^T=DRG'$, and $D^T_ZD^T=\one_Z\mathfrak F^T$ on $2$-forms.
\end{step}
\begin{step}[$\mathfrak F$ and $\mathfrak F^T$]
These are multiplications by the coefficients $c_{\mu\nu}$ of Lemma~\Nref{lem:curvature} (or their transposes), composed with shifts by at most $2\eta$. In the local gauge, $|c_{\mu\nu}|\le C\eps$. Since $\beta\le\alpha$ and $\Lam_1(c_y)$ has side $2\le6$, $[c_{\mu\nu}]_{\beta,\Lam_1(c_y)}\le C[c_{\mu\nu}]_{\alpha,\Lam_3(c_y)}\le C\eps$. Hence:
\begin{itemize}
\item $\mathfrak F$ and $\mathfrak F^T$ map $t\to t$ for $t\in\{0,1\}$, with radius $1$, no far part, and constant $C\eps$;
\item $\mathfrak F$ maps $2\to2$, by the Leibniz rule for the H\"older quotient;
\item a shift by at most $4\eta$ stays within adjacent blocks, so $\mathfrak F^T\mathfrak F$ also has radius $1$.
\end{itemize}
\end{step}
\begin{step}[The tier bounds]
Let $y\in Z_{(48)}$. We use the bounds of Detail~\ref{d9:lift}:
\begin{itemize}
\item (i): $DP_\cV$ maps $0\to1$ by (Z2).
\item (ii): $\mathfrak FG'R=\mathfrak FG'-\mathfrak FG'P_\cV$. $G'$ maps $1\to2$ (values) with radius $8$; $G'P_\cV$ maps $0\to2$ (by (Z2) and (Z1) with tier~$2\subset$ tier~$1$ input); $\mathfrak F$ maps $2\to2$ with radius $1$. So $M$ maps $1\to2$ with radius $9$.
\item (iii): $DRG'\mathfrak F^T=DG'\mathfrak F^T-DP_\cV G'\mathfrak F^T$. $\mathfrak F^T$ maps $1\to1$ (radius $1$) and $DG'$ maps $1\to1$ (radius $8$). $DP_\cV$ maps $0\to1$, and then $G'\mathfrak F^T$ is measured in the decay norm. So $M$ maps $1\to1$ with radius $9$.
\item (iv): expand both factors $R$. The term $DG'\mathfrak F^T\mathfrak FG'$ maps $1\to1$ with radius $8+1+8=17$: $G'$ maps $1\to2\subset1$, $\mathfrak F^T\mathfrak F$ maps $1\to1$, $DG'$ maps $1\to1$. Each of the three terms with $P_\cV$ either has $DP_\cV$ ($0\to1$) on the left, followed by decay norms, or has $G'P_\cV$ ($0\to2$) on the right, followed on the left by $\mathfrak F^T\mathfrak F$ and $DG'$. These terms map $0\to1$.
\end{itemize}
The radii add up to at most $17$. In each product, the factor at position $p$ is applied on $Z_{(48-r_p)}$, where $r_p\le17$ is the sum of the radii of the factors to its left. Since $Z_{(48-17)}=Z_{(31)}\subset Z_{(24)}$, Lemma~\Nref{lem:lift} applies to it, and the inclusions of Lemma~\Nref{lem:split-chain} hold. Local factors are composed term by term with Lemma~\Nref{lem:tier-calculus}(c) and Lemma~\Nref{lem:weak-local-ops}(a). The number of products is bounded. So $C'_{\mathfrak N}(\kappa,\delta)$ is a product of constants of Lemma~\Nref{lem:lift} and $C\eps\le C$, and is free of $M$.
\end{step}
\begin{step}[(N3), seams]
Telescope the products defining $M_i^{\Tf}$ and $M_i^{X'}$ (Detail~\ref{d9:wi-b}, Step~1). The distinguished factor is one of three kinds:
\begin{itemize}
\item a difference of one of the families of Lemma~\Nref{lem:lift}, bounded by its seam part on $X'_{(24)}$;
\item a cross term, supported on the seam;
\item a difference of the bounded local operators $\mathfrak F$, $\mathfrak F^T$, $\one_Z$, supported within distance $1$ of $\Tf\setminus X'$. A derivative is never such a factor: every $D$ and $D^T$ in a middle factor stands next to a scalar family of Lemma~\Nref{lem:lift}, and is compared with it.
\end{itemize}
The other factors are at rate $2\gamma_0$, and the seam version of Lemma~\Nref{lem:split-chain} composes. The radii are as in Step~3. So the near blocks lie in $X'_{(31)}\subset X'_{(24)}$.
\end{step}
\status{Expanded.}
\end{detail}

\begin{detail}[Proposition~\Nref{prop:nl-tiers}]\label{d9:nl-tiers}\Ntarget{prop:nl-tiers}
\emph{The idea.} The non-local parts enter the tier walk through the same terms as in the $L^2$ walk: commutators with $h_j$ and a seam term. Each is a product of tier-bounded factors (Steps~1 and~2). The commutators gain $(\gamma_0M)^{-1}$ from $\nabla h_j$ (Steps~3 and~4), and the seam term gains $e^{-\gamma_0(\rho-1)M/16}$ because it lives far from the support of $h_j$ (Step~5). Step~6 treats an external derivative, Step~7 sums, and Step~8 fixes the thresholds, with doubled constants at complex references.
Fix $j$, let $X'=\Lam^v_j$, $G_j=\bG^{X'}$, and let $F$ be a term of $M_i=M_i^{\Tf}(s)$ or of the difference in the fourth part. By the regrouping of Detail~\ref{d9:step-small}, the non-local parts of step $j$ are:
\[
X_i^T[M_i,h_j]Y_iG_j,\qquad[X_i^T,h_j]M_iY_iG_j,\qquad X_i^TM_i[Y_i,h_j]G_j,\qquad h_jX_i^T\big(PM_iP-M_i^{X'}\big)Y_iG_j .
\]
Here $Y_i\in\{D^T_{X'},D\}$ (since $W=0$). Write $b=2$ for the term~(ii) and $b=1$ otherwise. For (ii), $h_jD^T=D^Th_j-[D^T,h_j]$. So $\mathcal A^{\rm nl}$ consists of the first, third and fourth parts of (ii), without the left $D^T$, and must map $1\to2$. $\mathcal B^{\rm nl}$ consists of the rest and must map $1\to1$.
\begin{step}[The input]
Let $g$ have tier~$1$ and $u=Y_iG_jg$.
\begin{itemize}
\item[(a)] For $y'\in(\Lam^v_j)_{(16)}$: $(\Lam^v_j)_{(16)}\subset Y_j$, because the blocks within $8$ of $y'$ are at distance $>8$ from the complement. Lemma~\Nref{lem:int-factors} with $\varphi\equiv1$ gives $G_j\colon1\to3$, i.e.\ Corollary~\Nref{cor:lp-dirichlet}(a) at $p=2d=8$, with $\beta\le\frac12$. The field $u$ is a combination of components of $\nabla G_jg$, up to shifts, and $\mathfrak n^2(\nabla h)\le C\mathfrak n^3(h)$. So $\mathfrak n^3_{y'}(G_jg)+\mathfrak n^2_{y'}(u)\le\sum_{y''}k_1(y',y'')\mathfrak n^1_{y''}(g)$. Here the near and far parts are merged, using $\mathfrak n^0\le\mathfrak n^1$.
\item[(b)] At every block, $\mathfrak n^0_{y'}(u)\le\sum_{y''}k_0(y',y'')\mathfrak n^0_{y''}(g)$, with $k_0(y',y'')=\|\one_{y'}Y_iG_j\one_{y''}\|$ (W2).
\end{itemize}
The weighted row sums of $k_0$ and $k_1$ at $\gamma_0/2$ are at most $C\max(C_T,C_G)$. With $D^T$ on the right, (a) holds with $\mathfrak n^2_{y'}(G_jD^Tg)+\mathfrak n^1_{y'}(u)$ on the left. This comes from the entries $1\to2$ (values) of $\varphi G_jD^T$ and $1\to1$ of $\varphi DG_jD^T$, the latter for every first-order derivative.
\end{step}
\begin{step}[Inclusions]
$M\ge(2R_N+22)/(\rho-1)$ gives $(\rho-1)M-1\ge2R_N+21$. So:
\begin{itemize}
\item every block within $2R_N+2$ of $K_j$ is at distance $\ge19$ from $\Tf\setminus\Lam^v_j$, hence lies in $(\Lam^v_j)_{(16)}$;
\item every block within $2$ of $K_j$ is at distance $\ge2R_N+19=67$, hence lies in $(\Lam^v_j)_{(48)}$.
\end{itemize}
\end{step}
\begin{step}[The commutator with $M_i$]
The bound $1\to b$ of $F$ (Lemma~\Nref{lem:N-tier}, $Z=\Tf$, so it holds at every block) has kernels $k,k'$ at $2\gamma_0$ and radius $R_N=24$.
\begin{itemize}
\item \emph{$y$ at distance $>R_N+2$ from $K_j$.} Then $h_j$ vanishes on $\Lam_1(c_y)$ and on the blocks within $R_N$ of $y$, so $[F,h_j]u=F(h_ju)$ on $\Lam_1(c_y)$. The near part sees $h_ju=0$, and in the far part $|h_j(x)|=|h_j(x)-h_j(c_y)|\le CM^{-1}|x-c_y|$. So
\[
\mathfrak n^b_y([F,h_j]u)\le CM^{-1}\sum_{y'}(1+|y-y'|)k'(y,y')\mathfrak n^0_{y'}(u) .
\]
\item \emph{$y$ within $R_N+2$ of $K_j$.} Write $[F,h_j]u=F(\psi u)-\psi Fu$ with $\psi=h_j-h_j(c_y)$. Multiplication by $\psi$ changes $\mathfrak n^t_{y'}$ by at most $CM^{-1}(1+|y-y'|)$ (Detail~\ref{d9:tier-calc}, (d)). The near part involves $\mathfrak n^1_{y'}(u)\le\mathfrak n^2_{y'}(u)$ at blocks within $R_N$ of $y$, hence within $2R_N+2$ of $K_j$. These lie in $(\Lam^v_j)_{(16)}$ by Step~2, where Step~1(a) applies. The far part uses Step~1(b).
\end{itemize}
With $(1+t)e^{-\gamma_0t}\le C\gamma_0^{-1}$ (passing from $2\gamma_0$ to $\gamma_0$, and from $\gamma_0$ to $\gamma_0/2$ in the composition with $k_0,k_1$), $[F,h_j]Y_iG_j$ satisfies a global bound $1\to b$ at $\gamma_0/2$ with constant
\[
C(\gamma_0M)^{-1}\big(|F|^{1\to b}_{2\gamma_0;R_N,\Tf}+|F|_{2\gamma_0}\big)\max(C_T,C_G) .
\]
\end{step}
\begin{step}[The commutators with $X_i^T$ and $Y_i$]
$[D^T,h_j]$ and $[Y_i,h_j]$ have size $CM^{-1}$ and are supported in $K_j^{(1)}$ (up to shifts by $\eta$).
\begin{itemize}
\item $[Y_i,h_j]G_jg$ is supported in $K_j^{(1)}$. There $G_jg$ has tier~$3$, by Steps~1(a) and~2, so this field has tier~$1$ everywhere, with kernel $CM^{-1}k_1$. Since $F$ satisfies its bound at every block, $X_i^TF[Y_i,h_j]G_j$ maps $1\to b$ with constant $CM^{-1}(|F|^{1\to b}+|F|)C_T$.
\item For (ii), $[D^T,h_j]Fu$ vanishes at blocks farther than $2$ from $K_j$. At the other blocks $y$, $\mathfrak n^1_y([D^T,h_j]Fu)\le CM^{-1}\sum_{|y'-y|_\infty\le1}\mathfrak n^2_{y'}(Fu)$. The near part of $F$ at $y'$ involves $u$ within $R_N+3\le2R_N+2$ of $K_j$, where Step~1(a) applies. So this part maps $1\to1$.
\end{itemize}
\end{step}
\begin{step}[The seam term]
Let $\Delta$ be a term of $PM_iP-M_i^{X'}$. By Lemma~\Nref{lem:N-tier}(N3), $\Delta$ has a seam bound on $X'_{(48)}$ at $\gamma_0$ for $\Sigma=\Sigma_w(\Tf,X')$, with kernels $k_\Delta,k'_\Delta$ and constant $c_\Delta$. Only blocks $y$ within $2$ of $K_j$ contribute, because of $h_j$, or $[D^T,h_j]$ for (ii). These lie in $X'_{(48)}$ by Step~2, and
\[
d_\Sigma(y)\ge(\rho-1)M-1-2-w\ge\tfrac14(\rho-1)M-3
\]
by \Neqref{eq:rho-choice}. So $\sum_{y'}e^{\gamma_0|y-y'|}(k_\Delta+k'_\Delta)(y,y')\le e^{-\gamma_0d_\Sigma(y)}c_\Delta\le e^{3\gamma_0}e^{-\gamma_0(\rho-1)M/4}c_\Delta$. The near blocks lie within $R_N+2$ of $K_j$ (Step~1(a)); the far part uses Step~1(b). So the seam parts satisfy a global bound $1\to b$ at $\gamma_0/2$, with constant
\[
Ce^{-\gamma_0(\rho-1)M/16}\big(c_\Delta+|\Delta|_{\gamma_0,\Sigma}\big)\max(C_T,C_G) .
\]
\end{step}
\begin{step}[An external $D$]
Let $\lambda$ be a scalar of tier~$2$ and $u=Y_iG_jD\lambda$, with $G_jD\lambda=\bG^{X'}(D(\one_{X'}\lambda)+c_{X'}(\lambda))$ (Section~\Nref{ss:44.2}, paragraph ``Scalar gradients at the cut''). Steps~3--5 are said to ``apply verbatim''. The only changes are in the input, Step~1:
\begin{itemize}
\item On $(\Lam^v_j)_{(16)}$, $G_jD\lambda$ has tier~$3$. For $\lambda$ away from the cut this is the entry $2\to3$ of $\varphi\bG^{Z_j}D$ (Lemma~\Nref{lem:int-factors}). For $\lambda$ near the cut it is Lemma~\Nref{lem:cut-term}, whose far part is in $\mathfrak n^2$.
\item At every block, $\mathfrak n^0_{y'}(u)\le C\sum_{y''}e^{-\gamma_1|y'-y''|}\mathfrak n^2_{y''}(\lambda)$. This follows from the decay of $Y_i\bG^{X'}D\one_{X'}$, from Lemma~\Nref{lem:cut-term}, and from $|\lambda|_{B(y'')^{(1)}}\le\max_{|y'''-y''|_\infty\le1}\mathfrak n^2_{y'''}(\lambda)$.
\end{itemize}
So $\mathfrak n^1_{y''}(g)$ and $\mathfrak n^0_{y''}(g)$ are replaced by $\mathfrak n^2_{y''}(\lambda)$. Steps~3--5 then use only these two input bounds, and they give $2\to2$ for $\mathcal A^{\rm nl}D$ and $2\to1$ for $\mathcal B^{\rm nl}D$, with the same factors. Only a constant is claimed.
\end{step}
\begin{step}[Summation]
A part is non-zero only if $X_\omega^{(r)}$ meets $K_j^{(1)}$, so at most $C\deg\omega$ values of $j$ contribute for a fixed $\omega$. Evaluating at $\kappa+\delta$ costs $C\delta^{-1}$, the monomial $s^{X(j)}$ costs $e^{\kappa\mathsf N_{\rm b}}$, and $\#\{j:\Box\subset X(j)\}\le C_j$. With Lemma~\Nref{lem:N-tier} at $\kappa+\delta$ and $\gamma_0\le1$, Steps~3--5 give $\theta_{\mathfrak N}(\kappa)$ and Step~6 gives $C^D_{\mathfrak N}(\kappa)$. The output rate $\gamma_0/2$ exceeds $\gamma=\gamma_0/8$.

The proof gives one factor $\max(C_T,C_G)$; the note writes the square. With $C_T,C_G\ge1$ the stated bound dominates (compare Detail~\ref{d9:step-small}). Since $C$, $C'_{\mathfrak N}$, $C_T$, $C_G$ are free of $M$, so are $C^D_{\mathfrak N}$ and $M\theta_{\mathfrak N}$.
\end{step}
\begin{step}[The thresholds, and doubled constants]
The conditions in $M_{\mathfrak N}$ come from three places:
\begin{itemize}
\item Lemma~\Nref{lem:N-tier} at $\kappa+\delta$ needs Lemma~\Nref{lem:lift} at $\kappa+2\delta$, hence $M\ge M_1$ of (L1), (L2) at $\kappa+3\delta$, and $(\rho-1)M\ge16$;
\item Step~2 needs $(\rho-1)M\ge2R_N+22$;
\item tier norms raise the margins by $1$, and weakenable families need margins $\le M/16$, so $M\ge16(r+1)$.
\end{itemize}
Complex references. Four facts are needed:
\begin{enumerate}[label=(\roman*)]
\item $M_1(\kappa,\delta,C_G,C_{\rm loc},C_{\mathfrak N},\gamma_0,\rho,C_v)$ increases in each multiplicative constant, because $\theta$ does (Detail~\ref{d9:wi-a}, Step~2). So an $M_1$ computed with these constants doubled is valid for every set of constants up to twice the real ones.
\item $C'_{\mathfrak N}$, $C_T$, $C_G$ and $C_{\rm loc}$ are free of $M$ at real references (Step~7 and Detail~\ref{d9:N-tier}, Step~3).
\item $\theta_{\mathfrak N}$ contains $C'_{\mathfrak N}\max(C_T,C_G)^2$, a product of three inputs. If each input merely doubled, it could grow by a factor $8$. If each input changes by at most a factor $1+\sigma_{\rm cx}$, $\sigma_{\rm cx}=2^{1/3}-1$, it changes by at most $(1+\sigma_{\rm cx})^3=2$. The same holds for $\theta'$, a sum of products of at most three inputs.
\item The rates are fixed. The coercivity constants are the fixed ones of Proposition~\Nref{prop:complex-ext}, and every decay bound that enters is taken at the fixed rate of the differences there, a quarter of the minimum of the perturbation rate and the two chosen inverse decay rates in Lemma~\Nref{lem:perturbed-inverse}. Neither depends on $\delta_c$.
\end{enumerate}
Step~(vii) of the proof of Proposition~\Nref{prop:complex-ext} fixes $C_{\rm cx}$ level by level. For $\delta_c\le1/C_{\rm cx}$, every input of every level changes by at most a factor $1+\sigma_{\rm cx}$ at a complex reference of size $\delta_c$, and every constant by at most $1+C_{\rm cx}\delta_c\le2$ (Detail~\ref{xa:d-majorant}). The constants of the walks that enter later levels as inputs are taken as their explicit bounds, valid once the walk ratio is at most $\frac12$; for the tier walk this is $4C_d\delta^{-2}C_je^{(\kappa+\delta)\mathsf N_{\rm b}}\max(C_T,C_G)$ (Theorem~\Nref{thm:tier-walk}). These are polynomials in the inputs of their level, and the convexity bound of step~(vii) is applied to them and to $\theta'$. So the note computes $M_{\mathfrak N}$ with the multiplicative constants of Part~IX in $M_1$ doubled, and $M_1'$ with $\theta'$ replaced by $2\theta'$; both serve at every complex reference with $\delta_c<\delta_P$, and neither depends on $\bar\delta$ (Proposition~\Nref{prop:nl-tiers}, Theorem~\Nref{thm:tier-walk}). Detail~\ref{xa:d-majorant} also checks that $M_{\mathfrak N}$ and $M_1'$ are not circular.
\end{step}
\status{Expanded at real references. At complex references the constants are read as in step~(vii) of the proof of Proposition~\Nref{prop:complex-ext}, expanded in Detail~\ref{xa:d-majorant}.}
\end{detail}

\begin{detail}[Remark~\Nref{rem:nl-scope}]\label{d9:nl-scope}\Ntarget{rem:nl-scope}
The results are used only with $W=0$ at (L3). Theorems~\Nref{thm:tier-walk}, \Nref{thm:HC-pointwise} and Proposition~\Nref{prop:nl-tiers} are cited only in Sections~\Nref{sec:pointwise-weak} (discussion), \Nref{sec:tiers} and~\Nref{sec:weak-background}, in Proposition~\Nref{prop:complex-ext}(b), ``with $W=0$ at level~(L3)'', and in Remark~\Nref{rem:domain-constants} (thresholds). All have $W=0$. The exponent: (I1) at $p=8$ gives $\beta\le1-\frac48=\frac12$, and (I2) at $p=\frac83$ gives $\beta\le2-\frac{4}{8/3}=\frac12$.
\status{Expanded.}
\end{detail}

\begin{detail}[Theorem~\Nref{thm:tier-walk}]\label{d9:tier-walk}\Ntarget{thm:tier-walk}
\begin{step}[The input of a step]
A step acts on $u=G_jg$ with $g$ of tier~$1$. On $(\Lam^v_j)_{(16)}$, $u$ has tier~$3$ and $Du$, $D^Tu$ have tier~$2$ (Corollary~\Nref{cor:lp-dirichlet}(a), $p=8$). Near the cut only $L^2$ norms are controlled. On $\Tf$, $K_j^{(2)}\subset Y_j$ (Detail~\ref{d9:int-factors}, Step~5), so all local factors satisfy global tier bounds.
\end{step}
\begin{step}[Local parts]
The terms are those of Detail~\ref{d9:L3}, Step~4, with $W=0$, together with the expansion of $D[D^T,h_j]G_j$ into $\varphi'G_j$ and $\varphi\nabla G_j$. All are supported in $K_j^{(1)}$. By Lemma~\Nref{lem:int-factors}:
\begin{itemize}
\item the $\mathcal A$-part $\varphi G_j$ maps $1\to3$, and $\varphi G_jD^T$ maps $1\to2$;
\item the $\mathcal B$-parts $\varphi\nabla G_j$ and $\varphi'G_j$ map $1\to2\subset1\to1$, also with $D^T$ on the right ($1\to1$);
\item the averaging commutators map $1\to1$ by the argument of Lemma~\Nref{lem:endpoint-factors} ($\bQ^T\bQ$ has range one block).
\end{itemize}
All carry the factor $CM^{-1}$ (or $CM^{-2}$). The counting of Lemma~\Nref{lem:step-small} gives the first term of $\theta'$. At (L1), $\mathcal A=0$, and these are all the parts.
\end{step}
\begin{step}[The walk]
A word is $\Phi_0\cdots\Phi_n$ (Detail~\ref{d9:wstep}). Measure the factors as follows:
\begin{itemize}
\item $\Phi_i$, $i\ge2$, in $1\to1$;
\item $\Phi_1$ in $1\to2$ if $\Phi_0=\mathcal F_0D^T$, which forces $\Phi_1$ to be of type $\mathcal A$, and in $1\to1$ otherwise;
\item $\Phi_0$ in $2\to3$ or $1\to3$.
\end{itemize}
The tiers match: the output tier of each factor is at least the input tier of its left neighbour. Lemma~\Nref{lem:weak-chain} holds in global tier norms (Detail~\ref{d9:tier-calc}, (f)). With $2C_d\delta^{-2}\theta'(\kappa+\delta)\le\frac12$,
\[
\|G(s)\|^{1\to3}_{\kappa,\gamma}\le C_d\delta^{-2}C_je^{(\kappa+\delta)\mathsf N_{\rm b}}\max(C_T,C_G)\sum_n(n+1)2^{-n}=4C_d\delta^{-2}C_je^{(\kappa+\delta)\mathsf N_{\rm b}}\max(C_T,C_G)
\]
\end{step}
\begin{step}[The other statements]
\begin{itemize}
\item \emph{External $D^T$.} It is carried by $\Phi_n$. For $n=0$, $\mathcal F_0D^T$ maps $1\to2$ and $2\to3$. For $n\ge1$, $\mathcal BD^T$ maps $1\to1$ and $\mathcal AD^T$ maps $1\to2$ (local parts: Lemma~\Nref{lem:int-factors}; non-local parts: Proposition~\Nref{prop:nl-tiers}). For $G(s)D^T\colon2\to3$ the last factor needs only $2\to1$, since the rest maps $1\to3$.
\item \emph{External $D$, part (b).} For $n=0$ this is the entry $\varphi\bG^{Z_j}D\colon2\to3$ and Lemma~\Nref{lem:cut-term}. For $n\ge1$, the local parts of $\Phi_nD$ ($\varphi\nabla G_jD$, $\varphi G_jD$, $[a\bQ^T\bQ,h_j]G_jD$) follow from Lemma~\Nref{lem:int-factors} in the interior and from Lemma~\Nref{lem:cut-term} at the cut. The non-local parts follow from Proposition~\Nref{prop:nl-tiers} with $C^D_{\mathfrak N}$. They occur once per term, so a constant suffices.
\item \emph{Left $D$.} It uses $D\mathcal F_0\colon1\to2$ and $D\mathcal F_0D^T\colon2\to2$.
\end{itemize}
The pointwise bound is (a) together with $\mathfrak n^1_{y'}(f)\le|f|_{B(y')}$.
\end{step}
\begin{step}[$M_1'$]
Both terms of $\theta'(\kappa+\delta)$ are at most a constant, free of $M$, times $M^{-1}$: the first explicitly, and $\theta_{\mathfrak N}$ by Proposition~\Nref{prop:nl-tiers}, since $e^{-\gamma_0(\rho-1)M/16}\le C(\gamma_0,\rho)M^{-1}$. So $M_1'$ exists. It is taken with $2\theta'$. At complex references of size $\delta_c<\delta_P$ the inputs of $\theta'$ change by at most a factor $1+\sigma_{\rm cx}$, $\sigma_{\rm cx}=2^{1/3}-1$, and $\theta'$ is a sum of products of at most three of them, so $2\theta'$ bounds $\theta'$ there (step~(vii) of the proof of Proposition~\Nref{prop:complex-ext}; Detail~\ref{xa:d-majorant}). The rates are fixed as in Detail~\ref{d9:nl-tiers}, Step~8, item~(iv). So $M_1'$ does not depend on $\bar\delta$.
\end{step}
\status{Expanded.}
\end{detail}

\begin{detail}[The identity \Neqref{eq:C-identity} and $\bH(s)$, $\bC(s)$, Section~\Nref{ss:44.5}]\label{d9:C-identity}\Ntarget{eq:C-identity,ss:44.5}
The torus identities:
\begin{itemize}
\item By Lemma~\Nref{lem:Api-identities}(a), $\bG DR=DG'R$, hence $\bG DRG'D^T=DG'RG'D^T$.
\item By Lemma~\Nref{lem:slice}(c) with $E=0$, $\bQ\Pi=\bQ$, so $\bQ_\pi=\bQ$ and $\boldsymbol\Gamma_\pi=\boldsymbol\Gamma$.
\item The definition (Part~V) is $\bC=\bG-DG'RG'D^T-\bH\bQ\bG$.
\end{itemize}
So
\[
\bC=\bG(1-DRG'D^T)-\bH\bQ\bG=\bG\Pi^T-\bH\bQ\bG,\qquad\Pi^T=1-S^TD^T,\quad S^T=DRG' .
\]
Expanding $R=1-P_\cV$ gives the weakened operator
\[
\bC(s)=\bG(s)-\bG(s)DG'(s)D^T+\bG(s)DP_\cV(s)G'(s)D^T-\bH(s)\bQ\bG(s) .
\]
No identity is weakened alone. At $s=1$ every family represents its operator, so $\bC(1)=\bC$ and $\bH(1)=\bG\bQ^T\boldsymbol\Gamma=\bH$. The products are weakenable by Lemmas~\Nref{lem:weak-chain} and~\Nref{lem:weak-local-ops}(a), each costing $\delta$.
\status{Expanded.}
\end{detail}

\begin{detail}[Theorem~\Nref{thm:HC-pointwise}]\label{d9:HC}\Ntarget{thm:HC-pointwise}
All tier norms below are global.
\begin{itemize}
\item $P_\cV(s)$ maps $0\to3$. Indeed $Q'G'(s)$ maps $L^2$ to coarse fields, $(M')^{-1}(s)$ is a unit-lattice family, $Q'^T$ maps coarse fields to tier~$1$, and $G'(s)$ maps $1\to3$ (Theorem~\Nref{thm:tier-walk} at (L1)). So $R(s)=1-P_\cV(s)$ maps $2\to2$ inside products.
\item $\bH(s)$: $\boldsymbol\Gamma(s)$ is a unit-lattice family, $\bQ^T$ maps coarse fields to tier~$1$, and $\bG(s)$ maps $1\to3$.
\item $\bG(s)f$: tier $1\to3$, with $\mathfrak n^1_{y'}(f)\le|f|_{B(y')}$.
\item $\bG(s)DG'(s)D^Tf$ and $\bG(s)DP_\cV(s)G'(s)D^Tf$: $G'(s)D^T$ maps $1\to2$ (values), $P_\cV(s)$ maps $0\to3\subset2\to2$, and $\bG(s)D$ maps $2\to3$ (Theorem~\Nref{thm:tier-walk}(b)).
\item $\bH(s)\bQ\bG(s)f$: $\bQ\bG(s)f$ is a coarse field bounded by $L^2$ norms.
\item $\bC(s)D^Tg$: $\bG(s)D^T$ maps $2\to3$. In the other terms $G'(s)D^TD^Tg=G'(s)\mathfrak F^Tg$, because $D^TD^T=\mathfrak F^T$ is an identity between local operators. It holds inside every weakened product, since the families $G'(s)D^T$ have the terms $G'_\omega D^T$. And $\mathfrak F^Tg$ is bounded.
\end{itemize}
The H\"older quotient of $g$ is converted between local gauges as in the proof of Theorem~\Nref{thm:ID}(b), at a cost $C\eps|g|$.
\status{Expanded.}
\end{detail}

\begin{detail}[The closing list of Section~\Nref{sec:tiers}]\label{d9:closing}\Ntarget{sec:tiers}
\emph{``Not needed: $D^TD\bH(s)$''} is checked against the gradient formula \Neqref{eq:gradV-s}. There $D^TD\bH\mathcal D$ occurs only as $D^Tn_2$ with $n_2=-D\bH(s)\mathcal D_s$, which is handled by $\bC(s)D^T\colon2\to3$ and $|D\bH(s)\omega|_{(2)}\le C|\omega|_\infty$ (Lemma~\Nref{lem:weak-bg-linear}(b)). Also $\bH^TD^TDA'=(D\bH)^T(DA')$ and $\bH^TJ=(D\bH)^T\mathsf F$ are coarse fields bounded by $L^2$ norms (Detail~\ref{d9:gradV}). So the list is complete for $A_*(s)$.
\status{Expanded.}
\end{detail}

% IX.10
\subsection{The weakened background map}\label{d9:bg}

\begin{detail}[The setting of Section~\Nref{sec:weak-background} and Lemma~\Nref{lem:critical-current}]\label{d9:crit-current}\Ntarget{sec:weak-background,lem:critical-current}
\begin{step}[The Setting of Section~\Nref{sec:weak-background}]
\begin{itemize}
\item $U_0$ is $(M,\eps)$-regular, hence $(1,\eps)$-regular by Lemma~\Nref{lem:regular-monotone}. So Proposition~\Nref{prop:action} applies for $\eps\le c_5$. It gives $\mathcal N$ for complex $A$ with $|A|_\infty\le c_5$, with $|\nabla\mathcal N(A)|_\infty\le C(\eps+|A|_\infty)|A|_{C^1}$.
\item $|\mathsf F|=|F^s_{U_0}|\le2|F_{U_0}|\le C\eps$ (Lemma~\Nref{lem:current}(c)).
\item With $M\ge64$ and $r\le M/16$: $M/4-2r-2\ge M/8-2\ge6>2$.
\item The comparison of $|\cdot|_*$ with the graded-domain quantities, with the constant $C_{\rm conv}$, is Lemma~\Nref{lem:norm-conversion} (forward; Section~\ref{xa:sec}).
\end{itemize}
\end{step}
\begin{step}[The identity]
By Corollary~\Nref{cor:analytic}, $A_*(0)=0$ relative to $U'$. The Euler--Lagrange equation (Theorem~\Nref{thm:variational}(d)) at $A_*=0$ reads $\nabla\Phi(0)=d\mathcal Q(0)^T\omega_0$. Here $\nabla\Phi(0)=J_{U'}$ (Proposition~\Nref{prop:action}) and $d\mathcal Q(0)=\bQ$ (Proposition~\Nref{prop:Q-expansion}). So $J_{U'}=\bQ^T\omega_0$. Since $\bQ\bH=I$, we get $\omega_0=\bH^T\bQ^T\omega_0=\bH^TJ_{U'}$. With $J=D^TF^s$ (Lemma~\Nref{lem:current}(a)), $\bH^TJ=(D\bH)^TF^s$.
\end{step}
\begin{step}[The bound]
$|\omega_0(c)|\le\sum_y\|\one_c(D\bH)^T\one_y\|\,\|F^s\|_{B(y)}$, and $\|F^s\|_{B(y)}\le\binom d2^{1/2}|F^s|_\infty$, since blocks have volume $1$. The kernel sum is bounded by the decay of $D\bH$ (Theorem~\Nref{thm:reg-min}), a transpose-invariant block norm. Then $|J|_\infty\le|\bQ^T\omega_0|_\infty\le2|\omega_0|_\infty$, for two reasons:
\begin{itemize}
\item By Proposition~\Nref{prop:linearized-averages}(e), $\bQ^T$ differs from the flat transpose by $C_L\eps$ in sup norm.
\item The flat transpose has a non-negative kernel with $\sum_yw_\mu(y,x)=1$ (Lemma~\Nref{lem:avg}(a),(c)), so its sup norm is at most $1$.
\end{itemize}
With $|F^s|\le2|F|$ this gives $|J|_\infty\le C_J|F|_\infty$.
\end{step}
\begin{step}[Constants]
$C_J$ is the decay constant of $D\bH$ at regular backgrounds. By Remark~\Nref{rem:general-averaging} and Lemma~\Nref{lem:uniformity}(a), it depends on the averaging pair only through properties (i)--(iv). So it is free of $L$, for $\eps$ small depending on $L$ and $M$, and free of $M\ge M_0$. The expansion of Lemma~\Nref{lem:uniformity} (``no other property is used'') is Section~\ref{d7:unif} of the chapter on Part~VII, and its Step~0 ($M$-independence) is Detail~\ref{d7:step0}. This item cites them and does not re-derive them.
\end{step}
\begin{step}[The paragraph after the lemma]
With Proposition~\Nref{prop:sharp-curvature}, $|F_{U_k(V)}|\le(C_*+\delta)\eps_1$, so both the curvature and the current are at most $(1+C_J)(C_*+\delta)\eps_1$. In the units of step $k$, a field strength in the units of step $k+1$ is multiplied by $L^{-2}$, and a current, one more derivative, by $L^{-3}$. The use of this in the graded domain of step $k+1$ is Remark~\Nref{rem:domain-constants}(iii) (chapter on Part~VIII).
\end{step}
\status{Expanded (the constants rest on the expansion of Lemma~\Nref{lem:uniformity} in Section~\ref{d7:unif}).}
\end{detail}

\begin{detail}[The gradient \Neqref{eq:gradV-s}, Section~\Nref{ss:45.2}]\label{d9:gradV}\Ntarget{eq:gradV-s,ss:45.2}
$\bH=\bH(s)$ does not depend on $A'$. Write $\mathcal D=\mathcal D_s(A')$ and $d\mathcal D=d\mathcal D_s(A')$, and differentiate in a direction $a$, transposing with the bilinear form $\mathrm b$:
\begin{itemize}
\item $-\ip{\bH d\mathcal D\,a}{J}$ gives $-d\mathcal D^T\bH^TJ=-d\mathcal D^T(D\bH)^T\mathsf F$, using $J=D^T\mathsf F$;
\item $-\ip{Da}{D\bH\mathcal D}$ gives $-D^TD\bH\mathcal D=D^Tn_2$;
\item $-\ip{DA'}{D\bH d\mathcal D\,a}$ gives $-d\mathcal D^T(D\bH)^TDA'$;
\item $\frac12\cdot2\ip{D\bH d\mathcal D\,a}{D\bH\mathcal D}$ gives $d\mathcal D^T(D\bH)^TD\bH\mathcal D$;
\item $\mathcal N(A'-\bH\mathcal D)$ gives $\nabla\mathcal N(A)-d\mathcal D^T\bH^T\nabla\mathcal N(A)$.
\end{itemize}
The sum is $n_0+D^Tn_2$, as stated. Compared with Lemma~\Nref{lem:V-bounds} there are three changes: $D^TD\bH\mathcal D$ is kept as $D^T(D\bH\mathcal D)$, $\bH^TD^T$ is written as $(D\bH)^T$, and $\bH^TJ$ as $(D\bH)^T\mathsf F$.
\status{Expanded.}
\end{detail}

\begin{detail}[Lemma~\Nref{lem:weak-bg-linear}]\label{d9:bg-linear}\Ntarget{lem:weak-bg-linear}
\begin{itemize}
\item (a), and the first part of (b): Theorem~\Nref{thm:HC-pointwise}, with $\sum_{y'}e^{-\gamma|y-y'|}=S_\gamma$.
\item (b), second part: $D\bH(s)=D\bG(s)\,\bQ^T\boldsymbol\Gamma(s)$. Here $\boldsymbol\Gamma(s)$ is a unit-lattice family with $|\boldsymbol\Gamma(s)\omega|_\infty\le C|\boldsymbol\Gamma(s)|_\gamma|\omega|_\infty$, and $|\bQ^T\omega|_{B(y)}\le C|\omega|$ near $y$ gives tier~$1$. $D\bG(s)$ maps $1\to2$ (Theorem~\Nref{thm:tier-walk}(a) at (L3)) in a global tier norm of the family. By Lemma~\Nref{lem:weak-basic}(a) in tier norms (Detail~\ref{d9:tier-calc}(f)), this gives a pointwise bound for each $s\in\mathbb D_\kappa$, which is the bound in $|\cdot|_{(2)}$.
\item (c): the values of a coarse field on the bonds starting at $y$ are bounded by its block norm there. Then
\[
\|\one_y\bH(s)^Tf\|\le\sum_{y'}\|\one_y\bH(s)^T\one_{y'}\|\,\|f\|_{B(y')}\le C|\bH(s)|_\gamma|f|_\infty ,
\]
where $|\cdot|_\gamma$ is invariant under transposition (it is the maximum of row and column sums). Also $|\bH(s)|_\gamma+|D\bH(s)|_\gamma\le2^d\|\bH\|^D_{\kappa,\gamma}<\infty$, because $\bH(s)$ is a product of the (L3) output, the local $\bQ^T$, and the (L4) output.
\end{itemize}
\status{Expanded.}
\end{detail}

\begin{detail}[Lemma~\Nref{lem:Ds}]\label{d9:Ds}\Ntarget{lem:Ds}
Let $K_H$ be the constant with $|\bH(s)X|_*\le K_H|X|_\infty$. For $|X|_\infty,|X'|_\infty\le r_P$ and $|A'|_\infty\le r_P$, the argument $A=A'-\bH(s)X$ has $|A|_\infty\le(1+K_H)r_P\le c_3$, the domain of Proposition~\Nref{prop:Q-expansion}(b),(c). There $|\Cav_k(A)|\le C|A|^2$ and $|\Cav_k(A)-\Cav_k(\tilde A)|\le C\max(|A|,|\tilde A|)|A-\tilde A|$. So $X\mapsto \Cav_k(A'-\bH(s)X)$ maps the ball to itself and is a contraction with constant $2C(1+K_H)r_PK_H<1$ for $r_P$ small.

The iterates are analytic in $(s,A')$ and converge uniformly, so the fixed point is analytic.

For $|A'|\le r$: $|X|\le C(r+K_H|X|)^2\le2Cr^2+2CK_H^2|X|^2$. Since $2CK_H^2|X|\le2CK_H^2r_P\le\frac12$, this gives $|X|\le4Cr^2$.

Differentiating $X=\Cav_k(A'-\bH X)$ gives $(I+d\Cav_k(A)\bH)\,d\mathcal D=d\Cav_k(A)$. Transposing, $d\mathcal D^T=d\Cav_k(A)^T(I+\bH^Td\Cav_k(A)^T)^{-1}$. On coarse sup norms, $|d\Cav_k(A)^T\omega|_\infty\le C|A|_\infty|\omega|_\infty$, and $|\bH^T f|_\infty\le C_P|f|_\infty$ by Detail~\ref{d9:bg-linear}. So $\|\bH^Td\Cav_k(A)^T\|\le C_PC(1+K_H)r_P\le\frac12$, and the Neumann series gives $|d\mathcal D^T\omega|\le2C|A||\omega|\le C_Pr|\omega|$.
\status{Expanded.}
\end{detail}

\begin{detail}[Lemma~\Nref{lem:gradV-s}]\label{d9:gradVs}\Ntarget{lem:gradV-s}
\begin{step}[Size]
For $|A'|_{C^1}\le r\le2r'_P\le r_P$ we have $|\mathcal D_s|\le C_Pr^2$. Hence $|A|_{C^1}\le r+K_HC_Pr^2\le2r$, using $|\cdot|_{C^1}\le|\cdot|_*$ and $K_HC_Pr\le1$. The terms:
\begin{itemize}
\item $|n_2|_{(2)}\le C_P|\mathcal D_s|_\infty\le C_Pr^2$ (Detail~\ref{d9:bg-linear}(b)).
\item $|d\mathcal D^T(D\bH)^T(DA'-D\bH\mathcal D+\mathsf F)|\le C_Pr\cdot C_P(Cr+C_Pr^2+C\eps)$.
\item $|\nabla\mathcal N(A)|\le C(\eps+2r)2r$ for $2r\le c_5$ (Proposition~\Nref{prop:action}(a), complex $A$).
\item $|d\mathcal D^T\bH^T\nabla\mathcal N(A)|\le C_Pr\cdot C_P\cdot C(\eps+2r)2r$.
\end{itemize}
So $|n_0|_\infty\le C_P(\eps+r)r$.
\end{step}
\begin{step}[Lipschitz bounds, by Cauchy]
We use the form: if $f$ is analytic on the ball $\{|A'|_{C^1}<2r\}$ of a complex Banach space, with $\sup|f|\le K$, then $|f(A')-f(\tilde A')|\le\frac Kr|A'-\tilde A'|$ for $|A'|,|\tilde A'|\le r$. Indeed, $f(A')-f(\tilde A')=\int_0^1df(\xi_t)[h]\,dt$ with $\xi_t$ in the $r$-ball and $h=A'-\tilde A'$. Moreover $df(\xi)[h]=\frac1{2\pi\ii}\oint_{|z|=r/|h|}f(\xi+zh)z^{-2}\,dz$, and $|\xi+zh|\le2r$ on the circle. Apply this to $n_0$ with $K=C_P(\eps+2r)2r$, which gives $C_P(\eps+r)$, and to $n_2$ with $K=C_P(2r)^2$, which gives $C_Pr$. The radius $2r\le2r'_P$ is allowed by the first part.
\end{step}
\status{Expanded.}
\end{detail}

\begin{detail}[Theorem~\Nref{thm:weak-bg}]\label{d9:weak-bg}\Ntarget{thm:weak-bg}
\emph{The idea.} As in Part~V, the background map is the fixed point of a contraction, now for the weakened operators and in the weakened tier norm $|\cdot|_*$. Step~1 gives the contraction on $\{|A_1|_*\le\varrho_P\}$, with $4C_P(\eps_P+2\varrho_P)\le1$, and the a priori bounds; analyticity (Step~2) follows from the uniform convergence of the iterates; at $s=1$, for a critical reference $U_0$ and $|B|_\infty\le\min(R_P,\beta_V)$, the weakened map is the map of Part~V (Step~3); and the factorization (Step~4) follows from the block-diagonal form of the weakened families at $s_\Box=0$ off $Z$ (Lemma~\Nref{lem:weak-basic}(b)) and the locality of the maps of the iteration, while the dependence on $B$ only on $Z_i^{(r)}$ holds because $B$ enters only through $\bH(s)B$.
\begin{step}[(a)]
Let $b=|B|_\infty\le R_P=\varrho_P/K_H$ and $|A_1|_*\le\varrho_P$. Then $r:=|A'|_*\le2\varrho_P\le r'_P$. By \Neqref{eq:gradV-s}, Detail~\ref{d9:bg-linear} and Detail~\ref{d9:gradVs},
\[
|\mathcal T_s(A_1)|_*\le C_P(|n_0|_\infty+|n_2|_{(2)})\le C_P(\eps+r)r,\qquad\operatorname{Lip}\le C_P(\eps+2\varrho_P) .
\]
Choose $4C_P(\eps_P+2\varrho_P)\le1$.
\begin{itemize}
\item \emph{Self-map.} $|\mathcal T_s(A_1)|_*\le\frac14\cdot2\varrho_P\le\varrho_P$.
\item \emph{Contraction.} The constant is $\frac14$.
\item \emph{A priori bound.} At the fixed point, $|A_1|_*\le\frac14(|A_1|_*+K_Hb)$, so $|A_1|_*\le\frac13K_Hb$. Reinserting gives $|A_1|_*\le C_P(\eps+2K_Hb)2K_Hb\le C_P(\eps+b)b$, and then $|A_*|_*\le|A'|_*+K_HC_P|A'|^2\le C_Pb$.
\item \emph{$B=0$.} $\mathcal D_s(0)=0$ by uniqueness, $\nabla\mathcal N(0)=0$, and $d\mathcal D_s(0)=d\Cav_k(0)(\cdots)=0$, because $\Cav_k$ is the non-linear part. So $n_0(0)=n_2(0)=0$, and the fixed point is $0$.
\end{itemize}
The threshold $\eps_P$ also collects the input thresholds listed in the proof of Theorem~\Nref{thm:weak-bg}. The completeness of that list is checked in the parameter ledger.
\end{step}
\begin{step}[(b)]
The iterates $\mathcal T_s^m(0)$ are analytic in $(s,B)$, because the weakened operators are analytic in the interior of $\mathbb D_\kappa$ and bounded on it (Lemma~\Nref{lem:weak-basic}(a)), and $\mathcal D_s$, $\nabla\mathcal N$, $\Cav_k$ are analytic. They converge uniformly on the closed polydisc, where they are continuous, so the limit is analytic in the interior and continuous up to the boundary. At real $s$, $B$ and real $U_0$ every term is real.
\end{step}
\begin{step}[(c)]
At $s=1$: $\bC(1)=\bC$ (Detail~\ref{d9:C-identity}) and $\bH(1)=\bH$. Lemma~\Nref{lem:linearizing}'s $\mathcal D$ is the solution of the same equation. It is $O(r^2)$, so it lies in $\{|X|\le r_P\}$ and equals $\mathcal D_1$ by uniqueness. Hence $\nabla\cV_1=\nabla\cV$.

Critical $U_0$ gives $\bC J=0$. Indeed, $0=A_*(0)=A_{1,*}-\bH\mathcal D(A_{1,*})$ and $\bQ A_{1,*}=0$ give $\mathcal D(A_{1,*})=\bQ\bH\mathcal D(A_{1,*})=0$, so $A_{1,*}=0$. Then $0=\mathcal T(0)=-\bC(J+\nabla\cV(0))=-\bC J$, since $\nabla\cV(0)=n_0(0)+D^Tn_2(0)=0$. So $\mathcal T_1=\mathcal T$ of \Neqref{eq:fixed-point}. The range of $\bC$ lies in $\cX$, so the absence of the constraint is immaterial.

The estimate of the proof of Proposition~\Nref{prop:contraction}, without its first term $C_C\eps$ (which came from $\bC J$), reads
\[
|\mathcal T(A_1)|_{C^1}\le C_CC(\eps+\rho+C_Hb)(\rho+C_Hb),\qquad\rho=|A_1|_{C^1} .
\]
At the fixed point $\rho=|A_{1,*}|_{C^1}\le K(\eps+b)$, and $C_CC(\eps+\rho+C_Hb)\le\frac12$. So $|A_{1,*}|_{C^1}\le2C_CC(\eps+K(\eps+b)+C_Hb)C_Hb\le C(\eps+b)b\le K_Hb$, by the choice $C(\eps_P+R_P)\le K_H$.

Lemma~\Nref{lem:gradV-s} uses only $|A'|_{C^1}\le2K_Hb\le2\varrho_P$. So $|A_{1,*}|_*=|\mathcal T_1(A_{1,*})|_*\le C_P(\eps+2K_Hb)2K_Hb\le\varrho_P$. By uniqueness in the ball, $A_{1,*}=A_1(1;B)$, and $A_*(1;B)=A_*(B)$.

Both sides are analytic on the complex ball $|B|<\min(R_P,\beta_V)$ (Corollary~\Nref{cor:analytic} and (b)) and agree on its real points, so they agree on the ball.
\end{step}
\begin{step}[(d) and the region of $B$]
At $s_\Box=0$ off $Z$, Lemma~\Nref{lem:weak-basic}(b) applies to $\bC(s)$, $\bC(s)D^T$, $\bH(s)$, $D\bH(s)$ and their transposes: they vanish off $Z^{(r)}\times Z^{(r)}$, are block-diagonal over the $Z_i^{(r)}$, and their blocks depend on $U_0$ only on $Z_i^{(r)}$. The non-linear maps are local:
\begin{itemize}
\item $\Cav_k(A)(c)$ reads $A$ and $U_0$ on $S_c=B(c_-)\cup B(c_+)$;
\item $\nabla\mathcal N(A)(b)$ reads the plaquettes containing $b$;
\item $\mathsf F$ is local in $U_0$.
\end{itemize}
All vanish where their argument vanishes.

The statement claims dependence on $B$ only on $Z_i^{(r)}$, one unit less than the other data. The proof states the class with $Z_i^{(r)}$ and gives the reason (proof of Theorem~\Nref{thm:weak-bg}(d), ``The class preserved by the iteration''):
\begin{itemize}
\item In both iterations, $B$ enters only through $\bH(s)B$ in $A'_s=A_1+\bH(s)B$.
\item $\mathcal D_s$, $n_0$, $n_2$ and $\mathcal T_s$ see $B$ only through $A'$.
\item The block of $\bH(s)$ on $Z_i$ reads only the coarse bonds starting in $Z_i^{(r)}$.
\end{itemize}
So the class is stated with ``$B$ on the coarse bonds of $Z_i^{(r)}$'', keeping $U_0$ on $Z_i^{(r+1)}$ (because $\Cav_k(\cdot)(c)$ reads $S_c$) and $X$ on $Z^{(r+1)}$. Both iterations, started at $0$, preserve this class. The components are at distance $\ge M/4-2r-2>2$, so the local maps do not couple them.

Section~\Nref{sec:gen-locality} uses exactly this sharper form (the proof of Lemma~\Nref{lem:gl-bg}: ``$\bH(s)B$ depends on $B$ only on the coarse bonds of $Z^{(r)}$'').
\end{step}
\status{Expanded.}
\end{detail}

\begin{detail}[The remarks of Section~\Nref{sec:weak-background}]\label{d9:bg-remarks}\Ntarget{sec:weak-background}
\begin{itemize}
\item Criticality enters only Step~3 of Detail~\ref{d9:weak-bg}. So (a), (b), (d) hold at every real regular reference (Remark~\Nref{rem:weak-bg-noncritical}). Their extension to complex references is Proposition~\Nref{prop:complex-ext}(d) (Section~\ref{xa:sec}).
\item In the step, $B=\hat Z(Z)$, and $\hat Z$ has range $2L$ (Lemma~\Nref{lem:eliminate}). The factorization after composition with $\hat Z$ needs the components to stay separated after a $2L$-thickening: $M/4-2r-2-4L>2$, i.e.\ $M\ge CL$.
\end{itemize}
\status{Expanded.}
\end{detail}

% IX.11
\subsection{Summary}\label{d9:status}

\paragraph{Items expanded here.}
\begin{itemize}
\item The full gradient for the propagators with $W_2$ (Detail~\ref{d9:L4}, Step~4).
\item The dependence region of $W_j$, and the locality of holomorphic continuations (Detail~\ref{d9:inventory}).
\item The connectedness in Lemma~\Nref{lem:quadratic-bound} (Detail~\ref{d9:one-loop}).
\item The $\mathrm b$-symmetry of $\bG^Z$, proved algebraically (Detail~\ref{d9:cut-term}).
\item The hat in (W2): both readings are covered (Detail~\ref{d9:fam-setting}).
\item The identification of the walk at $s=1$ at the gauge levels, by words (Detail~\ref{d9:wi-a}, Step~4), and the clauses of (W2) used for the derivatives in the compatibility (Detail~\ref{d9:wi-b}, Step~5).
\item The lower bound at (L5) is for the operator, the bound of Theorem~\Nref{thm:weak-inverse} carries $C_G$, and additive terms carry no tier bound alone (Details~\ref{d9:L5}, \ref{d9:wi-b}, \ref{d9:fam-def}).
\item Theorem~\Nref{thm:weak-inverse}, Proposition~\Nref{prop:one-loop-local} with its Jacobian half, and Remark~\Nref{rem:one-loop-local-use}(a) (Sections~\ref{d9:walk}, \ref{d9:torus}; Detail~\ref{d9:ol-use}).
\item Lemma~\Nref{lem:torus-transfer} at the complex flat references, with comparison constant at most $2$ (Detail~\ref{d9:transfer}, Step~6).
\item The branch of $\log\det$ at complex references (Detail~\ref{d9:one-loop}).
\end{itemize}

\paragraph{Dependences on other chapters.}
\begin{itemize}
\item The reading of the doubled constants at complex references, $C_{\rm cx}$ level by level, and the non-circularity of $M_{\mathfrak N}$ and $M'_1$: Detail~\ref{xa:d-majorant}.
\item The comparison constant $C_{\rm cmp}$: Detail~\ref{xa:d-cmp}.
\item $W_2$ at complex references: Detail~\ref{d8:A1}.
\item The choice of $\eps_a$, $\delta_a$, $c_a$, also at $k=0$: the parameter ledger, Detail~\ref{det:ledger-eps-M}.
\item The global order of choices and the $M$-independence claims: the parameter ledger (Section~\ref{subsec:ledger-order}, Detail~\ref{det:ledger-M0}); the chapter on Part~VII (Section~\ref{d7:unif}, Detail~\ref{d7:step0}).
\end{itemize}

\paragraph{Not expanded.} These items carry no inference used later:
\begin{itemize}
\item Remark~\Nref{rem:L3-domains};
\item the example of Section~\Nref{ss:41.1}, paragraph ``Factorized families'' (why derivative norms of middle factors are not available);
\item Remark~\Nref{rem:one-loop-polymer}(a)--(c) (heuristics);
\item Route~1 of Section~\Nref{sec:pointwise-weak};
\item Proposition~\Nref{prop:schauder} and Remarks~\Nref{rem:schauder-averaging}, \Nref{rem:C-formula} (not used).
\end{itemize}

\setlist{itemsep=3pt,topsep=5pt}
